% =============================================================================
%  Documentacao tecnica do Rust VGA Kernel
%  Compile com: latexmk -pdf main.tex   (ou make)
% =============================================================================
\documentclass[11pt,a4paper,oneside]{report}

\usepackage[T1]{fontenc}
\usepackage[utf8]{inputenc}
\usepackage[brazilian]{babel}
\usepackage{lmodern}
\usepackage{microtype}
\usepackage[a4paper,top=2.6cm,bottom=2.6cm,left=2.6cm,right=2.6cm]{geometry}

\usepackage{xcolor}
\usepackage{graphicx}
\usepackage{booktabs}
\usepackage{longtable}
\usepackage{array}
\usepackage{tabularx}
\usepackage{enumitem}
\usepackage{amssymb}
\usepackage{titlesec}
\usepackage{fancyhdr}
\usepackage{listings}
\usepackage[most]{tcolorbox}
\usepackage{tikz}
\usetikzlibrary{arrows.meta,positioning,shapes.geometric,fit,backgrounds,calc}
\usepackage{caption}
\usepackage{float}
\usepackage[hidelinks,unicode]{hyperref}

% -----------------------------------------------------------------------------
% Cores e estilo
% -----------------------------------------------------------------------------
\definecolor{rinuxblue}{HTML}{1F4E79}
\definecolor{rinuxlight}{HTML}{EAF1F8}
\definecolor{rinuxgray}{HTML}{4D4D4D}
\definecolor{codegreen}{HTML}{2E7D32}
\definecolor{codegray}{HTML}{6A737D}
\definecolor{codepurple}{HTML}{8E24AA}
\definecolor{codebackground}{HTML}{F6F8FA}

\titleformat{\chapter}[display]
  {\normalfont\bfseries\color{rinuxblue}}
  {\filright\Large Cap\'itulo \thechapter}
  {0.6ex}
  {\Huge}
\titlespacing*{\chapter}{0pt}{-10pt}{24pt}

\titleformat{\section}
  {\normalfont\Large\bfseries\color{rinuxblue}}{\thesection}{0.8em}{}
\titleformat{\subsection}
  {\normalfont\large\bfseries\color{rinuxgray}}{\thesubsection}{0.8em}{}

\pagestyle{fancy}
\fancyhf{}
\fancyhead[L]{\small\textcolor{rinuxgray}{Rust VGA Kernel --- Documenta\c{c}\~ao T\'ecnica}}
\fancyhead[R]{\small\textcolor{rinuxgray}{\thepage}}
\renewcommand{\headrulewidth}{0.4pt}
\setlength{\headheight}{14pt}

\setlist[itemize]{topsep=2pt,itemsep=2pt,parsep=0pt}
\setlist[enumerate]{topsep=2pt,itemsep=2pt,parsep=0pt}

\captionsetup{font=small,labelfont=bf,width=0.95\linewidth}

% -----------------------------------------------------------------------------
% Listings (Rust)
% -----------------------------------------------------------------------------
\lstdefinelanguage{Rust}{
  keywords={as,break,const,continue,crate,else,enum,extern,false,fn,for,if,impl,in,let,loop,match,mod,move,mut,pub,ref,return,self,Self,static,struct,trait,true,type,unsafe,use,where,while,async,await,dyn,yield},
  keywordstyle=\color{rinuxblue}\bfseries,
  ndkeywords={u8,u16,u32,u64,usize,i32,bool,char,str,Option,Some,None,Result,Ok,Err,AtomicBool,UnsafeCell},
  ndkeywordstyle=\color{codepurple},
  sensitive=true,
  morecomment=[l]{//},
  morecomment=[s]{/*}{*/},
  morestring=[b]",
  morestring=[b]',
}

\lstset{
  language=Rust,
  backgroundcolor=\color{codebackground},
  basicstyle=\ttfamily\footnotesize,
  keywordstyle=\color{rinuxblue}\bfseries,
  commentstyle=\color{codegreen}\itshape,
  stringstyle=\color{codepurple},
  numbers=left,
  numberstyle=\tiny\color{codegray},
  numbersep=8pt,
  frame=single,
  rulecolor=\color{black!20},
  breaklines=true,
  breakatwhitespace=false,
  postbreak=\mbox{\textcolor{codegray}{$\hookrightarrow$}\space},
  showstringspaces=false,
  tabsize=4,
  captionpos=b,
  columns=fullflexible,
  keepspaces=true,
  upquote=true,
}

\lstdefinestyle{shell}{
  language={},
  basicstyle=\ttfamily\footnotesize,
  numbers=none,
  frame=single,
  backgroundcolor=\color{black!3},
}

% Caixas de destaque
\newtcolorbox{nota}[1][]{
  enhanced, colback=rinuxlight, colframe=rinuxblue,
  boxrule=0.5pt, left=6pt, right=6pt, top=4pt, bottom=4pt,
  fonttitle=\bfseries, title=#1
}
\newtcolorbox{atencao}[1][]{
  enhanced, colback=orange!7, colframe=orange!70!black,
  boxrule=0.5pt, left=6pt, right=6pt, top=4pt, bottom=4pt,
  fonttitle=\bfseries, title=#1
}

\newcommand{\code}[1]{\texttt{#1}}
\newcommand{\termo}[1]{\textit{#1}}

% Metadados
\title{\vspace{-2cm}\textbf{Rust VGA Kernel}\\[0.4em]
       \Large Documenta\c{c}\~ao T\'ecnica, Arquitetura e Compara\c{c}\~ao com o Linux}
\author{Documenta\c{c}\~ao gerada a partir do c\'odigo-fonte}
\date{Vers\~ao 1.0 --- 3 de novembro de 2026}

\begin{document}

\hypersetup{
  pdftitle={Rust VGA Kernel - Documentacao Tecnica},
  pdfauthor={Rust VGA Kernel},
  pdfsubject={Kernel minimalista x86-64 em Rust},
  pdfkeywords={Rust, kernel, x86-64, multiboot2, VGA, PS/2, Linux}
}

\maketitle

\begin{abstract}
Este documento descreve integralmente o \textbf{Rust VGA Kernel}, um kernel
did\'atico escrito em Rust com \code{no\_std} para PCs padr\~ao com BIOS/x86\_64.
O texto cobre a arquitetura, a cadeia de boot, o modelo de mem\'oria, os
drivers de v\'ideo VGA, teclado PS/2 e serial, o shell embutido, o processo de
compila\c{c}\~ao e empacotamento da ISO, os testes realizados e uma an\'alise
quantitativa e qualitativa do grau de similaridade com o kernel Linux.
O c\'odigo implementa o caminho m\'inimo entre o bootloader e um terminal
funcional, acrescido de IDT, timer, paginas de 4 KiB, \emph{demand paging} e
uma demonstra\c{c}\~ao de c\'odigo em ring 3. N\~ao \'e um porte do Linux e
n\~ao implementa escalonador, heap, sistema de arquivos ou um conjunto
completo de chamadas de sistema.
\end{abstract}

\tableofcontents
\listoffigures
\listoftables

\clearpage
\chapter{Introdu\c{c}\~ao}
\label{cap:introducao}

\section{Motiva\c{c}\~ao}

O kernel Linux \'e um sistema operacional completo, com milh\~oes de linhas de
c\'odigo e d\'ecadas de evolu\c{c}\~ao. Estudar seu funcionamento a partir do
c\'odigo real exige percorrer camadas de abstra\c{c}\~ao, configura\c{c}\~oes de
compila\c{c}\~ao, subsistemas e caminhos de c\'odigo que n\~ao s\~ao essenciais
para entender o n\'ucleo do problema ``como um computador sai do boot e abre um
terminal?''.

O \textbf{Rust VGA Kernel} foi criado como um recorte did\'atico desse
problema. Ele reimplementa, em Rust, o caminho m\'inimo necess\'ario para:

\begin{enumerate}
  \item ser carregado por um bootloader padr\~ao (GRUB via Multiboot2);
  \item sair do modo de 32 bits e entrar em modo longo x86\_64;
  \item configurar p\'aginas e uma pilha de kernel;
  \item inicializar um console de texto VGA 80$\times$25;
  \item ler teclas do controlador PS/2 8042 e da serial COM1;
  \item instalar uma IDT, receber exce\c{c}\~oes e o timer do PIT;
  \item alocar frames de 4 KiB e materializar p\'aginas sob demanda;
  \item executar c\'odigo em ring 3 com separa\c{c}\~ao usu\'ario/kernel;
  \item oferecer um pequeno shell com comandos b\'asicos.
\end{enumerate}

O resultado n\~ao \'e um Linux reduzido, nem um porte linha a linha. \'E um
kernel independente, sem depend\^encias externas de Rust, que reproduz
\emph{conceitos} dos subsistemas equivalentes do Linux em uma escala
compreens\'ivel.

\section{Objetivo do documento}

Esta documenta\c{c}\~ao tem quatro objetivos:

\begin{itemize}
  \item \textbf{Explicar a arquitetura} do kernel: camadas, fluxo de boot,
        modelo de mem\'oria, modelo de execu\c{c}\~ao e depend\^encias entre
        m\'odulos.
  \item \textbf{Documentar as principais fun\c{c}\~oes} de cada m\'odulo,
        incluindo contratos, pr\'e-condi\c{c}\~oes, efeitos colaterais e
        limita\c{c}\~oes. O Ap\^endice~\ref{ap:funcoes} traz o cat\'alogo
        completo das 96 fun\c{c}\~oes do projeto.
  \item \textbf{Registrar o processo de build, execu\c{c}\~ao e valida\c{c}\~ao},
        de forma que outra pessoa consiga reproduzir o ambiente e os testes.
  \item \textbf{Comparar o projeto com o kernel Linux}, apresentando
        correspond\^encias de arquivos e fun\c{c}\~oes e uma estimativa
        transparente do n\'ivel de similaridade.
\end{itemize}

\section{Escopo}

\subsection*{O que est\'a inclu\'ido}

\begin{itemize}
  \item Boot BIOS por CD/ISO usando GRUB e o protocolo Multiboot2.
  \item Transi\c{c}\~ao de 32 bits para modo longo (x86\_64).
  \item GDT com segmentos ring 0/ring 3, TSS e pilhas de privil\'egio;
        tabelas de p\'agina est\'aticas no boot e metade alta do kernel.
  \item IDT com 256 portas, exce\c{c}\~oes, PIC remapeado e timer PIT em IRQ0.
  \item Alocador de frames f\'isicos de 4 KiB e \emph{demand paging} para
        regi\~oes de usu\'ario e de kernel.
  \item Demos em ring 3 com prote\c{c}\~ao U/S e \emph{syscall} did\'atica
        \code{int \$0x80}.
  \item Console VGA em modo texto, com cursor de hardware, rolagem e
        atributos de cor.
  \item Driver de teclado PS/2 por \emph{polling}, com Shift, Enter, Tab e
        Backspace.
  \item Driver da UART 16550 em COM1, usado para depura\c{c}\~ao e entrada.
  \item Shell sem aloca\c{c}\~ao din\^amica com dez comandos.
  \item \emph{Panic handler} que imprime a falha e para a CPU.
  \item Build reprodut\'ivel com Cargo, \code{ld.lld}, \code{grub-mkrescue}
        e QEMU.
\end{itemize}

\subsection*{O que n\~ao est\'a inclu\'ido}

\begin{itemize}
  \item Heap, alocador de objetos, \code{alloc} ou cole\c{c}\~oes din\^amicas.
  \item Processos completos, threads, escalonador, troca de contexto,
        espa\c{c}os de endere\c{c}o por processo ou um conjunto amplo de
        chamadas de sistema.
  \item Sistema de arquivos, rede, USB, ACPI, SMP ou suporte UEFI nativo.
  \item Drivers para outros v\'ideos al\'em do VGA texto legado.
  \item Endurecimento de seguran\c{c}a como SMEP, SMAP, NX, KASLR ou
        valida\c{c}\~ao de ponteiros de usu\'ario.
\end{itemize}

\section{P\'ublico-alvo}

O documento \'e destinado a:

\begin{itemize}
  \item estudantes de sistemas operacionais e arquitetura de computadores;
  \item desenvolvedores Rust interessados em \code{no\_std} e bare metal;
  \item pessoas que desejam entender o caminho entre BIOS, bootloader e
        console antes de estudar o kernel Linux em profundidade;
  \item avaliadores que precisem julgar o grau de correspond\^encia entre
        este projeto e o Linux.
\end{itemize}

\section{Conven\c{c}\~oes}

\begin{itemize}
  \item Nomes de c\'odigo aparecem em \code{fonte monoespa\c{c}ada}.
  \item Endere\c{c}os hexadecimais usam o prefixo \code{0x}.
  \item ``Metade alta'' refere-se a endere\c{c}os acima de
        \code{0xFFFF800000000000}, conven\c{c}\~ao usada pelo Linux e pelo
        alvo \code{x86\_64-unknown-none}.
  \item ``Original'' ou ``Linux'' refere-se, salvo indica\c{c}\~ao contr\'aria,
        \`a \'arvore Linux 7.3.1 usada como refer\^encia de compara\c{c}\~ao.
  \item Valores de endere\c{c}o e tamanho de se\c{c}\~ao citados como
        ``observados'' foram extra\'idos da compila\c{c}\~ao de refer\^encia e
        podem variar ligeiramente entre vers\~oes do compilador.
\end{itemize}

\section{Estado do projeto}

A vers\~ao documentada \'e a \code{0.1.0}, com as seguintes caracter\'isticas:

\begin{table}[H]
\centering
\caption{Resumo do estado da vers\~ao documentada.}
\label{tab:estado}
\begin{tabularx}{\linewidth}{@{}lX@{}}
\toprule
\textbf{Item} & \textbf{Situa\c{c}\~ao} \\
\midrule
Compila\c{c}\~ao & \texttt{make} conclui sem erros; alvo \code{x86\_64-unknown-none}. \\
Imagem & \code{vga-kernel.iso} El Torito, inicializ\'avel por BIOS. \\
Boot & Validado em QEMU com SeaBIOS e GRUB Multiboot2. \\
Console & VGA texto 80$\times$25, cursor, rolagem e cores. \\
Entrada & PS/2 por polling e serial COM1. \\
Interrup\c{c}\~oes & IDT com 256 portas, exce\c{c}\~oes, PIC remapeado e PIT a
~100 Hz. \\
Mem\'oria & Alocador de frames de 4 KiB, tabelas de 4 KiB e demand paging. \\
Prote\c{c}\~ao & GDT/TSS, ring 3 e bits U/S para usu\'ario e kernel. \\
Shell & \code{help}, \code{clear}, \code{echo}, \code{info}, \code{demand},
\code{usertest}, \code{userprot}, \code{mem}, \code{reboot}, \code{halt}. \\
Testes & Boot, comandos, entrada PS/2/serial, timer, page faults e ring 3
verificados. \\
Cobertura Linux & Subconjunto pequeno; similaridade baixa (Cap\'itulo~\ref{cap:similaridade}). \\
\bottomrule
\end{tabularx}
\end{table}

\chapter{Vis\~ao geral do sistema}
\label{cap:visao}

\section{Resumo executivo}

O Rust VGA Kernel \'e um execut\'avel ELF x86\_64 carregado pelo GRUB. O GRUB
usa o protocolo Multiboot2 para transferir o controle a um pequeno \emph{stub}
em linguagem de montagem, ligado em 1 MiB. Esse stub constr\'oi tabelas de
p\'agina, ativa PAE e modo longo e salta para o c\'odigo Rust, ligado na metade
alta em \code{0xFFFFFFFF80000000}.

O c\'odigo Rust inicializa, nesta ordem, o renderizador VGA, a serial COM1, a
camada TTY com oito consoles virtuais, o alocador de frames, o heap do kernel,
o ramfs e o VFS, a GDT com ring 3/TSS, o teclado PS/2 e a IDT com PIC/PIT. Em
seguida liga as interrup\c{c}\~oes, cria a tabela de tarefas com o shell como
PID 1, imprime um \emph{banner}, valida a assinatura Multiboot2 e entra no
shell, que exibe o prompt \code{rinux>} no console ativo e processa comandos
lidos do teclado ou da serial.

O sistema inteiro cabe em cerca de 250 KiB de ELF e 6.797 linhas de Rust
(5.365 linhas efetivas de c\'odigo, sem coment\'arios e linhas em branco).
Inclui um heap \emph{first-fit} atr\'as de \code{\#[global\_allocator]}, um VFS
com dezesseis descritores sobre o ramfs e um \code{/proc} sint\'etico, um
escalonador \emph{round-robin} preemptivo por timer com \code{sleep}
bloqueante (\code{schedule\_timeout}) e filas de espera por evento
(\code{wait\_event}/\code{wake\_up}), e onze chamadas de sistema. Ainda
n\~ao h\'a \code{fork}/\code{exec}, espa\c{c}o de endere\c{c}os por processo
nem sistema de arquivos em disco.

\section{Funcionalidades}

\begin{table}[H]
\centering
\caption{Funcionalidades implementadas e seu estado.}
\label{tab:funcionalidades}
\begin{tabularx}{\linewidth}{@{}lXc@{}}
\toprule
\textbf{\'Area} & \textbf{Funcionalidade} & \textbf{Estado} \\
\midrule
Boot & Cabe\c{c}alho Multiboot2 v\'alido & Completo \\
Boot & Transi\c{c}\~ao 32 $\rightarrow$ 64 bits & Completo \\
Mem\'oria & PML4/PDPT/PD est\'aticos, 1 GiB & Completo \\
Mem\'oria & Metade alta do kernel & Completo \\
V\'ideo & Console VGA texto 80$\times$25 & Completo \\
V\'ideo & Rolagem, cursor, tab, backspace & Completo \\
V\'ideo & Oito consoles virtuais (\code{tty0}..\code{tty7}) & Completo \\
V\'ideo & Troca por Alt+F1..Alt+F8 e comando \code{tty} & Completo \\
V\'ideo & Framebuffer gr\'afico / GOP & Ausente \\
Teclado & PS/2 8042, conjunto 1, polling & Completo (subconjunto) \\
Teclado & Shift, Alt, Enter, Tab, Backspace & Completo \\
Teclado & Sele\c{c}\~ao de console com Alt+F1..Alt+F8 & Completo \\
Teclado & IRQ1, Caps Lock, AltGr, setas & Ausente \\
Serial & COM1 16550, 38400 8N1 & Completo (subconjunto) \\
Serial & Entrada pelo terminal do QEMU & Completo \\
Shell & \code{help}, \code{clear}, \code{echo} & Completo \\
Shell & \code{ls}, \code{cat}, \code{cp}, \code{mv}, \code{tree} & Completo \\
Shell & \code{info}, \code{tty}, \code{getty} & Completo \\
Shell & \code{reboot}, \code{halt} & Completo \\
Sistema & Panic handler & Completo \\
Sistema & IDT, exce\c{c}\~oes e IRQ0 (timer) & Completo (subconjunto) \\
Sistema & GDT/TSS, ring 3 e prote\c{c}\~ao U/S & Completo (demo) \\
Sistema & Syscalls \code{exit/write/read/get\_tty} & Completo (subconjunto) \\
Sistema & Syscalls de arquivo \code{open/close/read/write/lseek} & Completo (subconjunto) \\
Sistema & \code{getpid}/\code{getppid} vindos da tabela de tarefas & Completo \\
Usu\'ario & \code{getty}/login em ring 3 & Completo (demo) \\
Mem\'oria & Frames, p\'aginas 4 KiB e demand paging & Completo (subconjunto) \\
Mem\'oria & Heap do kernel, crate \code{alloc} e \code{Vec}/\code{String} & Completo (subconjunto) \\
Armazenamento & Ramfs em mem\'oria (48 n\'os, arquivos de 1 KiB) & Completo (subconjunto) \\
Armazenamento & VFS, 16 descritores e \code{lseek} & Completo (subconjunto) \\
Armazenamento & \code{/proc} sint\'etico (version, uptime, meminfo, ps) & Completo (subconjunto) \\
Processos & Tabela de tarefas, PID/PPID e estados & Completo (subconjunto) \\
Processos & Escalonador \emph{round-robin} preemptivo por timer & Completo (subconjunto) \\
Processos & \code{sleep} bloqueante (\code{schedule\_timeout}) & Completo (subconjunto) \\
Processos & Filas de espera por evento (\code{wait\_event}/\code{wake\_up}) & Completo (subconjunto) \\
Processos & \code{fork}/\code{exec}/\code{wait}, CFS/EEVDF & Ausente \\
Armazenamento & Driver de disco, mounts e permiss\~oes & Ausente \\
\bottomrule
\end{tabularx}
\end{table}

\section{Requisitos}

\subsection{Hardware}

\begin{itemize}
  \item PC compat\'ivel x86\_64 com BIOS legada ou CSM habilitado.
  \item Controlador de teclado PS/2 8042 (ou emula\c{c}\~ao equivalente).
  \item Adaptador VGA em modo texto, com buffer em \code{0xB8000}.
  \item Unidade de CD/DVD ou dispositivo equivalente para inicializar a ISO.
  \item Mem\'oria suficiente para carregar o kernel e a ISO; 128 MiB s\~ao
        suficientes no QEMU.
\end{itemize}

\subsection{Software de desenvolvimento}

\begin{itemize}
  \item Rust est\'avel com o alvo \code{x86\_64-unknown-none} instalado.
  \item \code{rustup} (recomendado; o Makefile prioriza o toolchain dele).
  \item \code{ld.lld} (LLVM linker).
  \item \code{grub-mkrescue}, \code{xorriso}, \code{libisoburn} e
        \code{mtools}.
  \item QEMU (\code{qemu-system-x86\_64}) para execu\c{c}\~ao e testes.
  \item Para esta documenta\c{c}\~ao: uma distribui\c{c}\~ao TeX Live com
        \code{pdflatex} e \code{latexmk}.
\end{itemize}

\section{Estrutura de diret\'orios}

\begin{lstlisting}[style=shell]
.
|-- Cargo.toml              metadados e perfis de compilacao
|-- Cargo.lock              lockfile (sem dependencias externas)
|-- rust-toolchain.toml     canal stable + alvo x86_64-unknown-none
|-- .cargo/config.toml      alvo padrao, linker e flags
|-- linker.ld               layout do ELF (stub baixo + kernel alto)
|-- Makefile                build do kernel, ISO e execucao no QEMU
|-- boot/grub/grub.cfg      menu Multiboot2 e console VGA
|-- src/
|   |-- main.rs             entrada Rust, macros e panic handler
|   |-- boot.rs             Multiboot2, GDT, paginacao e modo longo
|   |-- arch.rs             portas de I/O, interrupcoes, halt e reboot
|   |-- gdt.rs              GDT definitiva, TSS e pilhas de privilegio
|   |-- idt.rs              IDT, stubs, PIC, PIT e porta int 0x80
|   |-- paging.rs           frames, tabelas 4 KiB e demand paging
|   |-- heap.rs             heap do kernel e #[global_allocator]
|   |-- ramfs.rs            sistema de arquivos em memoria
|   |-- proc.rs             /proc virtual (uptime, meminfo, ps, ...)
|   |-- vfs.rs              VFS minimo e tabela de descritores
|   |-- task.rs             tarefas, escalonador, sleep e wait queues
|   |-- user.rs             ring 3, demos e getty/login
|   |-- syscall.rs          exit/write/read/open/close/lseek/getpid
|   |-- vga.rs              renderizador VGA 80x25
|   |-- tty.rs              consoles virtuais tty0..tty7
|   |-- keyboard.rs         teclado PS/2 (8042)
|   |-- serial.rs           UART 16550 / COM1
|   |-- shell.rs            prompt e comandos
|   `-- sync.rs             spinlocks (console, ramfs, heap, VFS, tarefas)
|-- docs/                   esta documentacao (LaTeX e PDF)
|-- build/                  arvore temporaria da ISO (gerada)
|-- target/                 artefatos do Cargo (gerado)
|-- vga-kernel.iso          imagem inicializavel (gerada)
`-- README.md               guia rapido de uso
\end{lstlisting}

\section{M\'etricas do c\'odigo}

\begin{table}[H]
\centering
\caption{M\'etricas dos m\'odulos Rust (contagem por linhas f\'isicas).}
\label{tab:metricas}
\begin{tabular}{@{}lrrl@{}}
\toprule
\textbf{M\'odulo} & \textbf{Linhas} & \textbf{Fun\c{c}\~oes} & \textbf{Responsabilidade} \\
\midrule
\code{boot.rs}     &  202 &  0 & Stub de boot e tabelas de p\'agina \\
\code{arch.rs}     &  101 &  7 & Portas de I/O, \code{irqsave}, halt, reboot \\
\code{gdt.rs}      &  159 &  6 & GDT, TSS e pilhas de privil\'egio \\
\code{idt.rs}      &  507 & 14 & IDT, PIC, PIT, exce\c{c}\~oes e IRQs \\
\code{paging.rs}   &  559 & 23 & Frames, tabelas e demand paging \\
\code{heap.rs}     &  270 &  9 & Heap do kernel e alocador global \\
\code{ramfs.rs}    &  534 & 31 & Sistema de arquivos em mem\'oria \\
\code{proc.rs}     &  191 &  9 & \code{/proc} sint\'etico \\
\code{vfs.rs}      &  302 & 11 & VFS e tabela de descritores \\
\code{task.rs}     &  746 & 29 & Tarefas, escalonador, sleep e wait queues \\
\code{user.rs}     &  446 &  9 & Ring 3, demos e getty \\
\code{syscall.rs}  &  244 &  9 & Syscalls de arquivo e processo \\
\code{vga.rs}      &   74 &  5 & Renderizador VGA texto \\
\code{tty.rs}      &  448 & 31 & Consoles virtuais e filas \\
\code{keyboard.rs} &  271 & 10 & Driver PS/2 8042 \\
\code{serial.rs}   &   70 &  5 & UART 16550 / COM1 \\
\code{shell.rs}    & 1468 & 61 & Linha de comando e applets \\
\code{main.rs}     &  139 &  2 & Entrada Rust, macros, panic \\
\code{sync.rs}     &   66 &  5 & Spinlocks seguros para IRQ \\
\midrule
\textbf{Total}     & \textbf{6.797} & \textbf{276} & 5.365 linhas efetivas de c\'odigo \\
\bottomrule
\end{tabular}
\end{table}

\begin{nota}[Sobre a contagem]
A contagem inclui 654 linhas de coment\'ario e 778 linhas em branco. O n\'umero
de fun\c{c}\~oes inclui m\'etodos de \code{impl}, fun\c{c}\~oes \code{const}
e fun\c{c}\~oes de trait (\code{Deref}, \code{DerefMut}, \code{Drop} e
\code{fmt::Write}); n\~ao inclui declara\c{c}\~oes \code{extern} sem corpo.
\code{boot.rs} n\~ao possui fun\c{c}\~oes Rust porque sua l\'ogica est\'a em
\code{global\_asm!}; o \emph{stub} \'e contado como bloco de montagem.
\end{nota}

\section{Fluxo resumido de execu\c{c}\~ao}

\begin{enumerate}
  \item A BIOS seleciona o CD e executa o El Torito.
  \item O GRUB carrega o ELF e l\^e o cabe\c{c}alho Multiboot2.
  \item O GRUB entra em modo protegido de 32 bits e salta para \code{\_start}.
  \item O stub salva \code{eax} e \code{ebx}, monta as p\'aginas, carrega a
        GDT e ativa o modo longo.
  \item O stub salta para \code{kernel\_main(magic, info)} na metade alta.
  \item O Rust limpa a tela VGA, configura a serial, inicializa a camada TTY
        com oito consoles virtuais, o alocador de frames, o heap do kernel, o
        ramfs e o VFS, a GDT/TSS e o teclado, e instala a IDT com PIC/PIT.
  \item O Rust liga as interrup\c{c}\~oes, cria a tabela de tarefas com o
        shell como PID 1 e o shell exibe \code{rinux>} no console ativo,
        processando comandos at\'e \code{halt}, \code{reboot} ou reset
        f\'isico.
\end{enumerate}

\chapter{Arquitetura}
\label{cap:arquitetura}

\section{Vis\~ao em camadas}

O kernel \'e organizado em seis camadas conceituais. As camadas inferiores
dependem apenas de hardware e de \code{core}; as superiores usam as
inferiores. N\~ao h\'a \code{std}, alocador de heap, \emph{runtime} de Rust
nem bibliotecas de terceiros.

\begin{figure}[H]
\centering
\begin{tikzpicture}[
  layer/.style={draw=rinuxblue, fill=rinuxlight, rounded corners=2pt,
                minimum width=11.4cm, minimum height=1.0cm,
                align=center, font=\small},
  hw/.style={draw=black!60, fill=black!8, rounded corners=2pt,
             minimum width=11.4cm, minimum height=1.0cm,
             align=center, font=\small},
  node distance=0.35cm
]
\node[layer] (shell) {\textbf{Interface} --- shell: linha de comando, \code{help}, \code{echo}, \code{info}, \code{demand}, \code{usertest}};
\node[layer, below=of shell] (services) {\textbf{Servi\c{c}os} --- console VGA, entrada PS/2, serial, macros de impress\~ao};
\node[layer, below=of services] (mem) {\textbf{Mem\'oria e interrup\c{c}\~oes} --- GDT/TSS, IDT/PIC/PIT, frames e \emph{demand paging}};
\node[layer, below=of mem] (core) {\textbf{N\'ucleo} --- \code{kernel\_main}, panic handler, spinlock, portas de I/O};
\node[layer, below=of core] (boot) {\textbf{Boot} --- stub 32/64 bits, GDT inicial, PML4/PDPT/PD, Multiboot2};
\node[hw, below=of boot] (hw) {\textbf{Hardware} --- CPU x86\_64, VGA, 8042, UART 16550, PIC/PIT, BIOS/GRUB};
\draw[-{Latex[length=2mm]}, thick, rinuxblue] (shell) -- (services);
\draw[-{Latex[length=2mm]}, thick, rinuxblue] (services) -- (mem);
\draw[-{Latex[length=2mm]}, thick, rinuxblue] (mem) -- (core);
\draw[-{Latex[length=2mm]}, thick, rinuxblue] (core) -- (boot);
\draw[-{Latex[length=2mm]}, thick, black!60] (boot) -- (hw);
\end{tikzpicture}
\caption{Camadas do Rust VGA Kernel.}
\label{fig:camadas}
\end{figure}

\section{Modelo de execu\c{c}\~ao}

O kernel executa em \textbf{ring 0}, em uma \'unica CPU, com
\textbf{interrup\c{c}\~oes desligadas}. Todo o processamento \'e feito por
\emph{polling}: o shell consulta o teclado e a serial em um la\c{c}o ativo,
usando \code{core::hint::spin\_loop()} quando n\~ao h\'a entrada.

\begin{table}[H]
\centering
\caption{Caracter\'isticas do modelo de execu\c{c}\~ao.}
\label{tab:modelo}
\begin{tabularx}{\linewidth}{@{}lX@{}}
\toprule
\textbf{Aspecto} & \textbf{Implementa\c{c}\~ao} \\
\midrule
Privil\'egio & Ring 0 para o kernel; ring 3 apenas nos demos de
\code{user.rs}, com GDT/TSS. \\
CPUs & Uma CPU; SMP n\~ao \'e inicializado. \\
Interrup\c{c}\~oes & IDT com 256 portas; PIC remapeado e timer IRQ0 a
~100 Hz. IRQ1 (teclado) permanece mascarada. \\
Entrada & Polling de \code{0x60}/\code{0x64} e de COM1; timer por IRQ. \\
Concorr\^encia & Spinlocks seguros para IRQ no console, ramfs, heap, VFS e
tabela de tarefas. O escalonador \emph{round-robin} preempta no quantum do
timer (5 ticks), troca em yield expl\'icito, bloqueia tarefas em \code{sleep}
at\'e um tick de prazo e as suspende em filas de espera at\'e um
\code{wake\_up}. \\
Mem\'oria & Frames f\'isicos de 4 KiB, tabelas sob demanda e page faults, um
heap \emph{first-fit} de 1 MiB e \code{alloc}/\code{Vec}/\code{String}. \\
Pilhas & 8 KiB no stub (baixa), 32 KiB no kernel (alta), 16 KiB para
interrup\c{c}\~oes de usu\'ario e 8 KiB de IST para double fault. \\
Tempo de vida & O shell nunca retorna; \code{halt} e \code{reboot} encerram o fluxo. \\
\bottomrule
\end{tabularx}
\end{table}

\section{Depend\^encias entre m\'odulos}

\begin{table}[H]
\centering
\caption{Depend\^encias principais entre m\'odulos.}
\label{tab:dependencias}
\begin{tabularx}{\linewidth}{@{}lX@{}}
\toprule
\textbf{M\'odulo} & \textbf{Usa / \'e usado por} \\
\midrule
\code{main.rs} & Declara todos os m\'odulos; chama \code{vga}, \code{serial},
\code{tty}, \code{paging}, \code{heap}, \code{ramfs}, \code{vfs}, \code{gdt},
\code{keyboard}, \code{idt}, \code{task}, \code{arch} e \code{shell};
implementa as macros de impress\~ao. \\
\code{boot.rs} & Assembly que chama \code{kernel\_main}; n\~ao \'e chamado por
Rust. \\
\code{arch.rs} & Portas de I/O, \code{irq\_save}/\code{irq\_restore},
\code{sti}, \code{halt} e \code{reboot}; usado por v\'arios m\'odulos. \\
\code{gdt.rs} & GDT/TSS e pilhas; usa apenas \code{core::arch}. \\
\code{idt.rs} & Usa \code{arch}, \code{paging}, \code{syscall} e
\code{user} para tratar exce\c{c}\~oes, IRQs e a \emph{syscall}. \\
\code{paging.rs} & Frames, tabelas e \emph{page faults}; usa apenas
\code{core::arch}. \\
\code{heap.rs} & Heap \emph{first-fit} e \code{\#[global\_allocator]}; usa
\code{paging} e \code{sync}. \\
\code{ramfs.rs} & \'Arvore em mem\'oria de 48 n\'os; usa \code{sync}. \\
\code{proc.rs} & \code{/proc} sint\'etico; renderiza a partir de \code{heap},
\code{idt}, \code{paging} e \code{task}. \\
\code{vfs.rs} & Tabela de descritores sobre ramfs, \code{/proc} e consoles;
usa \code{proc}, \code{ramfs}, \code{tty} e \code{sync}. \\
\code{task.rs} & Tabela de tarefas, pilhas, troca preemptiva por timer,
\code{sleep} bloqueante e filas de espera por evento; usa \code{sync} e
\code{arch}. \\
\code{user.rs} & Monta os programas, chama \code{paging} e \code{tty} e
retorna ao kernel por \code{user\_resume}. \\
\code{syscall.rs} & Implementa as onze chamadas; usa \code{paging},
\code{task}, \code{tty}, \code{user} e \code{vfs}. \\
\code{shell.rs} & Usa \code{tty}, \code{idt}, \code{paging}, \code{heap},
\code{ramfs}, \code{vfs}, \code{proc}, \code{task}, \code{user} e
\code{arch}. \\
\code{tty.rs} & Mant\'em os oito consoles; usa \code{vga}, \code{keyboard},
\code{serial} e \code{sync}. \\
\code{vga.rs} & Renderizador de baixo n\'ivel; usa \code{arch}. \\
\code{keyboard.rs} & Usa \code{arch} para acessar o 8042. \\
\code{serial.rs} & Usa \code{arch} para acessar a UART. \\
\code{sync.rs} & Spinlocks seguros para IRQ no console, ramfs, heap, VFS e
tarefas; usa \code{core} e \code{arch}. \\
\bottomrule
\end{tabularx}
\end{table}

A tabela mostra apenas as depend\^encias relevantes. \code{boot.rs} \'e um
caso especial: ele cont\'em o c\'odigo de montagem que chama \code{kernel\_main}
e n\~ao \'e chamado por Rust. \code{main.rs} declara todos os m\'odulos e
exporta as macros \code{print!} e \code{println!}.

\section{Cadeia de boot}

\begin{figure}[H]
\centering
\begin{tikzpicture}[
  stage/.style={draw=rinuxblue, fill=rinuxlight, rounded corners=3pt,
                minimum width=3.0cm, minimum height=1.25cm,
                align=center, font=\scriptsize},
  arrow/.style={-{Latex[length=2mm]}, thick, rinuxblue},
  node distance=0.55cm
]
\node[stage] (bios) {BIOS\\SeaBIOS};
\node[stage, right=of bios] (eltorito) {El Torito\\CD};
\node[stage, right=of eltorito] (grub) {GRUB\\i386-pc};
\node[stage, right=of grub] (mb) {Multiboot2\\kernel.elf};
\node[stage, below=1.0cm of mb] (stub) {\_start\\32 bits @ 1 MiB};
\node[stage, left=of stub] (paging) {PML4/PDPT/PD\\GDT + long mode};
\node[stage, left=of paging] (main) {kernel\_main\\metade alta};
\node[stage, left=of main] (shell) {vga / tty / kbd\\shell};
\draw[arrow] (bios) -- (eltorito);
\draw[arrow] (eltorito) -- (grub);
\draw[arrow] (grub) -- (mb);
\draw[arrow] (mb) -- (stub);
\draw[arrow] (stub) -- (paging);
\draw[arrow] (paging) -- (main);
\draw[arrow] (main) -- (shell);
\end{tikzpicture}
\caption{Cadeia de boot completa.}
\label{fig:bootchain}
\end{figure}

\section{Decis\~oes de projeto}

\begin{table}[H]
\centering
\caption{Decis\~oes arquiteturais e justificativas.}
\label{tab:decisoes}
\begin{tabularx}{\linewidth}{@{}p{3.3cm}X@{}}
\toprule
\textbf{Decis\~ao} & \textbf{Justificativa} \\
\midrule
\code{no\_std} & Elimina depend\^encia de sistema operacional, alocador e
\emph{runtime}; o c\'odigo usa apenas \code{core}. \\
Sem depend\^encias & Reduz superf\'icie de ataque e mant\'em o build
reprodut\'ivel; \code{Cargo.lock} n\~ao tem pacotes externos. \\
Alvo \code{x86\_64-unknown-none} & Alvo oficial bare metal, com modelo de
c\'odigo \emph{kernel} e sem biblioteca padr\~ao. \\
Kernel na metade alta & Conven\c{c}\~ao do Linux e do alvo; evita depender de
endere\c{c}os baixos e deixa o stub separado. \\
Stub em assembly & O controle inicial vem em 32 bits com pilha indefinida;
a transi\c{c}\~ao para 64 bits exige c\'odigo posicional e controle exato. \\
Tabelas de p\'agina est\'aticas & O stub n\~ao tem alocador; 3 p\'aginas de
4 KiB s\~ao suficientes para mapear 1 GiB com \emph{huge pages} de 2 MiB.
Depois do boot, \code{paging.rs} aloca tabelas de 4 KiB para as regi\~oes sob
demanda. \\
Polling no teclado & Mant\'em IRQ1 mascarada e evita um buffer de scancodes;
o timer IRQ0 j\'a demonstra o caminho de interrup\c{c}\~ao. \\
Serial al\'em do VGA & Permite depura\c{c}\~ao e testes automatizados sem
depender de captura gr\'afica. \\
\emph{Panic abort} & N\~ao h\'a \emph{unwinding} nem runtime de p\^anico;
\code{panic=abort} reduz c\'odigo gerado. \\
LTO e \code{codegen-units=1} & Gera um bin\'ario pequeno e previs\'ivel; o
linker descarta se\c{c}\~oes n\~ao usadas com \code{--gc-sections}. \\
\bottomrule
\end{tabularx}
\end{table}

\section{Tratamento de erros e falhas}

O kernel n\~ao possui um subsistema de erros. As estrat\'egias s\~ao:

\begin{itemize}
  \item \textbf{Timeout em esperas de hardware}: as fun\c{c}\~oes do 8042
        tentam por um n\'umero limitado de itera\c{c}\~oes e desistem, evitando
        \emph{hang} quando o controlador n\~ao responde.
  \item \textbf{Verifica\c{c}\~ao de presen\c{c}a da UART}: se o registrador de
        status devolver \code{0xFF}, a serial \'e considerada ausente.
  \item \textbf{Entrada inv\'alida}: bytes fora da faixa imprim\'ivel s\~ao
        ignorados no shell; UTF-8 inv\'alido em \code{echo} vira uma
        mensagem fixa.
  \item \textbf{Panic}: o handler imprime \code{[PANIC]} no console e entra em
        \code{halt\_loop()}. Exce\c{c}\~oes de CPU passam pela IDT; um
        \code{\#PF} sob demanda pode ser resolvido, mas uma falha fora das
        regi\~oes mapeadas imprime diagn\'ostico e para a CPU. Um
        \code{\#DF} usa a pilha IST e tamb\'em para o sistema.
  \item \textbf{Sem recupera\c{c}\~ao}: n\~ao h\'a \emph{kernel oops},
        reinicializa\c{c}\~ao controlada ou despejo de pilha.
\end{itemize}

\section{Concorr\^encia e sincroniza\c{c}\~ao}

Apesar de o sistema ser monoprocessado e rodar com interrup\c{c}\~oes
desligadas, o console \'e protegido por um \code{SpinLock}. Isso permite que
chamadas de impress\~ao vindas de contextos diferentes no futuro n\~ao
intercalem bytes no buffer VGA. O lock usa \code{AtomicBool::swap} com
ordena\c{c}\~ao \code{Acquire} na aquisi\c{c}\~ao e \code{Release} na
libera\c{c}\~ao. A serial n\~ao \'e protegida pelo lock; como o console
escreve VGA e serial em sequ\^encia, uma futura implementa\c{c}\~ao
multithread precisaria de um lock adicional para a UART.

\section{Ciclo de vida da inicializa\c{c}\~ao}

\begin{table}[H]
\centering
\caption{Ordem de inicializa\c{c}\~ao em \code{kernel\_main}.}
\label{tab:init}
\begin{tabular}{@{}clp{7.5cm}@{}}
\toprule
\textbf{\#} & \textbf{Chamada} & \textbf{Efeito} \\
\midrule
1 & \code{vga::init()} & Limpa a tela, zera o cursor, define forma do cursor. \\
2 & \code{serial::init()} & Configura COM1 em 38400 8N1 e detecta presen\c{c}a. \\
3 & \code{keyboard::init()} & Configura o 8042, tradu\c{c}\~ao e varredura. \\
4 & \code{println!} & Banner e mensagem de boas-vindas (VGA + serial). \\
5 & valida\c{c}\~ao & Confere \code{magic == 0x36d76289}. \\
6 & \code{shell::run()} & La\c{c}o eterno de prompt e leitura de comandos. \\
\bottomrule
\end{tabular}
\end{table}

\chapter{Boot e inicializa\c{c}\~ao de baixo n\'ivel}
\label{cap:boot}

\section{Responsabilidades de \code{boot.rs}}

O m\'odulo \code{src/boot.rs} concentra tudo o que precisa acontecer antes de
existir uma pilha Rust, mem\'oria virtual ou console:

\begin{enumerate}
  \item declarar o cabe\c{c}alho Multiboot2 reconhecido pelo GRUB;
  \item salvar os argumentos passados pelo bootloader;
  \item construir as tabelas de p\'agina;
  \item carregar uma GDT de 64 bits;
  \item habilitar PAE, \emph{long mode} e pagina\c{c}\~ao;
  \item transferir o controle para \code{kernel\_main} na metade alta.
\end{enumerate}

O arquivo combina uma estrutura Rust est\'atica (o cabe\c{c}alho) e um bloco
\code{global\_asm!} com o restante. N\~ao h\'a fun\c{c}\~oes Rust nesse
m\'odulo.

\section{Cabe\c{c}alho Multiboot2}

O GRUB procura a assinatura \code{0xE85250D6} nos primeiros 32 KiB do
arquivo. O cabe\c{c}alho \'e declarado em uma \code{static} alinhada a 8 bytes
e colocada na se\c{c}\~ao \code{.multiboot\_header}.

\begin{table}[H]
\centering
\caption{Campos do cabe\c{c}alho Multiboot2 implementado.}
\label{tab:mbheader}
\begin{tabular}{@{}rrl@{}}
\toprule
\textbf{Valor} & \textbf{Offset} & \textbf{Significado} \\
\midrule
\code{0xE85250D6} & 0x00 & Assinatura Multiboot2 \\
\code{0x00000000} & 0x04 & Arquitetura i386 (32 bits) \\
\code{0x00000018} & 0x08 & Tamanho total do cabe\c{c}alho (24 bytes) \\
\code{0x17ADAF12} & 0x0C & Checksum \\
\code{0x00000000} & 0x10 & Tag final: tipo 0 \\
\code{0x00000008} & 0x14 & Tag final: tamanho 8 \\
\bottomrule
\end{tabular}
\end{table}

O checksum \'e calculado de modo que a soma dos quatro primeiros inteiros de
32 bits seja zero:

\begin{verbatim}
-(0xE85250D6 + 0x00000000 + 0x00000018) = 0x17ADAF12
\end{verbatim}

A tag final \'e obrigat\'oria. N\~ao h\'a tags de informa\c{c}\~ao de
v\'ideo, mem\'oria ou m\'odulos; o kernel recebe apenas o ponteiro para a
estrutura Multiboot2 e ignora seu conte\'udo, exceto pelo valor de
\code{magic}, que \'e validado e impresso.

\section{Contrato de entrada}

Quando o GRUB transfere o controle para \code{\_start}:

\begin{itemize}
  \item a CPU est\'a em modo protegido de 32 bits;
  \item \code{eax} cont\'em \code{0x36D76289} (assinatura Multiboot2);
  \item \code{ebx} cont\'em o endere\c{c}o f\'isico da estrutura de
        informa\c{c}\~oes Multiboot2;
  \item a pilha n\~ao \'e confi\'avel;
  \item as interrup\c{c}\~oes podem estar em qualquer estado.
\end{itemize}

O stub imediatamente executa \code{cli} e \code{cld}, define uma pilha
tempor\'aria de 8 KiB em \code{boot\_stack\_top} e zera \code{ebp}.

\section{Sequ\^encia do stub}

\begin{enumerate}
  \item \textbf{Salvar argumentos}: \code{eax} e \code{ebx} s\~ao copiados
        para \code{mb\_magic} e \code{mb\_info}, vari\'aveis de 32 bits em
        \code{.boot.bss}. Como o c\'odigo ainda roda em endere\c{c}os baixos,
        esses endere\c{c}os s\~ao acess\'iveis diretamente.
  \item \textbf{Zerar p\'aginas}: \code{rep stosl} limpa 3 $\times$ 4096
        bytes (PML4, PDPT e PD).
  \item \textbf{Montar PML4}: a entrada 0 aponta para a PDPT e a entrada 511
        tamb\'em aponta para a mesma PDPT. A entrada 511 cobre a faixa
        \code{0xFFFFFF8000000000} em diante.
  \item \textbf{Montar PDPT}: a entrada 0 e a entrada 510 apontam para a
        mesma PD. A entrada 510, dentro do PML4[511], resulta em
        \code{0xFFFFFFFF80000000}.
  \item \textbf{Montar PD}: 512 entradas de 2 MiB, cada uma com os bits
        \emph{present} e \emph{page size} (\code{0x83}),
        mapeando fisicamente 0 a 1 GiB.
  \item \textbf{Carregar GDT}: \code{lgdt gdt64\_ptr}.
  \item \textbf{Preparar a CPU}: limpar \code{CR0.EM}, setar \code{CR0.MP},
        setar \code{CR4.PAE}, \code{CR4.OSFXSR} e \code{CR4.OSXMMEXCPT}.
  \item \textbf{Ativar long mode}: setar \code{EFER.LME} via \code{rdmsr} /
        \code{wrmsr} e, em seguida, \code{CR0.PG | CR0.PE}.
  \item \textbf{Salto de modo}: \code{ljmp \$0x08, \$long\_mode\_entry\_low},
        carregando um seletor de c\'odigo com o bit \code{L} ligado.
  \item \textbf{Modo 64 bits}: recarregar os segmentos de dados com
        \code{0x10}, carregar \code{rsp} com \code{stack\_top} (metade alta),
        preparar \code{edi}/\code{esi} e chamar \code{kernel\_main}.
\end{enumerate}

\begin{nota}[Por que duas pilhas?]
A pilha do stub precisa estar em endere\c{c}o baixo, porque a CPU ainda n\~ao
usa pagina\c{c}\~ao quando o stub come\c{c}a. A pilha definitiva do Rust fica
em \code{.bss}, na metade alta, e s\'o pode ser usada depois que o modo longo
e a pagina\c{c}\~ao est\~ao ativos.
\end{nota}

\section{Pagina\c{c}\~ao e layout de mem\'oria}

A tradu\c{c}\~ao de endere\c{c}os usa quatro n\'iveis, mas apenas tr\^es
tabelas s\~ao preenchidas. O page directory usa \emph{huge pages} de 2 MiB,
reduzindo o consumo a 12 KiB de tabelas.

\begin{figure}[H]
\centering
\begin{tikzpicture}[
  tab/.style={draw=rinuxblue, fill=rinuxlight, rounded corners=2pt,
              minimum width=3.3cm, minimum height=0.9cm,
              align=center, font=\scriptsize\ttfamily},
  box/.style={draw=black!50, fill=white, rounded corners=2pt,
              minimum width=3.8cm, minimum height=0.9cm,
              align=center, font=\scriptsize},
  arrow/.style={-{Latex[length=2mm]}, thick, rinuxblue},
  node distance=0.55cm and 0.55cm
]
\node[tab] (pml4) {PML4\\[0]=0x...003\\{[511]=0x...003}};
\node[tab, right=of pml4] (pdpt) {PDPT\\{[0]=PD, [510]=PD}};
\node[tab, right=of pdpt] (pd) {PD\\{512 x 2 MiB}};
\node[box, below=1.2cm of pd] (low) {Identidade\\0x00000000--0x3FFFFFFF};
\node[box, below=1.2cm of pml4] (high) {Metade alta\\0xFFFFFFFF80000000\\--0xFFFFFFFFBFFFFFFF};
\draw[arrow] (pml4) -- node[above,font=\tiny]{PDPT} (pdpt);
\draw[arrow] (pdpt) -- node[above,font=\tiny]{PD} (pd);
\draw[arrow] (pdpt.south) -- ++(0,-0.35) -| (low.north);
\draw[arrow] (pdpt.south) -- ++(0,-0.35) -| (high.north);
\end{tikzpicture}
\caption{Estrutura de pagina\c{c}\~ao: uma \'unica PDPT e uma \'unica PD
servem tanto \`a identidade baixa quanto \`a metade alta.}
\label{fig:paging}
\end{figure}

\subsection{C\'alculo dos \'indices}

Para um endere\c{c}o virtual de 48 bits, os \'indices s\~ao:

\begin{align*}
\text{PML4} &= \left\lfloor \frac{\text{VA}}{2^{39}} \right\rfloor \bmod 512,\\
\text{PDPT} &= \left\lfloor \frac{\text{VA}}{2^{30}} \right\rfloor \bmod 512,\\
\text{PD}   &= \left\lfloor \frac{\text{VA}}{2^{21}} \right\rfloor \bmod 512.
\end{align*}

Para \code{0xFFFFFFFF80000000}:

\begin{itemize}
  \item PML4 = 511 (\code{0xFFFFFFFF80000000 >> 39 \& 0x1FF});
  \item PDPT = 510 (\code{>> 30 \& 0x1FF});
  \item PD = 0.
\end{itemize}

Assim, a entrada 510 da PDPT aponta para a PD que descreve os primeiros 1 GiB
f\'isicos. O mesmo conjunto de tabelas atende aos dois mapeamentos:

\begin{itemize}
  \item \textbf{identidade baixa}: \code{0x00000000}--\code{0x3FFFFFFF}
        (stub, tabelas, estruturas do GRUB);
  \item \textbf{metade alta}: \code{0xFFFFFFFF80000000}--
        \code{0xFFFFFFFFBFFFFFFF} (c\'odigo e dados do kernel).
\end{itemize}

\begin{atencao}[Limite de 1 GiB]
O mapeamento cobre exatamente 1 GiB. Se o kernel ou o \emph{framebuffer}
estiverem acima disso, ou se a m\'aquina tiver mais de 1 GiB e o kernel
precisar acessar mem\'oria al\'em desse limite, o mapeamento precisa ser
estendido (mais entradas na PD ou tabelas adicionais).
\end{atencao}

\section{GDT e modo longo}

A GDT cont\'em tr\^es descritores:

\begin{table}[H]
\centering
\caption{Descritores da GDT de 64 bits.}
\label{tab:gdt}
\begin{tabular}{@{}clll@{}}
\toprule
\textbf{Seletor} & \textbf{Valor} & \textbf{Tipo} & \textbf{Uso} \\
\midrule
\code{0x00} & \code{0x0000000000000000} & Nulo & Obrigat\'orio \\
\code{0x08} & \code{0x00AF9A000000FFFF} & C\'odigo ring 0, L=1 & \code{CS} \\
\code{0x10} & \code{0x00CF92000000FFFF} & Dados ring 0 & \code{DS/ES/FS/GS/SS} \\
\bottomrule
\end{tabular}
\end{table}

O descritor de c\'odigo tem o bit \code{L} (long mode) ligado, indicando
c\'odigo de 64 bits. O descritor de dados \'e o tradicional de 32 bits, usado
apenas para os registradores de segmento; em modo longo, a segmenta\c{c}\~ao
de dados \'e praticamente ignorada.

\section{Registradores de controle}

\begin{table}[H]
\centering
\caption{Configura\c{c}\~ao dos registradores de controle.}
\label{tab:cr}
\begin{tabular}{@{}llp{7.5cm}@{}}
\toprule
\textbf{Registrador} & \textbf{Bits} & \textbf{Efeito} \\
\midrule
\code{CR0} & \code{EM=0}, \code{MP=1} & Habilita SSE/FXSR sem emula\c{c}\~ao de FPU. \\
\code{CR0} & \code{PG=1}, \code{PE=1} & Ativa pagina\c{c}\~ao e modo protegido. \\
\code{CR3} & \code{pml4} & Aponta para o PML4 f\'isico. \\
\code{CR4} & \code{PAE=1} & Habilita endere\c{c}os f\'isicos de 36 bits e p\'aginas de 2 MiB. \\
\code{CR4} & \code{OSFXSR=1} & Sistema operacional suporta \code{FXSAVE}/\code{FXRSTOR}. \\
\code{CR4} & \code{OSXMMEXCPT=1} & Exce\c{c}\~oes SIMD entregues ao SO. \\
\code{EFER} & \code{LME=1} & Habilita \emph{long mode}. \\
\bottomrule
\end{tabular}
\end{table}

\section{Linker script}

O arquivo \code{linker.ld} define dois mundos:

\begin{enumerate}
  \item \code{.boot}, carregado em \code{0x100000} (1 MiB), cont\'em o
        cabe\c{c}alho Multiboot2, o c\'odigo de 32/64 bits, a GDT, as tabelas
        de p\'agina e a pilha do stub;
  \item \code{.text}, \code{.rodata}, \code{.data} e \code{.bss}, ligados em
        \code{0xFFFFFFFF80000000} com \emph{load address} f\'isico
        \code{0x00100000 + deslocamento}.
\end{enumerate}

A macro \code{KERNEL\_VMA} vale \code{0xFFFFFFFF80000000}. Para cada
se\c{c}\~ao alta, o script usa \code{AT(ADDR(.secao) - KERNEL\_VMA)} para
informar ao linker que o endere\c{c}o de carga f\'isico \'e baixo, embora o
endere\c{c}o virtual seja alto.

\begin{lstlisting}[caption={Trecho central de \code{linker.ld}.},label={lst:linker}]
ENTRY(_start)
KERNEL_VMA = 0xFFFFFFFF80000000;

SECTIONS
{
    . = 1M;

    .boot : ALIGN(8)
    {
        KEEP(*(.multiboot_header))
        *(.text.boot)
        *(.boot.rodata)
        *(.boot.data)
        *(.boot.bss)
    }

    . = ALIGN(4096);
    . += KERNEL_VMA;

    .text : AT(ADDR(.text) - KERNEL_VMA)
    {
        *(.text .text.*)
    }
    /* .rodata, .data e .bss seguem o mesmo padrao */
}
\end{lstlisting}

Se\c{c}\~oes de metadados que n\~ao s\~ao necess\'arias em um kernel
(\code{.eh\_frame}, \code{.comment}, \code{.note}) s\~ao descartadas com
\code{/DISCARD/}, e o build usa \code{--gc-sections} para remover c\'odigo
n\~ao referenciado.

\section{Se\c{c}\~oes observadas no ELF}

A compila\c{c}\~ao de refer\^encia gerou o seguinte layout:

\begin{table}[H]
\centering
\caption{Se\c{c}\~oes do ELF de refer\^encia (endere\c{c}os virtuais).}
\label{tab:elf}
\begin{tabular}{@{}llrr@{}}
\toprule
\textbf{Se\c{c}\~ao} & \textbf{VMA} & \textbf{Offset} & \textbf{Tamanho} \\
\midrule
\code{.boot}   & \code{0x0000000000100000} & \code{0x1000} & \code{0x6008} \\
\code{.text}   & \code{0xFFFFFFFF80107000} & \code{0x8000} & \code{0x2516} \\
\code{.rodata} & \code{0xFFFFFFFF80109518} & \code{0xA518} & \code{0x0AE7} \\
\code{.data}   & \code{0xFFFFFFFF8010A000} & \code{0xB000} & \code{0x0020} \\
\code{.bss}    & \code{0xFFFFFFFF8010A020} & \code{0xB020} & \code{0x8003} \\
\bottomrule
\end{tabular}
\end{table}

O ponto de entrada \'e \code{0x100018}, dentro de \code{.boot}. O ELF tem
6 \emph{program headers}: um segmento baixo \code{RWE} para o stub e quatro
segmentos altos (c\'odigo, read-only, dados e BSS).

\section{Passagem para Rust}

O \'ultimo trecho do stub carrega os argumentos e chama a fun\c{c}\~ao Rust:

\begin{lstlisting}[caption={Chamada de \code{kernel\_main} a partir do stub.},label={lst:callmain}]
    mov mb_magic(%rip), %edi
    mov mb_info(%rip), %esi
    movabs $kernel_main, %rax
    call *%rax

3:
    cli
    hlt
    jmp 3b
\end{lstlisting}

A fun\c{c}\~ao Rust usa a conven\c{c}\~ao \code{extern "C"}: os dois primeiros
argumentos inteiros v\~ao em \code{edi} e \code{esi}. O retorno \'e
\code{!}, mas o stub ainda mant\'em um la\c{c}o de seguran\c{c}a com
\code{cli; hlt} caso a chamada retorne inesperadamente.

\section{Limita\c{c}\~oes do boot}

\begin{itemize}
  \item N\~ao h\'a suporte a UEFI nativo; depende de BIOS/CSM e de VGA
        legado.
  \item N\~ao h\'a real-mode setup, \emph{protected mode} pr\'oprio nem
        descompress\~ao; o GRUB faz o carregamento.
  \item O \emph{stub} n\~ao solicita informa\c{c}\~oes de mem\'oria,
        framebuffer ou m\'odulos. A estrutura Multiboot2 s\'o \'e interpretada
        depois, em Rust, por \code{paging::init}, para liberar frames.
  \item O \emph{stub} mapeia apenas 1 GiB por identidade e n\~ao instala
        demand paging, prote\c{c}\~ao \code{NX} ou \emph{guard pages}. Essas
        estruturas s\~ao criadas depois por \code{paging.rs}.
  \item O \emph{stub} n\~ao instala IDT; at\'e \code{idt::init()} rodar,
        qualquer exce\c{c}\~ao (incluindo \emph{page fault}) resulta em
        \emph{triple fault}.
\end{itemize}

\chapter{M\'odulos de base}
\label{cap:base}

Este cap\'itulo documenta os m\'odulos que n\~ao pertencem a um driver
espec\'ifico: \code{main.rs}, \code{arch.rs} e \code{sync.rs}. Eles fornecem o
ponto de entrada, as macros de impress\~ao, o acesso a portas de I/O, o
desligamento/reinicializa\c{c}\~ao e o spinlock do console.

\section{\code{main.rs}: ponto de entrada Rust}

\subsection{Atributos globais}

\begin{lstlisting}[caption={Atributos e m\'odulos declarados em \code{main.rs}.},label={lst:mainattrs}]
#![no_std]
#![no_main]

mod arch;
mod boot;
mod gdt;
mod idt;
mod keyboard;
mod paging;
mod serial;
mod shell;
mod sync;
mod user;
mod vga;
\end{lstlisting}

\begin{itemize}
  \item \code{no\_std}: o crate n\~ao usa a biblioteca padr\~ao; n\~ao h\'a
        \code{std::io}, \code{String}, \code{Vec}, sistema de arquivos nem
        \emph{threads}.
  \item \code{no\_main}: n\~ao existe a fun\c{c}\~ao \code{main} do runtime
        Rust. O ponto de entrada real \'e \code{\_start}, definido no
        assembly de \code{boot.rs}.
  \item \code{mod boot}: necess\'ario para que o linker inclua o c\'odigo de
        montagem e o cabe\c{c}alho Multiboot2.
\end{itemize}

\subsection{\code{kernel\_main}}

\begin{lstlisting}[caption={Assinatura e corpo principal de \code{kernel\_main}.},label={lst:kernelmain}]
#[unsafe(no_mangle)]
pub extern "C" fn kernel_main(multiboot_magic: u32, multiboot_info: u32) -> ! {
    vga::init();
    serial::init();
    paging::init(multiboot_magic, multiboot_info);
    gdt::init();
    keyboard::init();
    idt::init();
    arch::enable_interrupts();

    println!();
    println!("==============================================================");
    println!("  Rust VGA Kernel - x86_64 / modo longo");
    /* ... banner ... */

    if multiboot_magic != 0x36d7_6289 {
        println!("AVISO: assinatura multiboot2 invalida: {:#010x}",
                 multiboot_magic);
    }
    println!("Multiboot2: magic={:#010x} info={:#010x}",
             multiboot_magic, multiboot_info);
    println!();
    println!("Digite 'help' para ver os comandos.");
    println!();

    shell::run()
}
\end{lstlisting}

\begin{table}[H]
\centering
\caption{Contrato de \code{kernel\_main}.}
\label{tab:kernelmain}
\begin{tabularx}{\linewidth}{@{}lX@{}}
\toprule
\textbf{Aspecto} & \textbf{Descri\c{c}\~ao} \\
\midrule
Visibilidade & \code{pub extern "C"}, exportada sem \emph{mangling}. \\
Argumentos & \code{multiboot\_magic} (deve ser \code{0x36D76289}) e
\code{multiboot\_info} (endere\c{c}o f\'isico da estrutura Multiboot2). \\
Retorno & \code{!}: nunca retorna; termina em \code{shell::run()}. \\
Pr\'e-condi\c{c}\~oes & CPU em modo longo, pagina\c{c}\~ao ativa, pilha alta
v\'alida, \code{.bss} zerada pelo carregador. \\
Efeitos & Inicializa VGA, serial, alocador de frames, GDT/TSS, teclado,
IDT/PIC/PIT e liga interrup\c{c}\~oes; imprime banner; entra no shell. \\
Chamado por & Stub de assembly em \code{boot.rs}. \\
\bottomrule
\end{tabularx}
\end{table}

O valor \code{multiboot\_info} \'e usado por \code{paging::init} para
percorrer as tags Multiboot2 e liberar os frames de RAM dispon\'ivel. A
estrutura tamb\'em \'e reservada para n\~ao ser entregue pelo alocador.

\subsection{Macros \code{print!} e \code{println!}}

As macros s\~ao definidas com \code{macro\_rules!} e \code{\$crate}, para
funcionar em qualquer m\'odulo do crate sem importa\c{c}\~ao:

\begin{lstlisting}[caption={Macros de impress\~ao sem aloca\c{c}\~ao.},label={lst:macros}]
#[macro_export]
macro_rules! print {
    ($($arg:tt)*) => {
        $crate::tty::_print(format_args!($($arg)*))
    };
}

#[macro_export]
macro_rules! println {
    () => { $crate::print!("\n") };
    ($($arg:tt)*) => {
        $crate::print!("{}\n", format_args!($($arg)*))
    };
}
\end{lstlisting}

\code{format\_args!} cria uma descri\c{c}\~ao de formata\c{c}\~ao sem alocar;
a implementa\c{c}\~ao de \code{fmt::Write} em \code{tty.rs} consome os
fragmentos, anexa ao buffer do console ativo e espelha na serial.
N\~ao h\'a \code{format!} (que exigiria \code{alloc}).

\subsection{Panic handler}

\begin{lstlisting}[caption={\emph{Panic handler} do kernel.},label={lst:panic}]
#[panic_handler]
fn panic(info: &PanicInfo) -> ! {
    crate::println!();
    crate::println!("[PANIC] {}", info);
    crate::arch::halt_loop()
}
\end{lstlisting}

Como o perfil usa \code{panic = "abort"}, o handler n\~ao precisa de
\emph{unwinding}. Ele imprime a mensagem no console e para a CPU. Se o
p\^anico ocorrer antes de \code{vga::init()}, a sa\'ida pode n\~ao ser
vis\'ivel; se ocorrer dentro do console com o spinlock adquirido, pode
ocorrer \emph{deadlock} (o handler tentaria imprimir e esperaria o lock).

\section{\code{arch.rs}: arquitetura e portas de I/O}

\subsection{Portas de I/O}

O m\'odulo \code{arch::port} encapsula as instru\c{c}\~oes \code{in} e
\code{out} do x86:

\begin{lstlisting}[caption={Leitura e escrita de portas de 8 bits.},label={lst:portio}]
pub unsafe fn inb(port: u16) -> u8 {
    let value: u8;
    asm!("in al, dx",
         in("dx") port,
         out("al") value,
         options(nomem, nostack, preserves_flags));
    value
}

pub unsafe fn outb(port: u16, value: u8) {
    asm!("out dx, al",
         in("dx") port,
         in("al") value,
         options(nomem, nostack, preserves_flags));
}
\end{lstlisting}

As fun\c{c}\~oes s\~ao \code{unsafe} porque qualquer porta pode n\~ao existir
ou pode ter efeitos colaterais no hardware. As op\c{c}\~oes \code{nomem},
\code{nostack} e \code{preserves\_flags} informam ao compilador que a
instru\c{c}\~ao n\~ao altera mem\'oria, pilha ou flags, permitindo otimiza\c{c}\~oes
e evitando que o bloco seja removido.

\begin{table}[H]
\centering
\caption{API de \code{arch}.}
\label{tab:archapi}
\begin{tabularx}{\linewidth}{@{}lX@{}}
\toprule
\textbf{Fun\c{c}\~ao} & \textbf{Descri\c{c}\~ao} \\
\midrule
\code{inb(port)} & L\^e um byte de uma porta de I/O. \code{unsafe}. \\
\code{outb(port, value)} & Escreve um byte em uma porta. \code{unsafe}. \\
\code{halt\_loop()} & Executa \code{cli; hlt} em la\c{c}o infinito; n\~ao
retorna. Usado no panic e no comando \code{halt}. \\
\code{reboot()} & Espera o buffer de entrada do 8042 esvaziar e envia o
comando \code{0xFE} (reset). Se n\~ao funcionar, cai em \code{halt\_loop}. \\
\bottomrule
\end{tabularx}
\end{table}

\subsection{Reinicializa\c{c}\~ao}

O reset usa o controlador 8042, que historicamente controla a linha de reset
da CPU. O c\'odigo espera at\'e 100.000 itera\c{c}\~oes pelo bit
\code{INPUT\_FULL} (\code{0x02}) do status limpar e escreve \code{0xFE} na
porta de comando (\code{0x64}). Em m\'aquinas virtualizadas ou sem 8042, o
reset pode n\~ao ocorrer; nesse caso a CPU \'e parada.

\section{\code{sync.rs}: spinlock}

\subsection{Implementa\c{c}\~ao}

\begin{lstlisting}[caption={Spinlock com \code{AtomicBool}.},label={lst:spinlock}]
pub struct SpinLock<T> {
    locked: AtomicBool,
    data: UnsafeCell<T>,
}

unsafe impl<T: Send> Sync for SpinLock<T> {}
unsafe impl<T: Send> Send for SpinLock<T> {}

impl<T> SpinLock<T> {
    pub const fn new(data: T) -> Self { /* ... */ }

    pub fn lock(&self) -> SpinLockGuard<'_, T> {
        while self.locked.swap(true, Ordering::Acquire) {
            core::hint::spin_loop();
        }
        SpinLockGuard { lock: self }
    }
}
\end{lstlisting}

O \emph{guard} implementa \code{Deref}, \code{DerefMut} e \code{Drop}. Ao ser
destru\'ido, armazena \code{false} com ordena\c{c}\~ao \code{Release}. O
spinlock \'e usado apenas para o console:

\begin{lstlisting}[caption={Console global protegido por spinlock.},label={lst:consoleglobal}]
static CONSOLE: SpinLock<Console> = SpinLock::new(Console::new());
\end{lstlisting}

Como \code{Console::new()} \'e \code{const}, o lock pode ser inicializado
estaticamente sem c\'odigo de inicializa\c{c}\~ao em tempo de execu\c{c}\~ao.

\subsection{Propriedades e limita\c{c}\~oes}

\begin{itemize}
  \item \textbf{Bloqueio girat\'orio}: se outra CPU ou interrup\c{c}\~ao
        mantiver o lock, a CPU atual gira at\'e liberar. N\~ao h\'a
        \emph{sleep} nem timeout.
  \item \textbf{Sem IRQ save/restore}: desabilitar interrup\c{c}\~oes ao
        adquirir o lock n\~ao \'e feito. Com IRQs desligadas no boot, isso
        n\~ao \'e problema hoje; ao habilitar IRQs, um handler que use o
        console pode causar \emph{deadlock}.
  \item \textbf{Sem detec\c{c}\~ao de recurs\~ao}: adquirir o mesmo lock duas
        vezes na mesma CPU trava.
  \item \textbf{Ordem de mem\'oria}: \code{Acquire}/\code{Release} garante
        que escritas ao console n\~ao sejam reordenadas para fora da
        regi\~ao cr\'itica em arquiteturas com mem\'oria fraca; no x86,
        o modelo \'e mais forte, mas a anota\c{c}\~ao mant\'em o c\'odigo
        correto.
\end{itemize}

\begin{nota}[Escopo atual]
Com uma CPU, interrup\c{c}\~oes desligadas e apenas o shell chamando
\code{print!}, o spinlock nunca fica contended. Ele existe como prepara\c{c}\~ao
para extens\~oes e para documentar a inten\c{c}\~ao de exclus\~ao m\'utua.
\end{nota}

\chapter{Console VGA em modo texto}
\label{cap:vga}

\section{Vis\~ao geral}

O driver \code{src/vga.rs} \'e o equivalente simplificado do \code{vgacon} do
Linux. Ele opera no modo texto legado da VGA, no qual a tela \'e uma matriz
de 80 colunas por 25 linhas armazenada a partir do endere\c{c}o f\'isico
\code{0xB8000}. Cada c\'elula ocupa 16 bits.

\begin{table}[H]
\centering
\caption{Formato de uma c\'elula de texto VGA.}
\label{tab:vgacell}
\begin{tabular}{@{}lll@{}}
\toprule
\textbf{Bits} & \textbf{Campo} & \textbf{Valor usado} \\
\midrule
0--7  & Caractere ASCII & C\'odigo do caractere \\
8--11 & Cor de primeiro plano & \code{0x7} (cinza claro) \\
12--14 & Cor de fundo & \code{0x0} (preto) \\
15    & Piscar & \code{0} (desligado) \\
\midrule
\multicolumn{3}{@{}l@{}}{\textbf{Atributo completo}: \code{0x07}} \\
\bottomrule
\end{tabular}
\end{table}

O buffer tem $80 \times 25 = 2.000$ c\'elulas, ou 4.000 bytes. A posi\c{c}\~ao
do cursor \'e controlada pelos registradores da CRTC, acessados pelas portas
\code{0x3D4} (\'indice) e \code{0x3D5} (dados).

\section{Estado do console}

\begin{lstlisting}[caption={Estado mantido pelo console.},label={lst:consolestate}]
struct Console {
    column: usize,
    row: usize,
    color: u8,
}
\end{lstlisting}

O estado \'e protegido por \code{static CONSOLE: SpinLock<Console>}. N\~ao
h\'a buffer de \emph{scrollback}, nem mem\'oria separada da tela: cada
altera\c{c}\~ao \'e escrita diretamente no \emph{framebuffer} com
\code{write\_volatile}.

\section{Opera\c{c}\~oes}

\subsection{Inicializa\c{c}\~ao}

\code{vga::init()}:

\begin{enumerate}
  \item adquire o lock do console;
  \item preenche as 2.000 c\'elulas com espa\c{c}o e atributo \code{0x07};
  \item zera coluna e linha;
  \item configura o cursor como uma linha inferior, escrevendo
        \code{0x0E} no registrador \code{0x0A} (cursor start) e
        \code{0x0F} no registrador \code{0x0B} (cursor end);
  \item atualiza a posi\c{c}\~ao do cursor.
\end{enumerate}

\subsection{Escrita de bytes}

\begin{lstlisting}[caption={Tratamento de caracteres de controle e imprim\'iveis.},label={lst:writebyte}]
fn write_byte(&mut self, byte: u8) {
    match byte {
        b'\n' => self.newline(),
        b'\r' => { self.column = 0; self.update_cursor(); }
        b'\t' => {
            let next = (self.column + 8) & !7;
            if next >= BUFFER_WIDTH { self.newline(); }
            else { self.column = next; self.update_cursor(); }
        }
        0x08 => self.backspace(),
        0x20..=0x7E => self.write_char(byte as char),
        _ => self.write_char('?'),
    }
}
\end{lstlisting}

Caracteres imprim\'iveis ASCII s\~ao escritos diretamente. Qualquer outro
byte vira \code{?}; isso inclui acentos em UTF-8, que chegam como bytes
individuais e n\~ao s\~ao decodificados.

\subsection{Escrita de caracteres e quebra de linha}

\code{write\_char}:

\begin{itemize}
  \item se a coluna estiver em 80, chama \code{newline()};
  \item calcula o \'indice \code{row * 80 + column};
  \item grava \code{(ch as u16) | (color << 8)} com \code{write\_volatile};
  \item incrementa a coluna;
  \item se a coluna chegou a 80, chama \code{newline()}; caso contr\'ario,
        atualiza o cursor.
\end{itemize}

\code{newline()} zera a coluna e incrementa a linha. Se a linha atual for a
\'ultima (24), chama \code{scroll()}.

\subsection{Rolagem}

\code{scroll()} move cada c\'elula da linha $n$ para a linha $n-1$ e limpa a
\'ultima linha. O algoritmo \'e $O(80 \times 24)$ por linha adicional e n\~ao
usa mem\'oria tempor\'aria:

\begin{lstlisting}[caption={N\'ucleo da rolagem.},label={lst:scroll}]
for row in 1..BUFFER_HEIGHT {
    for column in 0..BUFFER_WIDTH {
        let src = row * BUFFER_WIDTH + column;
        let dst = (row - 1) * BUFFER_WIDTH + column;
        let cell = read_volatile(VGA_BUFFER.add(src));
        write_volatile(VGA_BUFFER.add(dst), cell);
    }
}
/* limpa a ultima linha */
\end{lstlisting}

Como origem e destino se sobrep\~uem e a c\'opia \'e feita de cima para baixo,
n\~ao h\'a perda de c\'elulas ainda n\~ao copiadas.

\subsection{Backspace}

O backspace recua a coluna; se estiver na coluna 0, passa para o fim da linha
anterior. Em seguida, grava uma c\'elula em branco na posi\c{c}\~ao e atualiza
o cursor. Ele n\~ao apaga al\'em da posi\c{c}\~ao (0,0).

\subsection{Cursor de hardware}

A posi\c{c}\~ao \'e calculada como \code{row * 80 + column} e dividida em dois
bytes:

\begin{lstlisting}[caption={Atualiza\c{c}\~ao do cursor via CRTC.},label={lst:cursor}]
outb(0x3D4, 0x0F);
outb(0x3D5, (position & 0xFF) as u8);
outb(0x3D4, 0x0E);
outb(0x3D5, (position >> 8) as u8);
\end{lstlisting}

O registrador \code{0x0E} recebe o byte alto e \code{0x0F} o byte baixo.

\section{Pipeline de impress\~ao}

\begin{figure}[H]
\centering
\begin{tikzpicture}[
  box/.style={draw=rinuxblue, fill=rinuxlight, rounded corners=2pt,
              minimum width=2.8cm, minimum height=1.0cm,
              align=center, font=\scriptsize},
  arrow/.style={-{Latex[length=2mm]}, thick, rinuxblue},
  node distance=0.55cm
]
\node[box] (macro) {\code{println!}};
\node[box, right=of macro] (args) {\code{format\_args!}};
\node[box, right=of args] (print) {\code{vga::\_print}};
\node[box, right=of print] (writer) {\code{ConsoleWriter}};
\node[box, below=0.8cm of writer] (serial) {\code{serial::write\_str}};
\node[box, left=of serial] (vga) {\code{Console::write\_string}};
\node[box, left=of vga] (fb) {VGA \code{0xB8000}};
\draw[arrow] (macro) -- (args);
\draw[arrow] (args) -- (print);
\draw[arrow] (print) -- (writer);
\draw[arrow] (writer) -- (vga);
\draw[arrow] (writer) -- (serial);
\draw[arrow] (vga) -- (fb);
\end{tikzpicture}
\caption{Fluxo de uma chamada de impress\~ao: VGA e serial recebem o mesmo texto.}
\label{fig:printpipeline}
\end{figure}

\code{ConsoleWriter} implementa \code{fmt::Write}. Ele adquire o lock,
escreve no VGA, libera o lock e s\'o ent\~ao escreve na serial. Isso evita
manter o lock durante a espera pela UART.

\section{API p\'ublica}

\begin{table}[H]
\centering
\caption{Fun\c{c}\~oes e constantes de \code{vga.rs}.}
\label{tab:vgaapi}
\begin{tabularx}{\linewidth}{@{}lX@{}}
\toprule
\textbf{Item} & \textbf{Descri\c{c}\~ao} \\
\midrule
\code{BUFFER\_WIDTH = 80} & Colunas da tela. \\
\code{BUFFER\_HEIGHT = 25} & Linhas da tela. \\
\code{init()} & Limpa a tela, configura forma e posi\c{c}\~ao do cursor. \\
\code{clear()} & Limpa a tela (comando \code{clear}). \\
\code{\_print(args)} & Ponto de entrada das macros; escreve VGA e serial. \\
\code{Console::clear()} & Preenche o buffer e reposiciona o cursor. \\
\code{Console::write\_byte()} & Interpreta controles e caracteres. \\
\code{Console::write\_char()} & Grava uma c\'elula e avan\c{c}a a coluna. \\
\code{Console::write\_string()} & Itera sobre os bytes de uma \code{\&str}. \\
\code{Console::newline()} & Avan\c{c}a linha, com rolagem se necess\'ario. \\
\code{Console::backspace()} & Retrocede e apaga uma c\'elula. \\
\code{Console::scroll()} & Move o conte\'udo uma linha para cima. \\
\code{Console::update\_cursor()} & Escreve a posi\c{c}\~ao na CRTC. \\
\bottomrule
\end{tabularx}
\end{table}

\section{Exemplo de uso}

\begin{lstlisting}[caption={Uso das macros no shell.},label={lst:usovga}]
crate::print!("rinux> ");
println!("Comandos disponiveis:");
crate::print!("\x08");
\end{lstlisting}

O primeiro comando imprime o prompt sem quebra de linha. O segundo imprime
uma linha completa. O terceiro envia o caractere de backspace para apagar o
\'ultimo caractere vis\'ivel.

\section{Limita\c{c}\~oes e observa\c{c}\~oes}

\begin{itemize}
  \item \textbf{Somente VGA texto}: n\~ao h\'a framebuffer gr\'afico, VBE,
        GOP ou console virtual (VT). Em UEFI sem CSM, \code{0xB8000} pode
        n\~ao existir.
  \item \textbf{ASCII apenas}: n\~ao h\'a decodifica\c{c}\~ao UTF-8 nem
        fonte bitmap pr\'opria; a fonte vem da VGA.
  \item \textbf{Sem atributos din\^amicos}: \code{color} existe no estado,
        mas n\~ao h\'a API p\'ublica para alter\'a-lo.
  \item \textbf{Sem sincroniza\c{c}\~ao com o retra\c{c}o vertical}: escritas
        podem causar \emph{tearing} ocasional, irrelevante para texto.
  \item \textbf{Sem prote\c{c}\~ao de mem\'oria}: \code{0xB8000} \'e acessado
        como ponteiro bruto; um endere\c{c}o inv\'alido causaria falha.
  \item \textbf{\'Indice sem verifica\c{c}\~ao}: \code{row} e \code{column}
        s\~ao mantidos dentro dos limites pelas fun\c{c}\~oes, mas n\~ao h\'a
        \emph{assert} de invariantes.
  \item \textbf{Scroll custa 1.920 leituras/escritas}: aceit\'avel para um
        terminal, mas poderia usar \emph{hardware scrolling} ou mem\'oria
        intermedi\'aria.
\end{itemize}

\chapter{Teclado PS/2}
\label{cap:keyboard}

\section{Vis\~ao geral}

O driver \code{src/keyboard.rs} implementa o subconjunto necess\'ario para
ler teclas de um teclado PS/2 conectado ao controlador 8042. Ele combina
conceitos de dois subsistemas do Linux:

\begin{itemize}
  \item \code{drivers/input/serio/i8042.c}: acesso ao controlador e ao
        protocolo de comandos;
  \item \code{drivers/input/keyboard/atkbd.c}: interpreta\c{c}\~ao de
        scancodes e teclas modificadoras.
\end{itemize}

N\~ao h\'a IRQ1, fila de eventos, \code{input\_device}, \emph{keycode} Linux
nem \code{evdev}. O driver faz \emph{polling} do buffer de sa\'ida do 8042 e
devolve bytes ASCII diretamente ao shell.

\section{Hardware: controlador 8042}

O 8042 exp\~oe duas portas:

\begin{table}[H]
\centering
\caption{Portas do controlador 8042.}
\label{tab:ports8042}
\begin{tabular}{@{}lll@{}}
\toprule
\textbf{Porta} & \textbf{Nome} & \textbf{Uso} \\
\midrule
\code{0x60} & Dados & Leitura de scancodes; escrita de dados ao teclado. \\
\code{0x64} & Status/Comando & Leitura de status; escrita de comandos ao 8042. \\
\bottomrule
\end{tabular}
\end{table}

O registrador de status (\code{0x64}) tem os bits usados pelo driver:

\begin{table}[H]
\centering
\caption{Bits de status do 8042.}
\label{tab:status8042}
\begin{tabular}{@{}lll@{}}
\toprule
\textbf{Bit} & \textbf{M\'ascara} & \textbf{Significado} \\
\midrule
0 & \code{0x01} & Buffer de sa\'ida cheio (h\'a dado para ler) \\
1 & \code{0x02} & Buffer de entrada cheio (n\~ao enviar comando) \\
5 & \code{0x20} & Dado veio do dispositivo auxiliar (mouse) \\
\bottomrule
\end{tabular}
\end{table}

\section{Inicializa\c{c}\~ao}

A inicializa\c{c}\~ao \'e mais do que limpar o buffer. A BIOS ou o GRUB podem
deixar o controlador em um estado em que o teclado envia scancodes do
conjunto 2, enquanto o driver espera o conjunto 1, ou com a varredura
desabilitada. O procedimento implementado \'e:

\begin{table}[H]
\centering
\caption{Sequ\^encia de inicializa\c{c}\~ao do 8042.}
\label{tab:kbdinit}
\begin{tabular}{@{}clp{7.3cm}@{}}
\toprule
\textbf{\#} & \textbf{Comando} & \textbf{Efeito} \\
\midrule
1 & leituras de \code{0x60} & Descarta bytes pendentes no buffer de sa\'ida. \\
2 & \code{0xAD} em \code{0x64} & Desabilita a porta do teclado. \\
3 & \code{0x20} em \code{0x64} & Solicita leitura do \emph{configuration byte}. \\
4 & leitura de \code{0x60} & Obt\'em o byte de configura\c{c}\~ao atual
(ou usa \code{0x45} se n\~ao houver resposta). \\
5 & ajuste & Liga tradu\c{c}\~ao (bit 6), desliga IRQ1 (bit 0) e liga o clock
(limpa bit 4). \\
6 & \code{0x60} em \code{0x64} & Prepara a escrita do novo byte de
configura\c{c}\~ao. \\
7 & byte em \code{0x60} & Envia o novo \emph{configuration byte}. \\
8 & \code{0xAE} em \code{0x64} & Reabilita a porta do teclado. \\
9 & \code{0xF4} em \code{0x60} & Envia ao teclado o comando \emph{enable
scanning}. \\
10 & leitura de \code{0x60} & Consome o ACK (\code{0xFA}), se presente. \\
11 & limpeza & Descarta qualquer res\'iduo do buffer. \\
\bottomrule
\end{tabular}
\end{table}

\begin{nota}[Por que traduzir?]
O teclado PS/2 moderno costuma falar o \emph{scan code set 2}, mas o
controlador 8042 pode apresentar ao sistema o \emph{set 1} quando o bit de
tradu\c{c}\~ao est\'a ligado. O Linux tamb\'em configura essa tradu\c{c}\~ao
no i8042. Sem ela, as teclas chegam em c\'odigos diferentes dos esperados
pela tabela do driver.
\end{nota}

As esperas usam at\'e 100.000 itera\c{c}\~oes, tanto para o buffer de entrada
esvaziar quanto para o buffer de sa\'ida encher. Se o controlador n\~ao
responder, a inicializa\c{c}\~ao continua em vez de travar.

\section{Scancodes e tradu\c{c}\~ao para ASCII}

O driver trabalha com o conjunto 1. A tabela \code{KEYMAP} tem 128 posi\c{c}\~oes
e cobre letras, d\'igitos, espa\c{c}o e alguns s\'imbolos. Posi\c{c}\~oes n\~ao
mapeadas valem zero e s\~ao ignoradas.

\begin{table}[H]
\centering
\caption{Exemplos de scancodes do conjunto 1.}
\label{tab:scancodes}
\begin{tabular}{@{}llll@{}}
\toprule
\textbf{Tecla} & \textbf{Make} & \textbf{Break} & \textbf{Resultado} \\
\midrule
\code{A} & \code{0x1E} & \code{0x9E} & \code{a} / \code{A} com Shift \\
\code{1} & \code{0x02} & \code{0x82} & \code{1} / \code{!} com Shift \\
\code{Enter} & \code{0x1C} & \code{0x9C} & \code{\textbackslash n} \\
\code{Backspace} & \code{0x0E} & \code{0x8E} & \code{0x08} \\
\code{Tab} & \code{0x0F} & \code{0x8F} & \code{\textbackslash t} \\
\code{Espa\c{c}o} & \code{0x39} & \code{0xB9} & espa\c{c}o \\
\code{Shift esquerdo} & \code{0x2A} & \code{0xAA} & ativa/desativa Shift \\
\code{Shift direito} & \code{0x36} & \code{0xB6} & ativa/desativa Shift \\
\code{Tecla estendida} & \code{0xE0} & --- & prefixo; c\'odigo seguinte ignorado \\
\bottomrule
\end{tabular}
\end{table}

\subsection{Shift}

O estado de Shift \'e mantido em \code{static SHIFT: AtomicBool}. As teclas
\code{0x2A} e \code{0x36} ligam o estado; os \emph{break codes} \code{0xAA}
e \code{0xB6} desligam. A fun\c{c}\~ao \code{apply\_shift} transforma letras
em mai\'usculas e aplica o s\'imbolo secund\'ario de d\'igitos e pontua\c{c}\~ao:

\begin{lstlisting}[caption={Mapeamento com Shift.},label={lst:shift}]
b'a'..=b'z' => byte - 32,
b'1'..=b'9' => b"!@#$%^&*("[(byte - b'1') as usize],
b'0' => b')',
b'-' => b'_',
b'=' => b'+',
/* ... */
\end{lstlisting}

\subsection{Teclas estendidas}

O prefixo \code{0xE0} liga a flag \code{EXTENDED}. O pr\'oximo scancode \'e
descartado e a flag \'e limpa. Isso evita que setas, Insert, Delete, Home,
End e teclas do teclado num\'erico sejam interpretadas como letras. A
consequ\^encia \'e que essas teclas n\~ao t\^em fun\c{c}\~ao no shell.

\subsection{Mouse}

Dados do mouse compartilham o mesmo 8042. Se o bit 5 do status estiver
ligado, o byte lido \'e descartado. Isso evita que movimentos do mouse
apare\c{c}am como letras.

\section{Fluxo de leitura}

\begin{figure}[H]
\centering
\begin{tikzpicture}[
  kstep/.style={draw=rinuxblue, fill=rinuxlight, rounded corners=2pt,
               minimum width=5.1cm, minimum height=0.8cm,
               align=center, font=\scriptsize},
  dec/.style={draw=orange!70!black, fill=orange!8, diamond, aspect=2.2,
              inner sep=1pt, align=center, font=\scriptsize},
  arrow/.style={-{Latex[length=2mm]}, thick, rinuxblue},
  node distance=0.38cm
]
\node[kstep] (status) {L\^e status (\code{0x64})};
\node[dec, below=of status] (full) {bit 0 setado?};
\node[kstep, below=of full] (read) {L\^e dado (\code{0x60})};
\node[dec, below=of read] (mouse) {bit 5 (mouse)?};
\node[dec, below=of mouse] (ext) {\code{0xE0} ou extended?};
\node[dec, below=of ext] (break) {bit 7 (break)?};
\node[kstep, below=of break] (key) {Traduz e retorna byte};
\node[kstep, right=2.4cm of full] (none) {Retorna \code{None}};
\draw[arrow] (status) -- (full);
\draw[arrow] (full) -- node[right,font=\tiny]{sim} (read);
\draw[arrow] (full.east) -- node[above,font=\tiny]{n\~ao} (none.west);
\draw[arrow] (read) -- (mouse);
\draw[arrow] (mouse) -- node[right,font=\tiny]{n\~ao} (ext);
\draw[arrow] (ext) -- node[right,font=\tiny]{n\~ao} (break);
\draw[arrow] (break) -- node[right,font=\tiny]{n\~ao} (key);
\draw[arrow] (mouse.east) -- ++(0.9,0) |- (none.east);
\draw[arrow] (ext.east) -- ++(0.5,0) |- (none.east);
\draw[arrow] (break.east) -- ++(0.2,0) |- (none.east);
\end{tikzpicture}
\caption{Fluxo de \code{read\_char()}. Caminhos de descarte levam a \code{None}.}
\label{fig:kbdflow}
\end{figure}

O shell chama \code{keyboard::read\_char()} repetidamente. Quando a fun\c{c}\~ao
retorna \code{None}, o shell tenta a serial e, se tamb\'em n\~ao houver dado,
executa \code{spin\_loop()}.

\section{API e fun\c{c}\~oes internas}

\begin{table}[H]
\centering
\caption{Fun\c{c}\~oes de \code{keyboard.rs}.}
\label{tab:kbdapi}
\begin{tabularx}{\linewidth}{@{}lX@{}}
\toprule
\textbf{Fun\c{c}\~ao} & \textbf{Descri\c{c}\~ao} \\
\midrule
\code{init()} & Configura tradu\c{c}\~ao, IRQ, clock e varredura do 8042. \\
\code{read\_char()} & L\^e e traduz um scancode; retorna \code{Option<u8>}. \\
\code{build\_keymap()} & Constr\'oi a tabela de 128 entradas em tempo de
compila\c{c}\~ao. \\
\code{apply\_shift(byte)} & Aplica o efeito de Shift a um caractere. \\
\code{wait\_input\_clear()} & Espera o buffer de entrada do 8042 esvaziar. \\
\code{wait\_output\_full()} & Espera um byte ficar dispon\'ivel na sa\'ida. \\
\code{flush\_output()} & Descarta at\'e 64 bytes pendentes. \\
\code{send\_command(cmd)} & Envia um comando ao 8042 (porta \code{0x64}). \\
\code{send\_data(data)} & Envia um byte ao teclado (porta \code{0x60}). \\
\bottomrule
\end{tabularx}
\end{table}

\section{Limita\c{c}\~oes}

\begin{itemize}
  \item \textbf{Polling}: n\~ao h\'a IRQ1; o shell precisa consultar
        constantemente. A CPU fica ocupada mesmo sem entrada.
  \item \textbf{Sem Caps Lock, Num Lock e Scroll Lock}: os LEDs e os estados
        n\~ao s\~ao tratados.
  \item \textbf{Sem AltGr}: layouts que dependem de AltGr n\~ao produzem
        caracteres adicionais.
  \item \textbf{Sem teclas estendidas funcionais}: setas, Delete, Home, End,
        Page Up/Down e teclado num\'erico s\~ao descartados.
  \item \textbf{Sem layout ABNT2}: a tabela \'e baseada em US/QWERTY; a
        posi\c{c}\~ao f\'isica das teclas e o conjunto de s\'imbolos podem
        diferir.
  \item \textbf{Sem \emph{key release} \'util}: al\'em de Shift, os
        \emph{break codes} s\~ao ignorados; n\~ao h\'a detec\c{c}\~ao de tecla
        pressionada para jogos ou atalhos.
  \item \textbf{Sem tratamento de erros}: bytes de erro (\code{0x00},
        \code{0xFF}, \code{0xEE}) n\~ao t\^em tratamento espec\'ifico.
  \item \textbf{Depend\^encia de hardware}: assume 8042 compat\'ivel. Em
        sistemas apenas USB-HID, o teclado pode n\~ao funcionar sem
        \emph{legacy USB} habilitado na BIOS.
\end{itemize}

\chapter{Serial COM1}
\label{cap:serial}

\section{Vis\~ao geral}

O m\'odulo \code{src/serial.rs} implementa um subconjunto do driver 8250 do
Linux para a UART 16550 em COM1. A serial tem dois pap\'eis:

\begin{itemize}
  \item \textbf{depura\c{c}\~ao}: todo texto impresso no VGA tamb\'em \'e
        enviado para COM1;
  \item \textbf{entrada}: al\'em do teclado PS/2, o shell aceita comandos
        digitados no terminal serial, o que viabiliza testes automatizados.
\end{itemize}

N\~ao h\'a interrup\c{c}\~oes, FIFO baseada em IRQ, \emph{termios}, \emph{tty}
layer ou controle de fluxo. O acesso \'e feito por polling de registradores
de I/O.

\section{Hardware}

A UART 16550 em COM1 usa uma base de 8 portas a partir de \code{0x3F8}:

\begin{table}[H]
\centering
\caption{Registradores da UART 16550 (COM1).}
\label{tab:uartregs}
\begin{tabular}{@{}llll@{}}
\toprule
\textbf{Offset} & \textbf{Nome} & \textbf{Leitura} & \textbf{Escrita} \\
\midrule
\code{+0} & RBR / THR & Receive Buffer & Transmit Holding \\
\code{+1} & IER / DLM & Interrupt Enable & Divisor Latch High \\
\code{+2} & IIR / FCR & Interrupt ID & FIFO Control \\
\code{+3} & LCR & Line Control & Line Control \\
\code{+4} & MCR & Modem Control & Modem Control \\
\code{+5} & LSR & Line Status & --- \\
\code{+6} & MSR & Modem Status & --- \\
\code{+7} & SCR & Scratch & Scratch \\
\bottomrule
\end{tabular}
\end{table}

O registrador LSR tem os bits usados pelo driver:

\begin{table}[H]
\centering
\caption{Bits do Line Status Register (LSR).}
\label{tab:lsr}
\begin{tabular}{@{}lll@{}}
\toprule
\textbf{Bit} & \textbf{M\'ascara} & \textbf{Significado} \\
\midrule
0 & \code{0x01} & Dado dispon\'ivel para leitura \\
5 & \code{0x20} & Transmitter Holding Register vazio (pode enviar) \\
\bottomrule
\end{tabular}
\end{table}

\section{Inicializa\c{c}\~ao}

A configura\c{c}\~ao implementada \'e 38400 baud, 8 bits de dados, sem
paridade e 1 bit de parada (8N1), com FIFO habilitada e IRQs desligadas.

\begin{table}[H]
\centering
\caption{Sequ\^encia de inicializa\c{c}\~ao da COM1.}
\label{tab:uartinit}
\begin{tabular}{@{}llp{7.2cm}@{}}
\toprule
\textbf{Offset} & \textbf{Valor} & \textbf{Efeito} \\
\midrule
\code{+1} & \code{0x00} & Desabilita interrup\c{c}\~oes. \\
\code{+3} & \code{0x80} & Liga DLAB para acessar o divisor. \\
\code{+0} & \code{0x03} & Divisor baixo = 3. \\
\code{+1} & \code{0x00} & Divisor alto = 0. \\
\code{+3} & \code{0x03} & 8 bits, sem paridade, 1 stop; desliga DLAB. \\
\code{+2} & \code{0xC7} & FIFO habilitada, limpa buffers, gatilho de 14 bytes. \\
\code{+4} & \code{0x0B} & DTR, RTS e OUT2; IRQs desligadas. \\
\bottomrule
\end{tabular}
\end{table}

O divisor 3 produz:

\[
\text{baud} = \frac{115200}{3} = 38400.
\]

Ap\'os a configura\c{c}\~ao, o driver l\^e o LSR. Se o valor for \code{0xFF},
considera que n\~ao h\'a UART no barramento e marca \code{PRESENT = false}.
Essa heur\'istica funciona porque um barramento sem dispositivo geralmente
devolve todos os bits em 1.

\begin{nota}[Por que 38400 e n\~ao 115200?]
O divisor 3 \'e um valor conservador e amplamente suportado por emula\c{c}\~oes
e cabos simples. A taxa n\~ao precisa ser alta para um terminal de texto.
Para 115200, o divisor seria 1; para 9600, 12. A escolha n\~ao afeta o
funcionamento do VGA.
\end{nota}

\section{Escrita}

\begin{lstlisting}[caption={Envio de um byte pela UART.},label={lst:uartwrite}]
fn write_byte(byte: u8) {
    unsafe {
        while inb(COM1 + 5) & 0x20 == 0 {}
        outb(COM1, byte);
    }
}
\end{lstlisting}

O la\c{c}o espera o THR esvaziar. Se n\~ao houver UART e o barramento devolver
\code{0xFF}, o bit \code{0x20} j\'a est\'a ligado e o la\c{c}o n\~ao trava;
o byte \'e simplesmente descartado. Se houver uma UART presente mas
travada, o la\c{c}o pode esperar indefinidamente.

\code{write\_str} converte \code{\textbackslash n} em \code{\textbackslash r
\textbackslash n}, porque terminais serial geralmente esperam CRLF:

\begin{lstlisting}[caption={Escrita de string com convers\~ao de nova linha.},label={lst:uartstr}]
pub fn write_str(text: &str) {
    for byte in text.bytes() {
        if byte == b'\n' {
            write_byte(b'\r');
        }
        write_byte(byte);
    }
}
\end{lstlisting}

\section{Leitura}

\begin{lstlisting}[caption={Leitura n\~ao bloqueante de um byte.},label={lst:uartread}]
pub fn read_byte() -> Option<u8> {
    if !PRESENT.load(Ordering::Relaxed) {
        return None;
    }
    unsafe {
        if inb(COM1 + 5) & 0x01 == 0 {
            return None;
        }
        Some(inb(COM1))
    }
}
\end{lstlisting}

A fun\c{c}\~ao \'e n\~ao bloqueante: se n\~ao h\'a byte, retorna \code{None}.
O shell a chama depois de \code{keyboard::read\_char()}, de modo que a
serial funciona como fonte alternativa de entrada.

\section{API}

\begin{table}[H]
\centering
\caption{API de \code{serial.rs}.}
\label{tab:serialapi}
\begin{tabularx}{\linewidth}{@{}lX@{}}
\toprule
\textbf{Fun\c{c}\~ao} & \textbf{Descri\c{c}\~ao} \\
\midrule
\code{init()} & Configura COM1 em 38400 8N1 e detecta presen\c{c}a. \\
\code{read\_byte()} & Retorna um byte recebido ou \code{None}. \\
\code{write\_byte(byte)} & Envia um byte, esperando o THR. Interna. \\
\code{write\_str(text)} & Envia uma string, convertendo \code{\textbackslash n}
em CRLF. \\
\bottomrule
\end{tabularx}
\end{table}

\section{Uso no QEMU e em hardware f\'isico}

No QEMU, a serial pode ser exposta de v\'arias formas:

\begin{lstlisting}[style=shell,caption={Exemplos de exposi\c{c}\~ao da serial no QEMU.},label={lst:qemuserial}]
# Terminal interativo no stdio (usado por `make run`)
qemu-system-x86_64 -cdrom vga-kernel.iso -boot order=d -m 128M -serial stdio

# Socket Unix para testes automatizados
qemu-system-x86_64 -cdrom vga-kernel.iso -boot order=d -m 128M \
  -display none -serial unix:/tmp/vga-serial.sock,server=on,wait=off
\end{lstlisting}

Em um PC f\'isico, COM1 aparece em \code{0x3F8} e pode ser acessada com um
cabo USB-serial e um terminal em 38400 8N1. A BIOS pode desabilitar a porta
ou remape\'a-la; nesse caso o driver n\~ao a encontra.

\section{Limita\c{c}\~oes}

\begin{itemize}
  \item \textbf{Sem IRQs}: a leitura depende de polling.
  \item \textbf{Sem FIFO de entrada}: bytes podem ser perdidos se chegarem
        mais r\'apido do que o shell consulta.
  \item \textbf{Sem controle de fluxo}: RTS/CTS n\~ao s\~ao usados.
  \item \textbf{Sem \emph{echo} autom\'atico}: o pr\'oprio shell ecoa os
        caracteres recebidos.
  \item \textbf{Sem suporte a outras portas}: apenas COM1.
  \item \textbf{Detec\c{c}\~ao heur\'istica}: \code{LSR == 0xFF} pode n\~ao
        ser suficiente em todos os chipsets.
\end{itemize}

\chapter{Shell e linha de comando}
\label{cap:shell}

\section{Vis\~ao geral}

O \code{src/shell.rs} implementa um terminal de linha de comando sem
aloca\c{c}\~ao din\^amica. Ele n\~ao \'e um shell POSIX, n\~ao interpreta
pipes, redirecionamentos, vari\'aveis, aspas ou scripts. \'E um la\c{c}o que:

\begin{enumerate}
  \item imprime o prompt \code{rinux> };
  \item l\^e bytes do teclado PS/2 ou da serial;
  \item ecoa caracteres imprim\'iveis e trata Backspace;
  \item ao receber Enter, executa o comando;
  \item volta ao in\'icio.
\end{enumerate}

\section{La\c{c}o principal}

\begin{lstlisting}[caption={Estrutura do la\c{c}o do shell.},label={lst:shellloop}]
pub fn run() -> ! {
    loop {
        let mut line = [0u8; LINE_MAX];
        let mut len = 0usize;
        let mut console = tty::active();
        crate::print!("rinux> ");

        loop {
            tty::pump_input();
            if tty::active() != console {
                console = tty::active();
                len = 0;
                crate::print!("rinux> ");
                continue;
            }

            let byte = match tty::read_byte() {
                Some(byte) => byte,
                None => { core::hint::spin_loop(); continue; }
            };

            match byte {
                b'\n' | b'\r' => {
                    println!();
                    execute(&line[..len]);
                    break;
                }
                0x08 | 0x7F => {
                    if len > 0 {
                        len -= 1;
                        crate::print!("\x08");
                    }
                }
                b'\t' => {}
                0x20..=0x7E => {
                    if len < LINE_MAX {
                        line[len] = byte;
                        len += 1;
                        crate::print!("{}", byte as char);
                    }
                }
                _ => {}
            }
        }
    }
}
\end{lstlisting}

\subsection{Buffer de linha}

A linha \'e um array fixo de 128 bytes na pilha do shell:

\begin{lstlisting}[caption={Buffer de linha.},label={lst:linebuffer}]
const LINE_MAX: usize = 128;
let mut line = [0u8; LINE_MAX];
let mut len = 0usize;
\end{lstlisting}

Quando a linha est\'a cheia, novos caracteres s\~ao silenciosamente
ignorados. N\~ao h\'a beep, mensagem de erro nem crescimento din\^amico.

\subsection{Entrada e eco}

A entrada \'e bombeada por \code{tty::pump\_input}, que drena o teclado PS/2 e
a serial COM1 para a fila do console virtual ativo. Ele tamb\'em interpreta
\code{Alt+F1}..\code{Alt+F8} como troca de console. O shell ent\~ao consome
bytes do console ativo com \code{tty::read\_byte}.

Se nenhum byte estiver dispon\'ivel, o shell executa \code{spin\_loop()} e
tenta novamente. N\~ao h\'a \emph{sleep}, j\'a que a espera \'e feita por
polling.

Caracteres imprim\'iveis (0x20 a 0x7E) s\~ao armazenados no buffer e ecoados
no console. Backspace (\code{0x08} ou \code{0x7F}) decrementa o tamanho e
imprime \code{\textbackslash x08}, que a camada TTY interpreta como apagar
a c\'elula anterior. Tab (\code{\textbackslash t}) \'e ignorado.

Se o console ativo mudar durante a edi\c{c}\~ao, a linha parcial \'e
descartada e o shell imprime um novo prompt no console selecionado.

\section{Interpreta\c{c}\~ao}

\subsection{Trim}

Antes de comparar, a linha passa por \code{trim}, que remove espa\c{c}os
ASCII no in\'icio e no fim. A fun\c{c}\~ao retorna uma \emph{slice} da
pr\'opria linha, sem c\'opia:

\begin{lstlisting}[caption={Remo\c{c}\~ao de espa\c{c}os nas extremidades.},label={lst:trim}]
fn trim(bytes: &[u8]) -> &[u8] {
    let start = bytes.iter()
        .position(|byte| !byte.is_ascii_whitespace())
        .unwrap_or(bytes.len());
    let end = bytes.iter()
        .rposition(|byte| !byte.is_ascii_whitespace())
        .map_or(start, |index| index + 1);
    &bytes[start..end]
}
\end{lstlisting}

Uma linha vazia ou composta apenas de espa\c{c}os n\~ao executa nada.

\subsection{Convers\~ao para texto}

\code{as\_text} tenta interpretar os bytes como UTF-8. Como o console s\'o
aceita ASCII, na pr\'atica a convers\~ao sempre funciona para texto digitado;
se falhar, retorna a string fixa \code{"<bytes invalidos>"}.

\subsection{Comandos}

A execu\c{c}\~ao \'e uma cadeia de compara\c{c}\~oes \code{if}/\code{else if}.
N\~ao h\'a tabela de despacho, \emph{parsing} de argumentos nem an\'alise
sint\'atica.

\begin{table}[H]
\centering
\caption{Comandos implementados.}
\label{tab:comandos}
\begin{tabularx}{\linewidth}{@{}lX@{}}
\toprule
\textbf{Comando} & \textbf{Efeito} \\
\midrule
\code{help} & Lista os comandos dispon\'iveis. \\
\code{clear} & Limpa o console virtual ativo com \code{ESC[2J ESC[H}. \\
\code{echo <texto>} & Imprime o texto ap\'os o primeiro espa\c{c}o. \\
\code{info} & Mostra informa\c{c}\~oes do kernel, TTY e modelo de entrada. \\
\code{tty} & Mostra o console virtual ativo. \\
\code{tty [0..7]} & Troca para o console virtual indicado. \\
\code{getty} & Roda o login em ring 3 no console ativo. \\
\code{getty [0..7]} & Troca de console e roda o login nele. \\
\code{demand} & Exercita o demand paging do kernel. \\
\code{usertest} & Roda o demo de demand paging em ring 3. \\
\code{userprot} & Tenta acessar o kernel a partir do ring 3. \\
\code{mem} & Mostra estat\'isticas de frames, faults e bits U/S. \\
\code{reboot} & Reinicia o PC via comando \code{0xFE} ao 8042. \\
\code{halt} & Imprime uma mensagem e para a CPU com \code{cli; hlt}. \\
\bottomrule
\end{tabularx}
\end{table}

\begin{nota}[Detalhes de comportamento]
\begin{itemize}
  \item \code{echo} sem argumento n\~ao \'e reconhecido como comando; apenas
        \code{echo } seguido de texto casa com o prefixo.
  \item \code{echo   ola} preserva os espa\c{c}os internos ap\'os o prefixo
        \code{"echo "}; o \emph{trim} remove apenas as extremidades.
  \item A compara\c{c}\~ao \'e sens\'ivel a mai\'usculas: \code{HELP} n\~ao
        funciona.
  \item Comandos desconhecidos produzem duas linhas de mensagem.
\end{itemize}
\end{nota}

\section{Sa\'ida de \code{info}}

O comando \code{info} imprime:

\begin{lstlisting}[style=shell,caption={Sa\'ida do comando \code{info}.},label={lst:info}]
Kernel: Rust, no_std, x86_64
Boot:   Multiboot2 (BIOS) -> modo longo
Video:  VGA texto 80x25, buffer 0xB8000
TTY:    8 consoles virtuais, ativo = tty0 (Alt+F1..Alt+F8)
Entrada: teclado PS/2 e serial COM1 (polling)
IDT:    256 portas; excecoes + IRQ0 (timer ~100 Hz)
Memoria: paginas de 4 KiB, demand paging e frames fisicos
Usuario: ring 3 com GDT/TSS, protecao U/S e int 0x80
Ticks do timer: 366
\end{lstlisting}

A maioria das linhas \'e est\'atica; o n\'umero do console ativo e a contagem
de ticks do timer s\~ao lidos no momento em que o comando roda.

\section{Comandos especiais}

\subsection{\code{halt}}

Imprime \code{"CPU parada. Use reset ou feche o QEMU."} e chama
\code{arch::halt\_loop()}, que executa \code{cli; hlt} indefinidamente. A
\'unica sa\'ida \'e reset f\'isico ou fechar o emulador.

\subsection{\code{reboot}}

Imprime \code{"Reiniciando..."} e chama \code{arch::reboot()}. O reset pode
n\~ao funcionar em todos os ambientes; se o comando ao 8042 for ignorado, a
CPU termina em \code{halt\_loop()}.

\section{Como estender}

Para adicionar um comando:

\begin{enumerate}
  \item acrescente um ramo \code{else if line == b"meucomando"} em
        \code{execute};
  \item use \code{println!} para a sa\'ida;
  \item se precisar de argumentos, use \code{strip\_prefix} como em
        \code{echo};
  \item mantenha o comando sem aloca\c{c}\~ao e sem bloqueio;
  \item atualize a lista de \code{help} e esta documenta\c{c}\~ao.
\end{enumerate}

Exemplo:

\begin{lstlisting}[caption={Esqueleto de um novo comando.},label={lst:novocomando}]
} else if line == b"versao" {
    println!("Rust VGA Kernel 0.1.0");
} else if let Some(arg) = line.strip_prefix(b"upper ") {
    for byte in arg {
        crate::print!("{}", byte.to_ascii_uppercase() as char);
    }
    println!();
}
\end{lstlisting}

\section{Limita\c{c}\~oes}

\begin{itemize}
  \item Sem hist\'orico, setas, Ctrl+C, Ctrl+D ou edi\c{c}\~ao no meio da
        linha.
  \item Sem aspas, escape, pipes, redirecionamento ou vari\'aveis.
  \item Sem argumentos estruturados: cada comando decide como interpretar o
        restante.
  \item Sem execu\c{c}\~ao de programas: n\~ao h\'a \emph{loader} nem
        processos.
  \item Sem tradu\c{c}\~ao de UTF-8: acentos digitados podem virar \code{?}.
  \item A linha \'e limitada a 128 bytes e o excesso \'e descartado em
        sil\^encio.
\end{itemize}

\chapter{Mem\'oria virtual, demand paging e interrup\c{c}\~oes}
\label{cap:memoria-usuario}

Este cap\'itulo documenta os m\'odulos \code{gdt.rs}, \code{idt.rs},
\code{paging.rs} e \code{user.rs}, adicionados depois do console e do shell.
O objetivo \'e sair do modelo ``ring 0, interrup\c{c}\~oes desligadas e
mem\'oria est\'atica'' para um kernel capaz de:

\begin{itemize}
  \item receber exce\c{c}\~oes e IRQs por uma IDT;
  \item receber o timer do PIT em IRQ0;
  \item alocar frames f\'isicos de 4 KiB;
  \item materializar p\'aginas somente quando a CPU gera \code{\#PF};
  \item executar c\'odigo em ring 3 com uma barreira real entre usu\'ario e
        kernel.
\end{itemize}

\section{Ordem de inicializa\c{c}\~ao}

A sequ\^encia em \code{kernel\_main} passou a ser:

\begin{enumerate}
  \item \code{vga::init()} e \code{serial::init()};
  \item \code{paging::init(magic, info)} --- libera os frames dispon\'iveis
        descritos pelo Multiboot2;
  \item \code{gdt::init()} --- GDT definitiva, segmentos ring 3 e TSS;
  \item \code{keyboard::init()} --- teclado PS/2 por \emph{polling};
  \item \code{idt::init()} --- IDT, remapeamento do PIC e PIT;
  \item \code{sti} e entrada no shell.
\end{enumerate}

A IDT s\'o \'e instalada depois que a GDT tem a TSS, pois a porta de
\emph{double fault} usa uma pilha IST. O timer s\'o \'e desmascarado depois
que a IDT est\'a pronta.

\section{GDT, TSS e pilhas de privil\'egio}

A GDT do stub de boot cont\'em apenas descritores ring 0. Para entrar em
ring 3, \code{gdt.rs} monta uma tabela com sete entradas:

\begin{table}[H]
\centering
\caption{Descritores da GDT definitiva.}
\label{tab:gdt-usuario}
\begin{tabular}{@{}rll@{}}
\toprule
\textbf{Selector} & \textbf{Conte\'udo} & \textbf{DPL} \\
\midrule
\code{0x08} & c\'odigo 64 bits & 0 \\
\code{0x10} & dados & 0 \\
\code{0x18} & dados & 3 \\
\code{0x20} & c\'odigo 64 bits & 3 \\
\code{0x28} & TSS (duas entradas) & 0 \\
\bottomrule
\end{tabular}
\end{table}

A TSS \'e obrigat\'oria para a troca autom\'atica de pilha. Quando a CPU
passa de ring 3 para ring 0 por interrup\c{c}\~ao ou exce\c{c}\~ao, ela l\^e
\code{rsp0} da TSS e continua a empilhar nessa pilha. O kernel reserva:

\begin{itemize}
  \item \code{USER\_KERNEL\_STACK}: 16 KiB para interrup\c{c}\~oes vindas do
        usu\'ario;
  \item \code{DOUBLE\_FAULT\_STACK}: 8 KiB como \code{IST1}, para que um
        \code{\#DF} n\~ao tente usar uma pilha possivelmente corrompida.
\end{itemize}

O campo \code{iomap\_base} aponta para o fim da TSS, desativando o bitmap de
I/O. Assim, instru\c{c}\~oes \code{in}/\code{out} em ring 3 geram \code{\#GP}.

\section{IDT, stubs e normaliza\c{c}\~ao do erro}

A IDT tem 256 portas de 16 bytes. As 32 primeiras s\~ao exce\c{c}\~oes; as
16 seguintes s\~ao as IRQs do PIC remapeadas para \code{0x20..0x2F}; a porta
\code{0x80} \'e uma \emph{syscall} did\'atica com DPL = 3.

O problema cl\'assico \'e que a CPU empilha um c\'odigo de erro apenas em
algumas exce\c{c}\~oes (8, 10--14, 17, 21, 29 e 30). Se cada \emph{stub}
tratasse a pilha de forma diferente, o layout visto pelo Rust mudaria. Para
uniformizar, \code{idt.rs} gera stubs em assembly:

\begin{lstlisting}[caption={Stubs que normalizam o c\'odigo de erro.},label={lst:isrstubs}]
.macro isr_no_err n
.global isr_stub_\n
isr_stub_\n:
    pushq $0          # erro sintetico
    pushq $\n         # numero do vetor
    jmp isr_common
.endm

.macro isr_err n
.global isr_stub_\n
isr_stub_\n:
    pushq $\n         # erro real ja empilhado pela CPU
    jmp isr_common
.endm
\end{lstlisting}

O \emph{stub} comum salva os registradores de uso geral, preserva
\code{xmm0..xmm15} e \code{MXCSR}, alinha a pilha para chamar
\code{interrupt\_dispatch} e, no retorno, restaura tudo e executa
\code{iretq}. O frame visto pelo Rust \'e:

\begin{lstlisting}[caption={Layout de \code{InterruptFrame} ap\'os o stub comum.},label={lst:frame}]
#[repr(C)]
pub struct InterruptFrame {
    pub r15: u64, /* ... r8 */
    pub rdi: u64, /* ... rax */
    pub vector: u64,
    pub error: u64,
    pub rip: u64,
    pub cs: u64,
    pub rflags: u64,
    pub rsp: u64,   // so existe se veio de ring 3
    pub ss: u64,    // so existe se veio de ring 3
}
\end{lstlisting}

O vetor 14 chama o tratador de \emph{page fault}; o vetor 13 trata
\code{\#GP}; o vetor 0x80 encerra o programa de usu\'ario; o vetor 32
incrementa o contador de \emph{ticks} e envia EOI ao PIC.

\section{PIC e PIT}

O PIC 8259 \'e remapeado para n\~ao colidir com as exce\c{c}\~oes:

\begin{lstlisting}[caption={Remapeamento do PIC e inicializa\c{c}\~ao do PIT.},label={lst:picpit}]
outb(0x20, 0x11); outb(0xA0, 0x11);  // ICW1
outb(0x21, 0x20); outb(0xA1, 0x28);  // ICW2: vetores 0x20/0x28
outb(0x21, 0x04); outb(0xA1, 0x02);  // ICW3
outb(0x21, 0x01); outb(0xA1, 0x01);  // ICW4

let divisor = 11_931;                // 1,193182 MHz / 100 Hz
outb(0x43, 0x36);
outb(0x40, divisor as u8);
outb(0x40, (divisor >> 8) as u8);

outb(0x21, 0xFE);                    // so IRQ0
outb(0xA1, 0xFF);
\end{lstlisting}

A IRQ1 continua mascarada porque \code{keyboard.rs} l\^e o 8042 por
\emph{polling}. Habilit\'a-la exigiria um buffer de scancodes e coordenar o
acesso ao controlador.

\section{Alocador de frames f\'isicos}

\code{paging::init} percorre as tags Multiboot2. A tag 6 (mapa de mem\'oria)
fornece entradas de 24 bytes com base, tamanho e tipo. O alocador come\c{c}a
com todos os 262.144 frames do primeiro GiB marcados como usados e libera
somente as faixas do tipo 1 (RAM dispon\'ivel).

O bitmap tem 32 KiB:

\begin{lstlisting}[caption={Bitmap de frames e aloca\c{c}\~ao.},label={lst:bitmap}]
static mut FRAME_BITMAP: [u8; 32 * 1024] = [0xFF; 32 * 1024];

pub fn alloc_frame() -> Option<u64> {
    /* procura um bit zero, marca como usado e zera os 4096 bytes */
}
\end{lstlisting}

O primeiro MiB e o fim do kernel (s\'imbolo \code{\_\_kernel\_end} do
\emph{linker script}) ficam sempre reservados. A estrutura Multiboot2 tamb\'em
\'e reservada depois que as faixas dispon\'iveis s\~ao liberadas.

Como o stub mapeia o primeiro GiB por identidade, escrever no endere\c{c}o
f\'isico do frame \'e equivalente a escrever no endere\c{c}o virtual baixo.
Isso permite zerar frames sem precisar de um mapa direto adicional.

\section{Demand paging e prote\c{c}\~ao U/S}

Cada entrada de tabela tem bits de presen\c{c}a (P), escrita (W) e
usu\'ario (U/S). O x86\_64 exige U/S = 1 em \emph{todos} os n\'iveis
(PML4, PDPT, PD e PT) para permitir um acesso em ring 3. O kernel usa isso
para criar duas regi\~oes:

\begin{table}[H]
\centering
\caption{Regi\~oes de demand paging.}
\label{tab:demand}
\begin{tabularx}{\linewidth}{@{}lrrX@{}}
\toprule
\textbf{Regi\~ao} & \textbf{Base} & \textbf{Tamanho} & \textbf{U/S das folhas} \\
\midrule
Usu\'ario & \code{0x0000008000000000} & 2 MiB & 1 \\
Kernel   & \code{0xFFFFFFFFC0000000} & 2 MiB & 0 \\
\bottomrule
\end{tabularx}
\end{table}

Nenhuma folha dessas regi\~oes existe no boot. Quando o \code{\#PF} ocorre,
o tratador l\^e \code{CR2} e o c\'odigo de erro e decide:

\begin{enumerate}
  \item se o endere\c{c}o pertence \`a regi\~ao de usu\'ario e a p\'agina
        n\~ao est\'a presente, aloca um frame e mapeia com
        \code{P|W|U/S}; as tabelas intermedi\'arias tamb\'em recebem U/S = 1;
  \item se pertence \`a regi\~ao do kernel e o bit U/S do erro \'e zero,
        mapeia com \code{P|W} e U/S = 0;
  \item se um usu\'ario tentar a regi\~ao do kernel, o erro tem U/S = 1 e o
        tratador recusa o mapeamento; o programa de usu\'ario \'e encerrado;
  \item se o endere\c{c}o est\'a fora das duas regi\~oes, a falha \'e fatal
        para o kernel ou encerra o usu\'ario.
\end{enumerate}

\begin{nota}[Por que o kernel n\~ao pode ``abrir'' uma tabela do supervisor]
Mapear uma p\'agina de usu\'ario n\~ao consiste apenas em ligar U/S na folha.
Se qualquer n\'ivel intermedi\'ario for supervisor, o hardware nega o acesso.
Al\'em disso, transformar uma tabela do kernel em tabela de usu\'ario poderia
expor outros mapeamentos. Por isso \code{next\_table} recusa uma tabela
supervisor no caminho de uma p\'agina de usu\'ario.
\end{nota}

\section{Modo usu\'ario e retorno ao kernel}

\code{user.rs} monta um pequeno programa de m\'aquina em um buffer e o copia
para a p\'agina \code{USER\_CODE}, que \'e residente. A pilha e a p\'agina de
dados ficam sob demanda:

\begin{enumerate}
  \item o programa faz \code{push}/\code{pop}, gerando \code{\#PF} na pilha;
  \item escreve em \code{USER\_DATA}, gerando outro \code{\#PF};
  \item executa \code{int \$0x80} com \code{rax = 0}.
\end{enumerate}

O comando \code{userprot} usa outro programa: ele tenta ler
\code{0xFFFFFFFF80000000} (in\'icio do kernel). A p\'agina existe, mas tem
U/S = 0; o erro de \emph{page fault} tem P = 1 e U/S = 1. O tratador
reconhece a viola\c{c}\~ao, imprime os dados e n\~ao devolve o controle ao
programa.

A entrada em ring 3 \'e feita por \code{iretq}, que carrega CS/SS de DPL 3,
RIP, RSP e RFLAGS. A volta n\~ao usa o caminho normal do \code{iretq}:

\begin{itemize}
  \item \code{user\_enter} salva os registradores callee-saved e o RSP do
        kernel em \code{user\_saved\_rsp};
  \item uma exce\c{c}\~ao ou \code{int \$0x80} cai no \emph{stub} comum em
        ring 0;
  \item \code{user::abort} grava o status e chama \code{user\_resume}, que
        restaura o RSP salvo, recarrega os segmentos do kernel e retorna para
        a chamada de \code{user\_enter}.
\end{itemize}

Esse mecanismo \'e um \emph{longjmp} controlado. Ele evita ter de reconstruir
o frame de interrup\c{c}\~ao para voltar a ring 0 e mant\'em o shell simples.

\section{Comandos e evid\^encias}

\begin{table}[H]
\centering
\caption{Comandos que exercitam a nova infraestrutura.}
\label{tab:comandos-memoria}
\begin{tabularx}{\linewidth}{@{}lX@{}}
\toprule
\textbf{Comando} & \textbf{Efeito} \\
\midrule
\code{demand} & escreve na regi\~ao de demand paging do kernel; o primeiro
acesso gera \code{\#PF} \\
\code{usertest} & entra em ring 3, gera \code{\#PF} na pilha e nos dados e
sai por \code{int \$0x80} \\
\code{userprot} & entra em ring 3 e tenta ler o kernel; a prote\c{c}\~ao U/S
encerra o programa \\
\code{mem} & mostra frames alocados/livres, \emph{page faults}, \emph{ticks}
e os bits U/S de uma p\'agina de usu\'ario e do kernel \\
\bottomrule
\end{tabularx}
\end{table}

Uma execu\c{c}\~ao t\'ipica em QEMU:

\begin{lstlisting}[style=shell,caption={Sa\'ida dos comandos de mem\'oria e ring 3.},label={lst:testdemand}]
rinux> demand
[PF] demanda atendida: 0xffffffffc0000000 (kernel, erro=0x2)
[DEMAND] 0xffffffffc0000000 = 0xdeadbeefcafe1234 (page faults: 1)
rinux> usertest
[USER] entrando em ring 3 (demand paging)...
[PF] demanda atendida: 0x0000008000007ff8 (usuario, erro=0x6)
[PF] demanda atendida: 0x0000008000002000 (usuario, erro=0x6)
[USER] voltou ao kernel; page faults nesta execucao: 2
rinux> userprot
[USER] entrando em ring 3 (teste de protecao)...
[USER] page fault em 0xffffffff80000000: endereco fora das regioes de demand paging (erro=0x5, rip=0x000000800000100a)
[USER] programa encerrado com status 1
rinux> mem
Frames fisicos (4 KiB):
  alocados : 7
  livres   : 32217
  page faults atendidos: 3
  ticks do timer       : 372
  USER_CODE  flags=0x203067 (U/S=1)
  KERNEL_VMA flags=0xe3 (U/S=0)
\end{lstlisting}

O erro \code{0x6} (P = 0, W = 1, U/S = 1) confirma uma escrita de usu\'ario em
p\'agina ausente; o erro \code{0x5} (P = 1, U/S = 1, sem W) confirma uma
leitura de usu\'ario em p\'agina presente e supervisor.

\section{Limita\c{c}\~oes desta implementa\c{c}\~ao}

\begin{itemize}
  \item n\~ao h\'a espa\c{c}os de endere\c{c}o separados por processo: o
        programa de usu\'ario usa a mesma PML4 do kernel;
  \item n\~ao h\'a escalonador nem troca de contexto; o timer apenas conta
        \emph{ticks};
  \item a \'unica \emph{syscall} \'e \code{exit} (\code{rax = 0}); n\~ao
        h\'a valida\c{c}\~ao de ponteiros de usu\'ario;
  \item SMEP/SMAP e NX ainda n\~ao est\~ao habilitados; a separa\c{c}\~ao
        depende apenas de U/S;
  \item o teclado continua em \emph{polling} e IRQ1 permanece mascarada;
  \item as regi\~oes de demand paging t\^em tamanho fixo de 2 MiB e ficam
        fora do primeiro GiB mapeado por identidade.
\end{itemize}

\chapter{Camada TTY, syscalls e getty}
\label{cap:tty}

\section{Do renderizador VGA aos consoles virtuais}

At\'e aqui o driver VGA era dono tanto do conte\'udo da tela quanto da posi\c{c}\~ao
do cursor. Para suportar v\'arios terminais como no Linux, o c\'odigo foi
dividido em duas camadas:

\begin{itemize}
  \item \code{vga.rs} \'e um \emph{renderizador}. Ele conhece a matriz de
        c\'elulas 80$\times$25 em \code{0xB8000}, os registradores CRTC do
        cursor e nada sobre linhas de comando ou entrada;
  \item \code{tty.rs} mant\'em oito consoles virtuais, \code{tty0} a
        \code{tty7}. Cada console tem seu pr\'oprio buffer de c\'elulas,
        cursor e fila de entrada.
\end{itemize}

Somente o console ativo \'e copiado para o framebuffer de hardware. Trocar de
console \'e, portanto, uma c\'opia de buffer mais a atualiza\c{c}\~ao do cursor,
a mesma id\'eia usada pela camada de terminal virtual (\code{vt}) do Linux.

\begin{lstlisting}[caption={Estado de cada console em \code{tty.rs}.},label={lst:ttytty}]
#[derive(Clone, Copy)]
struct Tty {
    cells: [u16; vga::CELLS],
    row: usize,
    column: usize,
    input: [u8; INPUT_SIZE],
    head: usize,
    tail: usize,
}

static TTYS: SpinLock<[Tty; TTY_COUNT]> =
    SpinLock::new([Tty::new(); TTY_COUNT]);
static ACTIVE: AtomicUsize = AtomicUsize::new(0);
\end{lstlisting}

A fila de entrada \'e um pequeno buffer circular de 256 bytes. O
\emph{produtor} \'e \code{pump\_input}, que drena o teclado PS/2 e a COM1; o
\emph{consumidor} \'e o shell ou uma chamada de sistema lendo do console ativo.

\section{Troca de console}

O driver de teclado agora exp\~oe dois tipos de resultado:

\begin{lstlisting}[caption={Tipo de resultado do teclado.},label={lst:keyenum}]
pub enum Key {
    Char(u8),
    Console(usize),
}
\end{lstlisting}

O driver acompanha Shift e Alt. Com Alt pressionado, os scancodes
\code{0x3B}..\code{0x42} (F1..F8) viram \code{Key::Console(0..7)}. A camada
TTY traduz isso em \code{switch(index)} e redesenha o buffer escolhido. A
mesma opera\c{c}\~ao est\'a dispon\'ivel no shell:

\begin{lstlisting}[style=shell,caption={Trocando de console pelo shell.},label={lst:ttyshell}]
rinux> tty
console ativo: tty0 (disponiveis: tty0..tty7)
rinux> tty 3
rinux> tty
console ativo: tty3 (disponiveis: tty0..tty7)
\end{lstlisting}

Como o shell \'e um \'unico la\c{c}o s\'incrono, trocar de console durante a
edi\c{c}\~ao de uma linha descarta a linha parcial e imprime um novo prompt no
console novo. Essa \'e uma simplifica\c{c}\~ao deliberada: ainda n\~ao existe
um \code{task\_struct} por terminal.

\section{Sequ\^encias de escape ANSI/VT100}

A sa\'ida n\~ao se limita a caracteres imprim\'iveis. Cada \code{Tty} mant\'em
sua pr\'opria m\'aquina de estados, para que um programa use as mesmas
sequ\^encias de escape de um terminal Linux sem saber qual console est\'a ativo:

\begin{lstlisting}[caption={Estado de escape guardado em cada console.},label={lst:ttyescape}]
enum EscapeState {
    Ground, // texto normal
    Start,  // recebeu ESC
    Csi,    // recebeu ESC [ e le parametros
}

struct Tty {
    // ...
    escape: EscapeState,
    csi_params: [u16; 4],
    csi_count: usize,
    csi_value: u16,
    csi_digits: bool,
}
\end{lstlisting}

\code{write\_byte} encaminha \code{0x1B} para a m\'aquina de estados. No estado
CSI o interpretador aceita par\^ametros decimais separados por \code{;} e um
byte final; marcadores privados como \code{?} e bytes intermedi\'arios s\~ao
consumidos sem serem impressos:

\begin{lstlisting}[caption={Despacho de ESC e CSI em \code{tty.rs}.},label={lst:ttywritebyte}]
fn write_byte(&mut self, byte: u8) {
    match self.escape {
        EscapeState::Ground => {
            if byte == 0x1B {
                self.escape = EscapeState::Start;
            } else {
                self.write_plain(byte);
            }
        }
        EscapeState::Start => {
            if byte == b'[' {
                self.escape = EscapeState::Csi;
                self.csi_params = [0; 4];
                self.csi_count = 0;
                // ...
            } else {
                self.escape = EscapeState::Ground;
            }
        }
        EscapeState::Csi => self.write_csi(byte),
    }
}
\end{lstlisting}

\begin{table}[H]
\centering
\caption{Sequ\^encias CSI suportadas.}
\label{tab:ansi}
\begin{tabularx}{\linewidth}{@{}lX@{}}
\toprule
\textbf{Sequ\^encia} & \textbf{Efeito} \\
\midrule
\code{ESC[2J}, \code{ESC[3J} & Apaga a tela inteira. \\
\code{ESC[J} & Apaga do cursor at\'e o fim da tela. \\
\code{ESC[K} & Apaga do cursor at\'e o fim da linha. \\
\code{ESC[H}, \code{ESC[<linha>;<coluna>H} & Leva o cursor ao in\'icio ou a
uma posi\c{c}\~ao base 1, limitada a 80$\times$25. \\
\code{ESC[A}, \code{ESC[B}, \code{ESC[C}, \code{ESC[D} & Move o cursor para
cima, baixo, direita ou esquerda (1 por padr\~ao, ou a quantidade dada). \\
\code{ESC[...m} & SGR \'e interpretado e ignorado; cores ainda n\~ao existem. \\
\bottomrule
\end{tabularx}
\end{table}

O comando \code{clear} agora \'e sa\'ida comum, n\~ao um caso especial do
renderizador:

\begin{lstlisting}[caption={Implementa\c{c}\~ao do comando \code{clear}.},label={lst:ttyclear}]
pub fn clear() {
    write_bytes(b"\x1b[2J\x1b[H");
}
\end{lstlisting}

\code{ESC[2J} apaga as c\'elulas mas, seguindo o comportamento VT100, n\~ao
move o cursor; a segunda sequ\^encia \code{ESC[H} o leva ao canto superior
esquerdo. Como \code{clear} passa por \code{write\_bytes}, a sequ\^encia
tamb\'em \'e espelhada na COM1, de modo que um terminal serial tamb\'em limpa.
Cores SGR, tela alternativa, scrollback e visibilidade do cursor ainda est\~ao
fora do escopo.

\section{ABI das chamadas de sistema}

A porta \code{0x80} j\'a existia para os demos em ring 3. Agora ela despacha
para \code{syscall.rs}, com a conven\c{c}\~ao usual de registradores:

\begin{table}[H]
\centering
\caption{Chamadas de sistema implementadas.}
\label{tab:syscalls}
\begin{tabularx}{\linewidth}{@{}clX@{}}
\toprule
\textbf{N\'umero} & \textbf{Nome} & \textbf{Argumentos e resultado} \\
\midrule
0 & \code{exit} & \code{rdi} = status; n\~ao retorna. \\
1 & \code{write} & \code{rdi} = buffer, \code{rsi} = tamanho;
devolve o n\'umero de bytes. \\
2 & \code{read} & \code{rdi} = buffer, \code{rsi} = tamanho,
\code{rdx} = flags (bit 0 desliga o eco); devolve os bytes lidos. \\
3 & \code{get\_tty} & sem argumentos; devolve o \'indice do console ativo. \\
\bottomrule
\end{tabularx}
\end{table}

\subsection{Valida\c{c}\~ao de ponteiros do usu\'ario}

Uma chamada de sistema roda em ring 0 com as mesmas tabelas de p\'agina do
usu\'ario. Tocar em um endere\c{c}o de usu\'ario n\~ao mapeado n\~ao geraria
um page fault de usu\'ario; seria um acesso do kernel. Por isso \code{write}
e \code{read} chamam \code{paging::ensure\_user\_range} antes de
desreferenciar o buffer:

\begin{lstlisting}[caption={Valida\c{c}\~ao e mapeamento sob demanda de um buffer do usu\'ario.},label={lst:ensureuser}]
pub fn ensure_user_range(addr: u64, len: u64) -> Result<(), &'static str> {
    let end = addr.checked_add(len).ok_or("intervalo invalido")?;
    if addr < USER_BASE || end > USER_BASE + USER_REGION_SIZE {
        return Err("ponteiro fora da regiao do usuario");
    }
    // Mapeia com U/S = 1 cada pagina ausente de [addr, end).
    // ...
    Ok(())
}
\end{lstlisting}

A fun\c{c}\~ao rejeita intervalos fora da regi\~ao do usu\'ario e materializa
cada p\'agina ausente com \code{PAGE\_USER | PAGE\_WRITE}. \'E o equivalente
did\'atico dos cuidadosos auxiliares \code{copy\_from\_user}/\code{copy\_to\_user}
do Linux.

\subsection{Leitura bloqueante}

A chamada \code{read} \'e s\'incrona. Enquanto n\~ao h\'a byte dispon\'ivel,
ela chama \code{tty::pump\_input} e \code{spin\_loop}. A porta de interrup\c{c}\~ao
limpa \code{IF}, portanto o timer n\~ao avança durante a espera na
implementa\c{c}\~ao atual. Bytes imprim\'iveis s\~ao gravados no buffer do
usu\'ario e ecoados quando a flag de sem-eco est\'a limpa; Backspace remove o
\'ultimo byte; Enter termina a linha e devolve o caractere de nova linha.

\section{getty e login em ring 3}

O primeiro programa de usu\'ario de verdade \'e um \emph{getty}/login. Ele \'e
escrito em assembly dentro de \code{user.rs}, em uma se\c{c}\~ao pr\'opria, e
copiado como um blob independente de posi\c{c}\~ao para
\code{paging::USER\_CODE}:

\begin{lstlisting}[caption={Fluxo de login do programa em ring 3.},label={lst:gettyasm}]
getty_start:
    GETTY_WRITE getty_login_msg, getty_login_end   # write(2)
    subq $128, %rsp                                # buffer na pilha
    movq %rsp, %rdi                                # buf
    movq $32, %rsi                                 # len
    xorq %rdx, %rdx                                # echo
    movq $2, %rax                                  # read
    int $0x80
    # compara "root" e pede "Password: "
    # segundo read com flags = 1 (sem eco)
    # compara "rinux" e escreve a mensagem final
    # exit(0) ou exit(1)
\end{lstlisting}

Usar uma se\c{c}\~ao dedicada com c\'odigo e strings mant\'em todas as
refer\^encias RIP-relativas v\'alidas depois da c\'opia. O buffer de leitura
fica na pilha do usu\'ario, cuja p\'agina \'e mapeada sob demanda pela primeira
chamada de sistema que a toca.

As credenciais da demonstra\c{c}\~ao s\~ao:

\begin{lstlisting}[style=shell,caption={Login bem-sucedido.},label={lst:gettyout}]
rinux> getty
[GETTY] iniciando login em tty0...
rinux login: root
Password:
Bem-vindo ao rinux, root!
[GETTY] login concluido, voltando ao shell
\end{lstlisting}

\code{getty n} primeiro troca para o console \code{n} e depois executa o
programa nele. Em um teste automatizado pela serial, as credenciais devem ser
enviadas somente depois que o prompt de login daquele console aparecer, pois a
porta serial \'e um \'unico fluxo global.

\section{Rela\c{c}\~ao com o Linux}

\begin{table}[H]
\centering
\caption{Correspond\^encia com o subsistema TTY do Linux.}
\label{tab:ttflinux}
\begin{tabularx}{\linewidth}{@{}lX@{}}
\toprule
\textbf{Linux} & \textbf{Aqui} \\
\midrule
\code{drivers/tty/tty\_io.c} & \code{src/tty.rs}: filas de leitura/escrita. \\
\code{drivers/tty/vt/vt.c} & \code{src/tty.rs}: consoles virtuais e troca. \\
\code{drivers/video/console/vgacon.c} & \code{src/vga.rs}: renderizador de
c\'elulas. \\
\code{drivers/tty/n\_tty.c} & N\~ao implementado; as linhas s\~ao tratadas
diretamente. \\
\code{arch/x86/entry/syscalls/syscall\_64.tbl} & \code{src/syscall.rs}. \\
\code{init/getty.c}, \code{login} & Blob ring 3 em \code{src/user.rs}. \\
\bottomrule
\end{tabularx}
\end{table}

\section{Limita\c{c}\~oes atuais}

\begin{itemize}
  \item Existe um shell, n\~ao um shell por terminal. Trocar de console move o
        mesmo la\c{c}o para outro buffer.
  \item O \code{getty} n\~ao \'e residente: ele roda uma vez, de forma
        s\'incrona, e volta ao shell. Um sistema real mant\'em um
        \code{getty} esperando em cada TTY.
  \item A serial \'e um \'unico fluxo; o teclado \'e o que realmente distingue
        os consoles no uso interativo.
  \item N\~ao h\'a disciplina de linha: nem modo can\^onico, sinais, controle
        de jobs, \code{Ctrl+C}, \code{Ctrl+D} ou ioctls de terminal.
  \item \code{read} espera com interrup\c{c}\~oes desligadas; o timer congela
        durante um prompt de login.
\end{itemize}

\chapter{Build, imagem e execu\c{c}\~ao}
\label{cap:build}

\section{Toolchain e alvo}

O projeto compila para \code{x86\_64-unknown-none}, alvo oficial de Rust para
bare metal x86\_64. Esse alvo:

\begin{itemize}
  \item n\~ao possui \code{std}, \code{alloc} nem \emph{runtime};
  \item usa o modelo de c\'odigo \emph{kernel}, que assume endere\c{c}os no
        ter\c{c}o superior do espa\c{c}o de endere\c{c}amento;
  \item produz c\'odigo posicional apenas quando necess\'ario, com
        \code{relocation-model=static}.
\end{itemize}

O arquivo \code{rust-toolchain.toml} fixa o canal \code{stable} e solicita o
alvo:

\begin{lstlisting}[caption={\code{rust-toolchain.toml}.},label={lst:toolchain}]
[toolchain]
channel = "stable"
targets = ["x86_64-unknown-none"]
\end{lstlisting}

\section{Configura\c{c}\~ao do Cargo}

\begin{lstlisting}[caption={\code{.cargo/config.toml}.},label={lst:cargoconfig}]
[build]
target = "x86_64-unknown-none"
rustc = "scripts/rustc"

[target.x86_64-unknown-none]
rustflags = [
    "-C", "relocation-model=static",
    "-C", "linker=ld.lld",
    "-C", "link-arg=-Tlinker.ld",
    "-C", "link-arg=--gc-sections",
]
\end{lstlisting}

A op\c{c}\~ao \code{rustc} aponta para o \emph{wrapper} \code{scripts/rustc}
do projeto. Ele existe porque, em distribui\c{c}\~oes como Arch Linux, o
\code{cargo} e o \code{rustc} encontrados primeiro no \code{PATH} s\~ao os do
sistema, que n\~ao possuem o alvo \code{x86\_64-unknown-none}. O wrapper
executa:

\begin{lstlisting}[style=shell,caption={Conte\'udo de \code{scripts/rustc}.},label={lst:rustcwrapper}]
#!/bin/sh
if command -v rustup >/dev/null 2>&1; then
    exec rustup run stable rustc "$@"
fi
exec rustc "$@"
\end{lstlisting}

Assim, um \code{cargo build} comum usa o compilador do \code{rustup} mesmo
quando o \code{cargo} \'e o do sistema. Se o \code{rustup} n\~ao estiver
instalado, o wrapper recai no \code{rustc} do \code{PATH}. O caminho
\code{scripts/rustc} \'e relativo ao diret\'orio do projeto, de modo que
n\~ao depende do nome do usu\'ario nem de \code{\$HOME}.

\begin{table}[H]
\centering
\caption{Flags de compila\c{c}\~ao e liga\c{c}\~ao.}
\label{tab:flags}
\begin{tabularx}{\linewidth}{@{}lX@{}}
\toprule
\textbf{Flag} & \textbf{Efeito} \\
\midrule
\code{relocation-model=static} & Evita \emph{PIC}/GOT; o kernel \'e ligado em
endere\c{c}o fixo. \\
\code{linker=ld.lld} & Usa o linker LLVM, compat\'ivel com o script. \\
\code{-Tlinker.ld} & Aplica o layout de mem\'oria do projeto. \\
\code{--gc-sections} & Remove se\c{c}\~oes n\~ao referenciadas. \\
\bottomrule
\end{tabularx}
\end{table}

\section{Perfis de compila\c{c}\~ao}

\begin{lstlisting}[caption={\code{Cargo.toml}.},label={lst:cargo}]
[package]
name = "vga-kernel"
version = "0.1.0"
edition = "2021"
description = "Kernel minimalista em Rust: boot BIOS (multiboot2), console VGA 80x25 e teclado PS/2"
license = "MIT OR Apache-2.0"

[dependencies]

[profile.dev]
panic = "abort"

[profile.release]
panic = "abort"
opt-level = 2
lto = true
codegen-units = 1
\end{lstlisting}

O perfil \code{release} \'e usado na ISO. \code{panic=abort} \'e obrigat\'orio
para um kernel sem \emph{unwinding}; \code{lto} e \code{codegen-units=1}
permitem otimiza\c{c}\~oes entre m\'odulos e reduzem o tamanho.

\section{Makefile}

\begin{lstlisting}[caption={Alvos principais do \code{Makefile}.},label={lst:makefile}]
TARGET  := x86_64-unknown-none
KERNEL  := target/$(TARGET)/release/vga-kernel
ISO     := vga-kernel.iso
ISODIR  := build/iso

RUSTC   ?= $(shell command -v rustup >/dev/null 2>&1 && rustup which rustc)
ifneq ($(RUSTC),)
export RUSTC
endif

CARGO   ?= $(shell command -v rustup >/dev/null 2>&1 && echo "rustup run stable cargo" || echo cargo)

all: iso

kernel:
	$(CARGO) build --release --target $(TARGET)

iso: kernel
	rm -rf $(ISODIR)
	mkdir -p $(ISODIR)/boot/grub
	cp $(KERNEL) $(ISODIR)/boot/kernel.elf
	cp boot/grub/grub.cfg $(ISODIR)/boot/grub/grub.cfg
	grub-mkrescue -o $(ISO) $(ISODIR)

run: iso
	qemu-system-x86_64 -cdrom $(ISO) -boot order=d -m 128M -serial stdio

clean:
	$(CARGO) clean
	rm -rf build $(ISO)
\end{lstlisting}

\begin{nota}[Compatibilidade com o sistema]
Em algumas distribui\c{c}\~oes, \code{/usr/bin/cargo} e \code{/usr/bin/rustc}
n\~ao s\~ao os do \code{rustup} e n\~ao t\^em o alvo
\code{x86\_64-unknown-none}. H\'a duas prote\c{c}\~oes independentes:

\begin{itemize}
  \item o \code{.cargo/config.toml} aponta \code{rustc} para
        \code{scripts/rustc}, que usa o \code{rustup};
  \item o \code{Makefile} detecta \code{rustup}, chama
        \code{rustup run stable cargo} e exporta \code{RUSTC}.
\end{itemize}

Com isso, tanto \code{cargo build} quanto \code{make} funcionam mesmo que o
\code{cargo} do sistema seja o primeiro do \code{PATH}.
\end{nota}

\subsection{Compila\c{c}\~ao direta com Cargo}

Antes do ajuste, \code{cargo build} usava o alvo do \emph{host}
(\code{x86\_64-unknown-linux-gnu}); o linker tentava produzir um execut\'avel
Linux e falhava com \code{duplicate symbol: \_start}, porque o stub do kernel
e o \code{Scrt1.o} da libc definem \code{\_start}. Al\'em disso, o
\code{rustc} do sistema n\~ao tinha o alvo \code{x86\_64-unknown-none}.

Agora o alvo padr\~ao vem do \code{.cargo/config.toml} e o compilador vem do
wrapper. Os comandos usuais funcionam sem vari\'aveis de ambiente:

\begin{lstlisting}[style=shell,caption={Compila\c{c}\~ao direta com Cargo.},label={lst:cargodirect}]
cargo build              # ELF de depuracao em target/x86_64-unknown-none/debug/
cargo build --release    # ELF otimizado usado na ISO
cargo clean              # remove target/
\end{lstlisting}

N\~ao \'e necess\'ario passar \code{--target}: ele j\'a est\'a fixado no
arquivo de configura\c{c}\~ao. O bin\'ario resultante \'e um ELF64
\code{EXEC} para \code{x86-64}, com ponto de entrada
\code{0x100018}, e n\~ao um execut\'avel Linux.

\subsection{Alvos}

\begin{table}[H]
\centering
\caption{Alvos do Makefile.}
\label{tab:maketargets}
\begin{tabularx}{\linewidth}{@{}lX@{}}
\toprule
\textbf{Alvo} & \textbf{A\c{c}\~ao} \\
\midrule
\code{make} / \code{make all} & Compila o kernel e gera a ISO. \\
\code{make kernel} & Compila apenas o ELF release. \\
\code{make iso} & Compila e empacota a ISO com GRUB. \\
\code{make run} & Gera a ISO e inicia o QEMU com VGA e serial no stdio. \\
\code{make clean} & Remove \code{target/}, \code{build/} e a ISO. \\
\bottomrule
\end{tabularx}
\end{table}

\section{Configura\c{c}\~ao do GRUB}

\begin{lstlisting}[caption={\code{boot/grub/grub.cfg}.},label={lst:grubcfg}]
set timeout=0
set default=0

terminal_output console
set gfxpayload=text

insmod multiboot2

menuentry "Rust VGA Kernel (x86_64)" {
    multiboot2 /boot/kernel.elf
    boot
}
\end{lstlisting}

\begin{itemize}
  \item \code{terminal\_output console} mant\'em o GRUB no modo texto VGA,
        garantindo que o kernel encontre \code{0xB8000} como esperado.
  \item \code{gfxpayload=text} impede que o GRUB selecione um modo gr\'afico
        linear.
  \item \code{insmod multiboot2} carrega o m\'odulo do protocolo.
  \item \code{multiboot2} instrui o GRUB a carregar o ELF e passar a
        assinatura e a estrutura de informa\c{c}\~oes.
\end{itemize}

\section{ISO}

O alvo \code{iso} monta uma \'arvore tempor\'aria:

\begin{lstlisting}[style=shell,caption={Conte\'udo da \'arvore da ISO.},label={lst:isotree}]
build/iso/
|-- boot/
|   |-- grub/
|   |   `-- grub.cfg
|   `-- kernel.elf
`-- (arquivos auxiliares gerados pelo grub-mkrescue)
\end{lstlisting}

\code{grub-mkrescue} cria uma imagem El Torito com o GRUB para BIOS
(\code{i386-pc}) e tamb\'em pode incluir suporte a EFI, dependendo dos
m\'odulos instalados. O projeto usa apenas o caminho BIOS.

Observa\c{c}\~oes da compila\c{c}\~ao de refer\^encia:

\begin{table}[H]
\centering
\caption{Artefatos gerados na compila\c{c}\~ao de refer\^encia.}
\label{tab:artefatos}
\begin{tabular}{@{}ll@{}}
\toprule
\textbf{Artefato} & \textbf{Valor observado} \\
\midrule
ELF & \code{target/x86\_64-unknown-none/release/vga-kernel} \\
Tamanho do ELF & aproximadamente 48 KiB \\
ISO & \code{vga-kernel.iso} \\
Tamanho da ISO & aproximadamente 33,8 MB \\
Setores da ISO & 16.517 setores de 2.048 bytes \\
Sistema de arquivos & ISO 9660 com Rock Ridge \\
Boot & El Torito, \code{boot\_hybrid.img} \\
\bottomrule
\end{tabular}
\end{table}

\section{Execu\c{c}\~ao}

\subsection{\code{make run}}

\begin{lstlisting}[style=shell,caption={Execu\c{c}\~ao padr\~ao.},label={lst:makerun}]
make run
# equivale a:
qemu-system-x86_64 -cdrom vga-kernel.iso -boot order=d -m 128M -serial stdio
\end{lstlisting}

A janela do QEMU mostra a tela VGA. O terminal de onde o comando foi
executado recebe a mesma sa\'ida pela serial e pode enviar comandos.

\subsection{Execu\c{c}\~ao sem interface gr\'afica}

Para testes automatizados, use \code{-display none} e sockets:

\begin{lstlisting}[style=shell,caption={QEMU headless com sockets de serial e monitor.},label={lst:qemuheadless}]
qemu-system-x86_64 -cdrom vga-kernel.iso -boot order=d -m 128M \
  -display none \
  -serial unix:/tmp/vga-serial.sock,server=on,wait=off \
  -monitor unix:/tmp/vga-mon.sock,server=on,wait=off
\end{lstlisting}

O monitor permite inspecionar a mem\'oria VGA:

\begin{lstlisting}[style=shell,caption={Dump da mem\'oria de texto VGA pelo monitor.},label={lst:vgadump}]
xp /4000bx 0xb8000
\end{lstlisting}

Cada c\'elula tem dois bytes: caractere e atributo. Um script pode reconstruir
as 25 linhas e verificar o conte\'udo.

\subsection{Valida\c{c}\~ao do ELF e da ISO}

\begin{lstlisting}[style=shell,caption={Verifica\c{c}\~oes \'uteis.},label={lst:validacoes}]
grub-file --is-x86-multiboot2 target/x86_64-unknown-none/release/vga-kernel
file vga-kernel.iso
xorriso -indev vga-kernel.iso -report_el_torito plain
readelf -h target/x86_64-unknown-none/release/vga-kernel
readelf -S target/x86_64-unknown-none/release/vga-kernel
\end{lstlisting}

\section{Solu\c{c}\~ao de problemas}

\begin{table}[H]
\centering
\caption{Problemas comuns e solu\c{c}\~oes.}
\label{tab:troubleshooting}
\begin{tabularx}{\linewidth}{@{}p{4.5cm}X@{}}
\toprule
\textbf{Sintoma} & \textbf{Causa prov\'avel / solu\c{c}\~ao} \\
\midrule
\code{can't find crate for `core'} & Cargo/rustc do sistema sem o alvo.
Use \code{rustup target add x86\_64-unknown-none} e
\code{rustup run stable cargo}. \\
\code{linker `ld.lld' not found} & Instale \code{lld} ou ajuste
\code{-C linker}. \\
\code{grub-mkrescue: command not found} & Instale \code{grub} e
\code{libisoburn}/\code{xorriso}. \\
QEMU tenta o disco e falha & Use \code{-boot order=d}. \\
Tela preta ap\'os o GRUB & Verifique se \code{gfxpayload=text} est\'a no
\code{grub.cfg} e se o ELF \'e Multiboot2 v\'alido. \\
Teclado n\~ao responde & Verifique \emph{legacy USB}/PS/2 na BIOS; teste a
serial. \\
Comandos aparecem na serial mas n\~ao no VGA & O modo texto VGA pode n\~ao
estar dispon\'ivel (UEFI sem CSM). \\
\code{make run} n\~ao abre janela & Falta servidor gr\'afico; use
\code{-display none} ou \code{-nographic}. \\
\bottomrule
\end{tabularx}
\end{table}

\section{Reprodutibilidade}

O build n\~ao depende de crates externos, mas depende de vers\~oes de Rust,
LLVM, GRUB e xorriso. Para maior reprodutibilidade, recomenda-se registrar:

\begin{itemize}
  \item a vers\~ao do Rust (\code{rustc --version});
  \item a vers\~ao do QEMU (\code{qemu-system-x86\_64 --version});
  \item a vers\~ao do GRUB (\code{grub-mkrescue --version});
  \item o hash do commit do projeto.
\end{itemize}

O \code{Cargo.lock} existe, mas n\~ao cont\'em depend\^encias externas; ele
apenas registra o pr\'oprio pacote.

\chapter{Testes e valida\c{c}\~ao}
\label{cap:testes}

\section{Estrat\'egia}

O projeto n\~ao possui testes unit\'arios automatizados no sentido usual
(\code{cargo test}), porque o c\'odigo depende de hardware, de um bootloader
e de um ambiente de execu\c{c}\~ao bare metal. A valida\c{c}\~ao foi feita por
\emph{testes de integra\c{c}\~ao} em QEMU, com verifica\c{c}\~ao de:

\begin{itemize}
  \item sa\'ida textual na serial;
  \item conte\'udo da mem\'oria VGA, lido pelo monitor do QEMU;
  \item entrada pelo teclado PS/2, simulada com \code{sendkey};
  \item entrada pela serial, enviada por socket Unix;
  \item validade do ELF e da ISO com ferramentas externas.
\end{itemize}

\section{Matriz de testes}

\begin{table}[H]
\centering
\caption{Matriz de testes executados.}
\label{tab:testmatrix}
\begin{tabularx}{\linewidth}{@{}lXc@{}}
\toprule
\textbf{ID} & \textbf{Teste} & \textbf{Resultado} \\
\midrule
T01 & \code{make clean \&\& make} conclui sem erros & Passou \\
T02 & \code{grub-file --is-x86-multiboot2} aceita o ELF & Passou \\
T03 & ISO reconhecida como El Torito BIOS & Passou \\
T04 & QEMU/SeaBIOS inicializa a ISO e chega ao banner & Passou \\
T05 & Prompt \code{rinux>} aparece na tela VGA & Passou \\
T06 & Entrada PS/2: \code{help} lista os comandos & Passou \\
T07 & Entrada PS/2: \code{echo ola} imprime \code{ola} & Passou \\
T08 & Entrada serial: \code{echo serial ok} funciona & Passou \\
T09 & Prompt retorna ap\'os cada comando & Passou \\
T10 & \code{info} mostra as configura\c{c}\~oes do kernel & Passou \\
T11 & \code{clear} envia \code{ESC[2J ESC[H} e limpa a tela & Passou \\
T12 & \code{reboot} e \code{halt} encerram o fluxo & Verificado manualmente \\
T13 & IDT instalada e timer IRQ0 incrementa \code{ticks} & Passou \\
T14 & \code{demand} gera \code{\#PF} e mapeia a p\'agina do kernel & Passou \\
T15 & \code{usertest} roda em ring 3, gera dois page faults e retorna & Passou \\
T16 & \code{userprot} viola U/S e \'e encerrado pelo tratador & Passou \\
T17 & \code{mem} mostra U/S = 1 para usu\'ario e 0 para o kernel & Passou \\
T18 & \code{tty} reporta \code{tty0} e \code{tty 3} troca de console & Passou \\
T19 & \code{getty} autentica \code{root}/\code{rinux} em ring 3 & Passou \\
T20 & Senha n\~ao \'e ecoada; senha errada encerra com status 1 & Passou \\
\bottomrule
\end{tabularx}
\end{table}

\section{Evid\^encias de execu\c{c}\~ao}

\subsection{Banner e prompt}

\begin{lstlisting}[style=shell,caption={Tela VGA ap\'os o boot.},label={lst:testboot}]
==============================================================
  Rust VGA Kernel - x86_64 / modo longo
  Boot: Multiboot2 (BIOS)
  Video: VGA texto 80x25 @ 0xB8000
  Memoria: paginas de 4 KiB + demand paging
  Interrupcoes: IDT + PIC + timer IRQ0 (~100 Hz)
==============================================================
Multiboot2: magic=0x36d76289 info=0x00112040

Digite 'help' para ver os comandos.

rinux>
\end{lstlisting}

O valor de \code{info} varia a cada execu\c{c}\~ao, porque depende da
aloca\c{c}\~ao feita pelo GRUB. O \code{magic} deve ser sempre
\code{0x36d76289}.

\subsection{Comando \code{help} via PS/2}

\begin{lstlisting}[style=shell,caption={Sa\'ida de \code{help} ap\'os entrada PS/2.},label={lst:testhelp}]
rinux> help
Comandos disponiveis:
  help          mostra esta ajuda
  clear         limpa o console ativo
  echo <texto>  imprime o texto
  info          informacoes do kernel
  tty [0..7]    mostra ou troca o console virtual
  getty [0..7]  roda o login em ring 3 no console
  demand        testa demand paging no kernel
  usertest      roda um programa em ring 3 com page fault
  userprot      testa a protecao usuario/kernel
  mem           estatisticas de frames e paginas
  reboot        reinicia o PC
  halt          para a CPU
rinux>
\end{lstlisting}

\subsection{Comando \code{echo} via PS/2}

\begin{lstlisting}[style=shell,caption={Sa\'ida de \code{echo ola} ap\'os entrada PS/2.},label={lst:testecho}]
rinux> echo ola
ola
rinux>
\end{lstlisting}

\subsection{\code{clear} via ANSI}

\begin{lstlisting}[style=shell,caption={Sequ\^encia de escape observada na COM1.},label={lst:testclear}]
rinux> clear
^[[2J^[[H
rinux>
\end{lstlisting}

A nota\c{c}\~ao com circunflexo \'e como um monitor serial exibe o byte
\code{0x1B}. No console VGA o comando apaga as 2.000 c\'elulas e devolve o
cursor para \code{(0,0)}, sem imprimir os caracteres de escape como texto.

\subsection{Entrada serial}

\begin{lstlisting}[style=shell,caption={Entrada e sa\'ida pela COM1.},label={lst:testserial}]
echo serial ok
serial ok
rinux>
\end{lstlisting}

O texto digitado no socket serial aparece no VGA porque o shell ecoa cada
caractere e o console escreve tanto no VGA quanto na serial.

\section{Script de teste headless}

O procedimento abaixo pode ser adaptado para CI. Ele inicia o QEMU sem
interface gr\'afica, envia um comando pela serial e l\^e a sa\'ida.

\begin{lstlisting}[language=Python,caption={Esqueleto de teste headless com socket serial.},label={lst:testscript}]
import socket, subprocess, time, os

ISO = "vga-kernel.iso"
SER = "/tmp/vga-serial.sock"
MON = "/tmp/vga-mon.sock"

for path in (SER, MON):
    try: os.unlink(path)
    except FileNotFoundError: pass

qemu = subprocess.Popen([
    "qemu-system-x86_64", "-cdrom", ISO, "-boot", "order=d",
    "-m", "128M", "-display", "none",
    "-serial", f"unix:{SER},server=on,wait=off",
    "-monitor", f"unix:{MON},server=on,wait=off",
])

time.sleep(7)
ser = socket.socket(socket.AF_UNIX, socket.SOCK_STREAM)
ser.connect(SER)
ser.sendall(b"echo serial ok\r")
time.sleep(1)
data = ser.recv(65536).decode(errors="replace")
assert "serial ok" in data
ser.close()
qemu.terminate()
qemu.wait(timeout=10)
\end{lstlisting}

Para verificar o VGA, conecte ao monitor, envie
\code{xp /4000bx 0xb8000}, interprete os pares (caractere, atributo) e
reconstrua as linhas. O teste PS/2 usa \code{sendkey}:

\begin{lstlisting}[style=shell,caption={Envio de teclas pelo monitor do QEMU.},label={lst:sendkey}]
sendkey h
sendkey e
sendkey l
sendkey p
sendkey ret
\end{lstlisting}

\begin{atencao}[N\~ao envie \code{quit} antes de terminar]
O comando \code{quit} no monitor encerra o QEMU. Se o script ainda precisar
inspecionar o VGA ou enviar mais teclas, use uma conex\~ao separada para cada
comando e s\'o encerre o processo ao final.
\end{atencao}

\section{Testes de build e formato}

\begin{lstlisting}[style=shell,caption={Comandos de valida\c{c}\~ao estrutural.},label={lst:testbuild}]
make clean && make
grub-file --is-x86-multiboot2 target/x86_64-unknown-none/release/vga-kernel
readelf -h target/x86_64-unknown-none/release/vga-kernel
readelf -l target/x86_64-unknown-none/release/vga-kernel
xorriso -indev vga-kernel.iso -report_el_torito plain
\end{lstlisting}

Resultados esperados:

\begin{itemize}
  \item \code{grub-file} retorna c\'odigo 0;
  \item \code{readelf -h} mostra \code{Type: EXEC}, \code{Machine: Advanced
        Micro Devices X86-64} e \code{Entry point address: 0x100018};
  \item \code{readelf -l} mostra um segmento baixo em \code{0x100000} e
        segmentos altos em \code{0xFFFFFFFF8010...};
  \item \code{xorriso} reporta \emph{El Torito} com
        \code{boot\_hybrid.img}.
\end{itemize}

\section{Limita\c{c}\~oes dos testes}

\begin{itemize}
  \item N\~ao h\'a cobertura automatizada de todos os caminhos do teclado
        (Shift, Tab, Backspace, teclas estendidas).
  \item O teste PS/2 depende do \code{sendkey} do QEMU e do layout de teclas
        do monitor; nem todos os caracteres s\~ao f\'aceis de simular.
  \item A leitura do VGA pelo monitor depende do formato de dump e do
        endere\c{c}o \code{0xB8000}.
  \item O comportamento em hardware f\'isico pode diferir do QEMU,
        especialmente em UEFI sem CSM, teclados USB e chipsets sem 8042.
  \item \code{reboot} e \code{halt} foram verificados manualmente; n\~ao h\'a
        teste automatizado de reset.
  \item N\~ao h\'a testes de estresse, tempo ou concorr\^encia.
\end{itemize}

\section{Checklist de aceita\c{c}\~ao}

\begin{itemize}
  \item[$\square$] \code{make} gera ELF e ISO sem erros.
  \item[$\square$] \code{grub-file} reconhece Multiboot2.
  \item[$\square$] QEMU chega ao banner e ao prompt \code{rinux>}.
  \item[$\square$] \code{help}, \code{echo}, \code{info} e \code{clear}
        funcionam pelo teclado PS/2.
  \item[$\square$] Entrada pela serial funciona.
  \item[$\square$] \code{info} mostra o timer e \code{mem} mostra
        \emph{ticks} crescentes.
  \item[$\square$] \code{demand} mapeia a p\'agina do kernel ap\'os
        \code{\#PF}.
  \item[$\square$] \code{usertest} gera dois page faults e retorna ao shell.
  \item[$\square$] \code{userprot} \'e encerrado pela prote\c{c}\~ao U/S.
  \item[$\square$] \code{tty 3} troca para \code{tty3} e \code{tty 0}
        retorna ao \code{tty0}.
  \item[$\square$] \code{getty} aceita \code{root}/\code{rinux} e volta ao
        shell; senha errada encerra com status 1.
  \item[$\square$] \code{halt} para a CPU e \code{reboot} reinicia (quando o
        ambiente suporta).
  \item[$\square$] A documenta\c{c}\~ao LaTeX compila e gera o PDF.
\end{itemize}

\chapter{Similaridade com o kernel Linux}
\label{cap:similaridade}

\section{O que significa ``similaridade''}

Comparar este projeto com o Linux exige separar tr\^es perguntas diferentes:

\begin{enumerate}
  \item \textbf{Similaridade funcional}: quais servi\c{c}os do Linux existem
        aqui e com que profundidade?
  \item \textbf{Similaridade arquitetural}: as solu\c{c}\~oes de boot,
        mem\'oria e drivers usam os mesmos conceitos e mecanismos de
        hardware?
  \item \textbf{Similaridade de c\'odigo}: houve c\'opia, tradu\c{c}\~ao
        linha a linha ou reutiliza\c{c}\~ao de algoritmos?
\end{enumerate}

As respostas s\~ao bastante diferentes. O projeto \'e conceitualmente
pr\'oximo do Linux no caminho de boot e no uso de VGA/PS/2, mas cobre uma
fra\c{c}\~ao min\'uscula do kernel completo. N\~ao h\'a c\'odigo copiado: a
implementa\c{c}\~ao \'e independente e escrita em Rust.

\section{Metodologia}

A compara\c{c}\~ao usa como refer\^encia a \'arvore Linux 7.3.1 presente em
\code{linux-7.3.1/}. Foram coletadas:

\begin{itemize}
  \item linhas f\'isicas dos 12 arquivos Linux mais pr\'oximos;
  \item linhas f\'isicas dos 19 m\'odulos Rust;
  \item funcionalidades presentes, ausentes e parciais;
  \item mecanismos de hardware usados (Multiboot2, GDT, PAE, long mode,
        VGA texto, 8042, UART 16550);
  \item presen\c{c}a ou aus\^encia de subsistemas maiores (IDT, alocador,
        scheduler, syscalls, FS, rede).
\end{itemize}

A pontua\c{c}\~ao \'e uma estimativa fundamentada, n\~ao uma m\'etrica formal.
Ela serve para comunicar ordem de grandeza e evitar tanto a conclus\~ao
ing\^enua de que ``\'e um Linux em Rust'' quanto a conclus\~ao oposta de que
``n\~ao tem nada a ver''.

\section{Correspond\^encia de arquivos}

\begin{longtable}{@{}p{7.2cm}r p{4.3cm}@{}}
\caption{Arquivos Linux an\'alogos e m\'odulos Rust correspondentes.}
\label{tab:correspondencia}\\
\toprule
\textbf{Arquivo Linux 7.3.1} & \textbf{Linhas} & \textbf{Equivalente Rust} \\
\midrule
\endfirsthead
\toprule
\textbf{Arquivo Linux 7.3.1} & \textbf{Linhas} & \textbf{Equivalente Rust} \\
\midrule
\endhead
\code{arch/x86/boot/header.S} & 619 & \code{boot.rs} (cabe\c{c}alho) \\
\code{arch/x86/boot/pm.c} & 125 & \code{boot.rs} (stub) \\
\code{arch/x86/kernel/head\_64.S} & 686 & \code{boot.rs} (modo longo) \\
\code{arch/x86/kernel/cpu/common.c} & 2.688 & \code{boot.rs} (GDT) + \code{arch.rs} \\
\code{arch/x86/mm/init\_64.c} & 1.557 & \code{boot.rs} (p\'aginas) \\
\code{drivers/video/console/vgacon.c} & 1.226 & \code{vga.rs} + \code{tty.rs} \\
\code{drivers/input/serio/i8042.c} & 1.664 & \code{keyboard.rs} \\
\code{drivers/input/keyboard/atkbd.c} & 1.981 & \code{keyboard.rs} \\
\code{drivers/tty/serial/8250/8250\_port.c} & 3.638 & \code{serial.rs} \\
\code{kernel/printk/printk.c} & 5.188 & macros + \code{tty.rs} \\
\code{kernel/panic.c} & 1.274 & \emph{panic handler} em \code{main.rs} \\
\code{init/main.c} & 1.706 & \code{kernel\_main} em \code{main.rs} \\
\code{mm/slab\_common.c}, \code{mm/slub.c} & --- & \code{heap.rs} \\
\code{fs/ramfs/inode.c} & --- & \code{ramfs.rs} \\
\code{fs/open.c}, \code{fs/file\_table.c} & --- & \code{vfs.rs} \\
\code{fs/proc/} & --- & \code{proc.rs} \\
\code{kernel/sched/core.c} & --- & \code{task.rs} \\
\code{arch/x86/entry/syscalls/syscall\_64.tbl} & --- & \code{syscall.rs} \\
\midrule
\textbf{Total Linux (12 primeiros arquivos)} & \textbf{22.352} & --- \\
\textbf{Total Rust (19 m\'odulos)} & \textbf{7.973} & --- \\
\bottomrule
\end{longtable}

As linhas marcadas com \code{---} ampliam o mapeamento conceitual depois da
adi\c{c}\~ao do heap, do VFS, do \code{/proc}, das tarefas e das syscalls de
arquivo; elas n\~ao fazem parte da compara\c{c}\~ao de linhas, que continua
usando os doze primeiros arquivos Linux.

A raz\~ao de linhas \'e:

\[
\frac{7.973}{22.352} \approx 35{,}67\%.
\]

Em rela\c{c}\~ao \`a \'arvore Linux completa, que tem aproximadamente
38.062.672 linhas de \code{.c}, \code{.h} e \code{.S} na vers\~ao usada, a
raz\~ao \'e:

\[
\frac{7.973}{38.062.672} \approx 0{,}0209\%.
\]

\begin{atencao}[Linhas de c\'odigo n\~ao medem similaridade]
Os 12 arquivos Linux n\~ao t\^em o mesmo escopo que os 19 m\'odulos Rust. O
\code{vgacon.c}, por exemplo, suporta v\'arias arquiteturas, consoles
virtuais e modos gr\'aficos; o \code{vga.rs} implementa apenas texto 80x25.
Al\'em disso, grande parte das linhas Linux \'e coment\'ario, configura\c{c}\~ao
e tratamento de casos de hardware. A compara\c{c}\~ao de LOC serve apenas como
ordem de grandeza.
\end{atencao}

\section{Similaridade funcional por dimens\~ao}

A tabela abaixo distribui 100 pontos entre dez dimens\~oes e atribui uma nota
de 0 a 100 para o quanto o projeto cobre da dimens\~ao correspondente no
Linux. A \'ultima coluna \'e a contribui\c{c}\~ao ponderada.

\begin{table}[H]
\centering
\caption{Pontua\c{c}\~ao de similaridade funcional e arquitetural.}
\label{tab:score}
\begin{tabular}{@{}lcrrr@{}}
\toprule
\textbf{Dimens\~ao} & \textbf{Peso} & \textbf{Linux} & \textbf{Projeto} & \textbf{Contrib.} \\
\midrule
Boot e protocolo de firmware & 10 & 100 & 30 & 3,00 \\
Transi\c{c}\~ao de modo e GDT & 8 & 100 & 35 & 2,80 \\
Pagina\c{c}\~ao e MMU & 12 & 100 & 25 & 3,00 \\
Console e v\'ideo & 15 & 100 & 25 & 3,75 \\
Teclado e entrada & 12 & 100 & 18 & 2,16 \\
Serial e tty & 5 & 100 & 20 & 1,00 \\
Interrup\c{c}\~oes e exce\c{c}\~oes & 12 & 100 & 30 & 3,60 \\
Mem\'oria f\'isica e heap & 8 & 100 & 25 & 2,00 \\
Processos, scheduler e syscalls & 12 & 100 & 45 & 5,40 \\
Arquivos, rede e outros drivers & 6 & 100 & 10 & 0,60 \\
\midrule
\textbf{Total} & \textbf{100} & --- & --- & \textbf{27,31} \\
\bottomrule
\end{tabular}
\end{table}

O \'indice ponderado \'e \textbf{27,3\%}. Ele deve ser lido como:

\begin{quote}
``Dentro das dez dimens\~oes escolhidas, o projeto reproduz cerca de 27,3\%
do que o Linux oferece, com peso maior para boot, console e teclado.''
\end{quote}

N\~ao significa que 27,3\% do c\'odigo do Linux foi reimplementado. Em
n\'umero de subsistemas e servi\c{c}os, a fra\c{c}\~ao \'e muito menor: o
projeto tem apenas um heap \emph{first-fit}, um escalonador preemptivo justo
e ponderado com \code{nice}, pesos, \code{vruntime}, \code{sleep} bloqueante,
filas de espera por evento, sinais no estilo Unix e colheita de zumbis, um
sistema de arquivos em mem\'oria e um \code{/proc} sint\'etico, enquanto o
Linux tamb\'em tem buddy allocator, slab, CFS/EEVDF completo com classes de
tempo real, uma infraestrutura de sinais e wait queues muito mais rica, um VFS
completo com block devices e rede.

\subsection{Leitura por dimens\~ao}

\begin{itemize}
  \item \textbf{Boot (30/100)}: o projeto usa Multiboot2 e GRUB, enquanto o
        Linux usa o protocolo de boot pr\'oprio, setup em tempo real,
        descompress\~ao e stubs EFI. O conceito de ``carregar um ELF e
        transferir o controle'' \'e o mesmo; a implementa\c{c}\~ao \'e bem
        mais simples.
  \item \textbf{Modo longo e GDT (35/100)}: a sequ\^encia PAE $\rightarrow$
        LME $\rightarrow$ PG \'e a mesma do \code{head\_64.S}; h\'a GDT com
        ring 3 e TSS, mas faltam \emph{per-CPU}, \emph{stack canary},
        \emph{paravirt} e tratamento de CPUs secund\'arias.
  \item \textbf{Pagina\c{c}\~ao (25/100)}: h\'a \emph{huge pages} de 2 MiB no
        boot, alocador de frames por bitmap, tabelas de 4 KiB e demand
        paging; faltam buddy allocator, \code{vmalloc}, \code{mmap}, COW,
        swap, NUMA e prote\c{c}\~ao \code{NX}.
  \item \textbf{Console (25/100)}: o modo texto VGA e a rolagem s\~ao
        an\'alogos ao \code{vgacon} e h\'a oito consoles virtuais com troca
        por \code{Alt+F1}..\code{Alt+F8}; faltam fontes, Unicode,
        \emph{fbcon}, \code{printk} com n\'iveis e buffer de log.
  \item \textbf{Teclado (18/100)}: o acesso ao 8042 e a tradu\c{c}\~ao de
        scancodes s\~ao semelhantes, mas faltam IRQ1, \emph{input core},
        \code{evdev}, LEDs, layouts e teclas estendidas.
  \item \textbf{Serial (20/100)}: a programa\c{c}\~ao b\'asica da 16550 \'e
        parecida, mas faltam \emph{tty layer}, \emph{termios}, IRQs, FIFO
        de entrada, DMA e controle de fluxo.
  \item \textbf{Interrup\c{c}\~oes (30/100)}: h\'a IDT, \emph{stubs} para
        exce\c{c}\~oes, PIC remapeado e PIT em IRQ0; faltam APIC, IRQ1,
        IRQs esp\'urias, \emph{softirq}, workqueues e preemp\c{c}\~ao.
  \item \textbf{Mem\'oria/heap (25/100)}: h\'a alocador de frames por bitmap,
        demand paging e um heap \emph{first-fit} de 1 MiB atr\'as de
        \code{\#[global\_allocator]}, mas n\~ao h\'a buddy allocator, slab,
        page cache, \code{vmalloc}, swap ou NUMA.
  \item \textbf{Processos (45/100)}: h\'a tabela de tarefas com PID/PPID e
        estados, troca de contexto que preserva \code{RFLAGS}, escalonador
        preemptivo justo e ponderado com \code{nice}, pesos no estilo
        \code{sched\_prio\_to\_weight} e \code{vruntime}, \code{sleep}
        bloqueante com sem\^antica de \code{schedule\_timeout}, filas de
        espera por evento com \code{wait\_event}/\code{wake\_up}, sinais no
        estilo Unix com m\'ascaras de pendentes/bloqueados, disposi\c{c}\~oes e
        a\c{c}\~oes padr\~ao, tarefas paradas/continuadas, status de sa\'ida,
        zumbis, \code{wait}/\code{SIGCHLD} e doze syscalls; n\~ao h\'a
        \code{fork}/\code{exec}, espa\c{c}os de endere\c{c}o separados,
        classes de tempo real nem \'arvore rubro-negra completa do CFS/EEVDF.
  \item \textbf{FS/rede (10/100)}: h\'a ramfs, VFS com dezesseis descritores e
        um \code{/proc} sint\'etico, mas n\~ao h\'a block layer, driver de
        disco, tabela de mounts, modelo de permiss\~oes, page cache, TCP/IP ou
        netfilter.
\end{itemize}

\section{Similaridade arquitetural}

A tabela a seguir compara mecanismos concretos, independentemente do tamanho
do c\'odigo.

\begin{table}[H]
\centering
\caption{Compara\c{c}\~ao de mecanismos de hardware e arquitetura.}
\label{tab:mecanismos}
\begin{tabularx}{\linewidth}{@{}p{3.6cm}p{4.7cm}X@{}}
\toprule
\textbf{Mecanismo} & \textbf{Linux 7.3.1} & \textbf{Rust VGA Kernel} \\
\midrule
Carregamento & Protocolo Linux + setup + descompress\~ao & Multiboot2 via GRUB \\
Entrada inicial & \code{startup\_32}/\code{startup\_64} & \code{\_start} em assembly \\
Modo longo & PAE + LME + \code{CR0.PG} & Mesma sequ\^encia, simplificada \\
GDT & \code{gdt\_table}, \emph{per-CPU}, TSS & 7 descritores, ring 3 e
TSS/IST \\
Pagina\c{c}\~ao & 4/5 n\'iveis, alocadores, \code{vmalloc} & 4 n\'iveis,
2 MiB \emph{huge pages} + folhas 4 KiB sob demanda \\
Console & \code{vgacon}, VT, \code{fbcon}, \code{printk} & \code{vga.rs}
renderizador + \code{tty.rs} com oito consoles \\
Teclado & i8042 + serio + atkbd + input & 8042 por polling (IRQ1 mascarada) \\
Serial & 8250 + tty + termios & 16550 por polling \\
Entrada & IRQ1/IRQ4, \emph{softirq}, workqueues & IRQ0 (timer) + polling no
shell + \code{int 0x80} \\
Tempo & PIT/APIC, \code{jiffies}, timers & PIT a ~100 Hz, contador de
\emph{ticks}, contabilidade ponderada de \code{vruntime} \\
Mem\'oria & buddy, slab, \code{vmalloc}, page cache & Bitmap de frames,
demand paging, heap \emph{first-fit} \\
Processos & CFS/EEVDF, namespaces, cgroups & Tabela de tarefas,
escalonador justo e ponderado (\code{nice}/peso/\code{vruntime}), \code{sleep}
bloqueante, wait queues, sinais e \code{wait}/zumbis; sem \code{fork} ou
namespaces \\
Syscalls & \code{syscall}/\code{sysret}, \code{ptrace} & \code{int 0x80} com
doze chamadas (\code{exit}..\code{kill}) \\
Arquivos & VFS, inodes, dentries, page cache & ramfs, 16 descritores,
\code{/proc} sint\'etico; sem block layer \\
Rede & sockets, TCP/IP, netfilter & Ausente \\
\bottomrule
\end{tabularx}
\end{table}

\section{Similaridade de c\'odigo}

\begin{itemize}
  \item \textbf{C\'opia de c\'odigo}: nenhuma. O projeto foi escrito em Rust
        e n\~ao cont\'em trechos traduzidos do C ou do assembly do Linux.
  \item \textbf{Algoritmos}: as sequ\^encias de inicializa\c{c}\~ao do 8042,
        UART e modo longo seguem a documenta\c{c}\~ao de hardware e padr\~oes
        consolidados, que o Linux tamb\'em segue. Isso \'e converg\^encia
        t\'ecnica, n\~ao c\'opia.
  \item \textbf{Constantes}: valores como \code{0x3F8}, \code{0x60},
        \code{0x64}, \code{0x3D4}, \code{0xB8000} e os bits de \code{CR0},
        \code{CR4} e \code{EFER} s\~ao definidos pelo hardware e aparecem em
        qualquer driver, inclusive no Linux.
  \item \textbf{Licen\c{c}a}: o projeto se declara MIT/Apache-2.0. Se
        trechos GPL-2.0 do Linux fossem copiados, a licen\c{c}a precisaria
        ser compat\'ivel com GPL. Como a implementa\c{c}\~ao \'e
        independente, n\~ao h\'a essa obriga\c{c}\~ao.
\end{itemize}

\section{Similaridade global}

\begin{table}[H]
\centering
\caption{Resumo das diferentes medidas de similaridade.}
\label{tab:resumo-similaridade}
\begin{tabularx}{\linewidth}{@{}lX@{}}
\toprule
\textbf{Medida} & \textbf{Resultado} \\
\midrule
C\'odigo copiado & 0\% (implementa\c{c}\~ao independente) \\
Linhas vs.\ 12 arquivos an\'alogos & $\approx 35{,}67\%$ \\
Linhas vs.\ \'arvore Linux completa & $\approx 0{,}0209\%$ \\
Similaridade funcional ponderada (10 dimens\~oes) & $\approx 27{,}3\%$ \\
Cobertura de subsistemas do kernel completo & muito baixa; ordem de 2\% \\
Similaridade conceitual no boot & alta (mesmos mecanismos x86) \\
Similaridade conceitual em drivers & m\'edia/baixa (subconjunto por polling) \\
Similaridade conceitual em servi\c{c}os & baixa (MM, scheduler e VFS b\'asicos; sem rede) \\
\bottomrule
\end{tabularx}
\end{table}

\subsection{Classifica\c{c}\~ao}

Em uma escala qualitativa:

\begin{itemize}
  \item \textbf{Relacionamento com o Linux}: ``inspirado nos mesmos
        conceitos e no mesmo hardware'', n\~ao ``porte'' nem ``clone''.
  \item \textbf{Similaridade global}: \textbf{baixa}.
  \item \textbf{Valor did\'atico}: alto para entender boot, modo longo,
        VGA, PS/2 e serial; baixo para entender o Linux como sistema
        operacional completo.
  \item \textbf{Nota de 0 a 10 para ``quanto do Linux foi reimplementado''}:
        aproximadamente \textbf{2,7/10} na pontua\c{c}\~ao ponderada; menos de
        \textbf{0,1/10} se o crit\'erio for o kernel completo em
        funcionalidades.
\end{itemize}

\begin{nota}[Como usar esta an\'alise]
Se o objetivo for estudar o Linux, este projeto serve como ``mapa do
caminho inicial'': ele mostra, em c\'odigo pequeno, o que acontece antes de
o kernel Linux come\c{c}ar a usar seus subsistemas. Se o objetivo for
substituir o Linux, o projeto est\'a a v\'arias ordens de grandeza de
dist\^ancia: faltam CFS/EEVDF, \code{fork}/\code{exec}, espa\c{c}os de
endere\c{c}o por processo, sistema de arquivos em disco, rede, drivers e
endurecimento de seguran\c{c}a.
\end{nota}

\section{Diferen\c{c}as essenciais}

\begin{table}[H]
\centering
\caption{Diferen\c{c}as essenciais entre o projeto e o Linux.}
\label{tab:diferencas}
\begin{tabularx}{\linewidth}{@{}p{3.2cm}XX@{}}
\toprule
\textbf{Aspecto} & \textbf{Linux} & \textbf{Rust VGA Kernel} \\
\midrule
Linguagem & C e assembly & Rust e assembly \\
Tamanho & $\sim$38 milh\~oes de linhas & 7.973 linhas \\
Modelo de boot & Protocolo pr\'oprio, setup, EFI & Multiboot2/GRUB \\
Interrup\c{c}\~oes & IDT, PIC/APIC, timers & IDT, PIC/PIT; teclado em polling \\
Concorr\^encia & SMP, preemp\c{c}\~ao, RCU & Uma CPU, escalonador preemptivo
justo e ponderado com sleep bloqueante e wait queues \\
Mem\'oria & Buddy, slab, VM, page cache & Bitmap de frames, demand paging,
heap \emph{first-fit} \\
Processos & Scheduler, sinais, namespaces, cgroups & Tabela de tarefas,
escalonador justo e ponderado, wait queues, sinais e \code{wait}/zumbis; sem
namespaces \\
Entrada/sa\'ida & VFS, block, net, tty & ramfs, VFS, \code{/proc}, portas de I/O diretas \\
Seguran\c{c}a & Rings, SMEP/SMAP, KASLR, LSM & Rings 0/3 e U/S; sem NX,
SMEP ou SMAP \\
Portabilidade & V\'arias arquiteturas & Apenas x86\_64 BIOS \\
Licen\c{c}a & GPL-2.0 & MIT/Apache-2.0 \\
\bottomrule
\end{tabularx}
\end{table}

\chapter{Limita\c{c}\~oes, riscos e evolu\c{c}\~ao}
\label{cap:limitacoes}

\section{Limita\c{c}\~oes t\'ecnicas}

\begin{table}[H]
\centering
\caption{Limita\c{c}\~oes por \'area e consequ\^encias.}
\label{tab:limitacoes}
\begin{tabularx}{\linewidth}{@{}p{3.3cm}X@{}}
\toprule
\textbf{\'Area} & \textbf{Limita\c{c}\~ao e consequ\^encia} \\
\midrule
Interrup\c{c}\~oes & IDT com 256 portas, PIC remapeado e timer IRQ0. IRQ1
do teclado continua mascarada; n\~ao h\'a APIC nem SMP, e a IRQ0 alimenta o
escalonador preemptivo. \\
Mem\'oria & Alocador de frames, p\'aginas de 4 KiB sob demanda e um heap
\emph{first-fit} de 1 MiB com \code{Vec}/\code{String}. Sem slab ou buddy
allocator, sem swap, sem reclaim e sem NUMA. \\
Processos & Tabela de tarefas, PID/PPID, troca de contexto, escalonador
preemptivo justo e ponderado com \code{nice}, pesos no estilo
\code{sched\_prio\_to\_weight} e \code{vruntime}, \code{sleep} bloqueante
(\code{schedule\_timeout}), filas de espera por evento
(\code{wait\_event}/\code{wake\_up}), sinais no estilo Unix com m\'ascaras de
pendentes/bloqueados e disposi\c{c}\~oes, tarefas paradas/continuadas, status
de sa\'ida, zumbis e \code{wait}/\code{SIGCHLD}. N\~ao h\'a
\code{fork}/\code{exec}, \'arvore rubro-negra completa do CFS/EEVDF nem
espa\c{c}o de endere\c{c}o por processo; os programas em ring 3 ainda n\~ao
s\~ao processos escalonados. Doze syscalls
(\code{exit/write/read/get\_tty/open/close/read\_fd/write\_fd/lseek/
getpid/getppid/kill}). \\
Armazenamento & VFS com dezesseis descritores sobre o ramfs em mem\'oria e um
\code{/proc} sint\'etico. Sem driver de disco, block layer, mounts,
permiss\~oes, page cache, FAT ou ext2; a ISO \'e s\'o o meio de boot. \\
Rede & Sem pilha TCP/IP, netdevice ou drivers de rede. \\
V\'ideo & Apenas VGA texto legado. UEFI sem CSM, GPUs e modos gr\'aficos
n\~ao s\~ao suportados. \\
Entrada & Teclado PS/2 por polling, sem USB-HID nativo, Caps Lock, AltGr,
setas ou layout ABNT2. \\
Serial & Apenas COM1, polling, sem IRQ, DMA ou controle de fluxo. \\
Concorr\^encia & Uma CPU; o escalonador \'e preemptivo por timer com
contabilidade justa e ponderada por \code{nice}/\code{vruntime},
\code{sleep} bloqueante e filas de espera por evento, mas a preemp\c{c}\~ao
fica suspensa enquanto um programa de ring 3 est\'a ativo. Os spinlocks usam
\code{irqsave}, ent\~ao a tabela de tarefas \'e segura contra a IRQ0. \\
Seguran\c{c}a & U/S separa usu\'ario e kernel, mas n\~ao h\'a ASLR, NX,
SMEP/SMAP, \emph{stack canary} nem espa\c{c}os de endere\c{c}o por processo.
Ponteiros do usu\'ario s\~ao validados por intervalo, mas n\~ao h\'a um
auxiliar no estilo \code{copy\_from\_user}. \\
Robustez & IDT evita \emph{triple fault} na maioria dos casos; n\~ao h\'a
recupera\c{c}\~ao de falhas, watchdog ou logs persistentes. \\
Portabilidade & C\'odigo espec\'ifico para x86\_64 BIOS; n\~ao compila para
ARM, RISC-V ou UEFI puro. \\
\bottomrule
\end{tabularx}
\end{table}

\section{Riscos}

\begin{table}[H]
\centering
\caption{Riscos e mitiga\c{c}\~oes atuais.}
\label{tab:riscos}
\begin{tabularx}{\linewidth}{@{}p{3.7cm}p{3.0cm}X@{}}
\toprule
\textbf{Risco} & \textbf{Probabilidade} & \textbf{Mitiga\c{c}\~ao} \\
\midrule
\emph{Triple fault} por exce\c{c}\~ao & M\'edia em c\'odigo novo & IDT e
\emph{stubs} j\'a instalados; testar novos tratadores em QEMU antes de
habilitar rotas cr\'iticas. \\
\emph{Deadlock} no console & Baixa hoje & N\~ao adquirir o lock duas vezes;
o timer n\~ao imprime, mas futuras IRQs devem usar \code{irqsave}. \\
\emph{Hang} em espera de hardware & M\'edia & Timeouts nas esperas do 8042;
a UART ainda pode travar se presente e defeituosa. \\
Teclado n\~ao funciona em hardware real & M\'edia & Testar com \emph{legacy
USB}/PS/2; usar serial como alternativa. \\
Tela preta em UEFI & Alta sem CSM & Documentar requisito de BIOS/CSM;
implementar GOP no futuro. \\
Perda de bytes na serial & M\'edia & Polling r\'apido; futuramente FIFO/IRQ. \\
Corrup\c{c}\~ao por bug de \'indice & Baixa & Revisar invariantes de
\code{row}/\code{column}; adicionar testes. \\
\bottomrule
\end{tabularx}
\end{table}

\section{Roadmap sugerido}

A ordem abaixo prioriza ganhos de robustez e de compreens\~ao antes de
adicionar complexidade.

\begin{table}[H]
\centering
\caption{Roadmap priorizado.}
\label{tab:roadmap}
\begin{tabularx}{\linewidth}{@{}clX@{}}
\toprule
\textbf{Prior.} & \textbf{Item} & \textbf{Objetivo} \\
\midrule
1 & Sinais em ring 3 e pilhas por tarefa & A entrega de sinais j\'a funciona
para tarefas de kernel; estend\^e-la a programas de ring 3 e processos
escalonados com pilha de kernel por tarefa. \\
2 & Processos de verdade & \code{fork}/\code{exec} (com \code{wait} j\'a
presente), tabelas de p\'agina por processo e tarefas de ring 3 escalonadas
pelo kernel. \\
3 & VFS e sistema de arquivos & Inodes, dentries, mounts e permiss\~oes;
driver de disco e initramfs. \\
4 & Robustez de IRQ e teclado & IRQ1 com buffer de scancodes e tratamento de
IRQs esp\'urias. \\
5 & Disciplina de linha TTY & Modo can\^onico, sinais, controle de jobs,
\code{Ctrl+C}/\code{Ctrl+D} e um \code{getty} por console. \\
6 & Framebuffer GOP/VBE & Suporte a UEFI e modos gr\'aficos; fonte bitmap. \\
7 & Endurecimento de mem\'oria & NX, SMEP/SMAP, valida\c{c}\~ao de ponteiros
de usu\'ario e ASLR. \\
8 & Driver de teclado completo & USB-HID, Caps Lock, AltGr, layout ABNT2,
teclas estendidas. \\
9 & Rede & NE2000/virtio-net, ARP/IP/ICMP/UDP, testes com QEMU. \\
10 & Testes automatizados & CI com QEMU headless, dumps VGA e asserts. \\
\bottomrule
\end{tabularx}
\end{table}

\section{Crit\'erios para considerar o projeto ``um kernel''}

Do ponto de vista did\'atico, o projeto j\'a \'e um kernel: roda em ring 0,
controla hardware e n\~ao depende de sistema operacional. Do ponto de vista
de um sistema operacional utiliz\'avel, ainda faltam, no m\'inimo:

\begin{itemize}
  \item isolamento de processos e espa\c{c}o de usu\'ario por processo;
  \item entrega de sinais em programas de ring 3, espa\c{c}os de endere\c{c}o
        por processo e pilhas de kernel por tarefa para ring 3
        (\code{nice}, pesos, \code{vruntime}, sleep, filas de espera por
        evento, sinais em contexto de kernel e \code{wait} j\'a existem);
  \item sistema de arquivos persistente e \emph{loader} de programas;
  \item drivers de entrada/sa\'ida al\'em de VGA/PS/2/serial;
  \item tratamento de erros e recupera\c{c}\~ao.
\end{itemize}

O projeto atual \'e melhor classificado como \textbf{kernel monol\'itico
m\'inimo de demonstra\c{c}\~ao}, focado em boot, mem\'oria,
interrup\c{c}\~oes, arquivos, tarefas e console.

\appendix

\chapter{Cat\'alogo de fun\c{c}\~oes}
\label{ap:funcoes}

Esta se\c{c}\~ao lista as 129 fun\c{c}\~oes do c\'odigo Rust, agrupadas por
m\'odulo. O \emph{stub} \code{\_start} n\~ao \'e uma fun\c{c}\~ao Rust e
aparece separadamente no Cap\'itulo~\ref{cap:boot}. Declara\c{c}\~oes
\code{extern} sem corpo, como \code{gdt\_reload\_cs}, \code{user\_enter} e
\code{user\_resume}, n\~ao s\~ao contadas.

\section{\code{main.rs} (2 fun\c{c}\~oes)}

\begin{longtable}{@{}p{6.2cm}p{8.3cm}@{}}
\caption{Fun\c{c}\~oes de \code{main.rs}.}\\
\toprule
\textbf{Assinatura} & \textbf{Descri\c{c}\~ao} \\
\midrule
\endfirsthead
\toprule
\textbf{Assinatura} & \textbf{Descri\c{c}\~ao} \\
\midrule
\endhead
\code{kernel\_main(multiboot\_magic: u32, multiboot\_info: u32) -> !} &
Ponto de entrada Rust; inicializa drivers, imprime banner e entra no shell. \\
\code{panic(info: \&PanicInfo) -> !} &
\emph{Panic handler}; imprime a mensagem e para a CPU. \\
\bottomrule
\end{longtable}

\section{\code{arch.rs} (5 fun\c{c}\~oes)}

\begin{longtable}{@{}p{5.2cm}p{9.3cm}@{}}
\caption{Fun\c{c}\~oes de \code{arch.rs}.}\\
\toprule
\textbf{Assinatura} & \textbf{Descri\c{c}\~ao} \\
\midrule
\endfirsthead
\toprule
\textbf{Assinatura} & \textbf{Descri\c{c}\~ao} \\
\midrule
\endhead
\code{inb(port: u16) -> u8} &
L\^e um byte de uma porta. \code{unsafe}. \\
\code{outb(port: u16, value: u8)} &
Escreve um byte em uma porta. \code{unsafe}. \\
\code{enable\_interrupts()} &
Executa \code{sti} para habilitar interrup\c{c}\~oes mascar\'aveis. \\
\code{halt\_loop() -> !} &
La\c{c}o infinito com \code{cli; hlt}. \\
\code{reboot() -> !} &
Envia \code{0xFE} ao 8042; se falhar, para a CPU. \\
\bottomrule
\end{longtable}

\section{\code{sync.rs} (5 fun\c{c}\~oes)}

\begin{longtable}{@{}p{5.7cm}p{8.8cm}@{}}
\caption{Fun\c{c}\~oes de \code{sync.rs}.}\\
\toprule
\textbf{Assinatura} & \textbf{Descri\c{c}\~ao} \\
\midrule
\endfirsthead
\toprule
\textbf{Assinatura} & \textbf{Descri\c{c}\~ao} \\
\midrule
\endhead
\code{SpinLock::new(data: T) -> SpinLock<T>} &
Cria o lock em contexto \code{const}. \\
\code{SpinLock::lock(\&self) -> SpinLockGuard<'\_, T>} &
Adquire o lock girando at\'e conseguir. \\
\code{SpinLockGuard::deref(\&self) -> \&T} &
Acesso imut\'avel ao dado protegido. \\
\code{SpinLockGuard::deref\_mut(\&mut self) -> \&mut T} &
Acesso mut\'avel ao dado protegido. \\
\code{SpinLockGuard::drop(\&mut self)} &
Libera o lock com ordena\c{c}\~ao \code{Release}. \\
\bottomrule
\end{longtable}

\section{\code{vga.rs} (5 fun\c{c}\~oes)}

\begin{longtable}{@{}p{6.1cm}p{8.4cm}@{}}
\caption{Fun\c{c}\~oes de \code{vga.rs}.}\\
\toprule
\textbf{Assinatura} & \textbf{Descri\c{c}\~ao} \\
\midrule
\endfirsthead
\toprule
\textbf{Assinatura} & \textbf{Descri\c{c}\~ao} \\
\midrule
\endhead
\code{blank\_cell() -> u16} &
C\'elula de espa\c{c}o com atributo padr\~ao; \code{const fn}. \\
\code{init()} &
Limpa o hardware e configura a forma do cursor. \\
\code{clear\_hardware()} &
Preenche as 2.000 c\'elulas do framebuffer com c\'elulas vazias. \\
\code{render(cells, row, column)} &
Copia um buffer de console para \code{0xB8000} e posiciona o cursor. \\
\code{set\_cursor(row, column)} &
Escreve a posi\c{c}\~ao nos registradores da CRTC. \\
\bottomrule
\end{longtable}

\section{\code{tty.rs} (31 fun\c{c}\~oes)}

\begin{longtable}{@{}p{6.0cm}p{8.5cm}@{}}
\caption{Fun\c{c}\~oes de \code{tty.rs}.}\\
\toprule
\textbf{Assinatura} & \textbf{Descri\c{c}\~ao} \\
\midrule
\endfirsthead
\toprule
\textbf{Assinatura} & \textbf{Descri\c{c}\~ao} \\
\midrule
\endhead
\code{Tty::new() -> Tty} &
Inicializa um console vazio; \code{const fn}. \\
\code{Tty::clear(\&mut self)} &
Preenche as c\'elulas com espa\c{c}os e zera o cursor. \\
\code{Tty::push\_input(\&mut self, byte: u8)} &
Anexa um byte \`a fila de entrada do console. \\
\code{Tty::pop\_input(\&mut self) -> Option<u8>} &
Remove um byte da fila. \\
\code{Tty::write\_bytes(\&mut self, bytes: \&[u8])} &
Escreve uma fatia no buffer do console. \\
\code{Tty::write\_byte(\&mut self, byte: u8)} &
Despacha texto simples e sequ\^encias de escape ANSI. \\
\code{Tty::write\_plain(\&mut self, byte: u8)} &
Interpreta controles, backspace e ASCII imprim\'ivel. \\
\code{Tty::write\_csi(\&mut self, byte: u8)} &
Acumula par\^ametros CSI e chama o despachante final. \\
\code{Tty::push\_csi\_param(\&mut self)} &
Guarda um par\^ametro decimal da sequ\^encia CSI. \\
\code{Tty::csi\_param(\&self, index, default) -> u16} &
Devolve um par\^ametro ou seu valor padr\~ao. \\
\code{Tty::csi\_count\_param(\&self, index) -> usize} &
Devolve um par\^ametro de contagem, tratando zero como um. \\
\code{Tty::execute\_csi(\&mut self, final\_byte: u8)} &
Executa CUP, movimento do cursor e comandos de apagar. \\
\code{Tty::erase\_display(\&mut self, mode: u16)} &
Implementa ED (\code{ESC[J}, \code{ESC[2J}). \\
\code{Tty::erase\_line(\&mut self, mode: u16)} &
Implementa EL (\code{ESC[K}). \\
\code{Tty::blank\_row(\&mut self, row: usize)} &
Preenche uma linha com c\'elulas vazias. \\
\code{Tty::write\_char(\&mut self, ch: char)} &
Grava uma c\'elula e avan\c{c}a a coluna. \\
\code{Tty::newline(\&mut self)} &
Avan\c{c}a a linha, com rolagem na \'ultima. \\
\code{Tty::backspace(\&mut self)} &
Recua uma posi\c{c}\~ao e apaga a c\'elula. \\
\code{Tty::scroll(\&mut self)} &
Move o conte\'udo uma linha para cima. \\
\code{init()} &
Limpa os oito consoles e renderiza o \code{tty0}. \\
\code{active() -> usize} &
Devolve o \'indice do console ativo. \\
\code{switch(index: usize) -> bool} &
Troca o console ativo e o redesenha. \\
\code{clear()} &
Escreve \code{ESC[2J ESC[H} no console ativo. \\
\code{write\_bytes(bytes: \&[u8])} &
Escreve no console ativo e o renderiza. \\
\code{write\_to(index: usize, bytes: \&[u8])} &
Escreve em um console espec\'ifico sem trocar o ativo. \\
\code{read\_byte() -> Option<u8>} &
Retira um byte da fila do console ativo. \\
\code{read\_byte\_from(index: usize) -> Option<u8>} &
Retira um byte da fila de um console espec\'ifico. \\
\code{pump\_input()} &
Drena teclado e serial para o console ativo; trata
\code{Alt+F1}..\code{Alt+F8}. \\
\code{push\_input(byte: u8)} &
Auxiliar interno que enfileira um byte por console. \\
\code{<TtyWriter as fmt::Write>::write\_str} &
Anexa texto formatado ao console ativo e espelha na serial. \\
\code{\_print(args: fmt::Arguments<'\_>)} &
Ponto de entrada das macros \code{print!}/\code{println!}. \\
\bottomrule
\end{longtable}

\section{\code{keyboard.rs} (10 fun\c{c}\~oes)}

\begin{longtable}{@{}p{5.3cm}p{9.2cm}@{}}
\caption{Fun\c{c}\~oes de \code{keyboard.rs}.}\\
\toprule
\textbf{Assinatura} & \textbf{Descri\c{c}\~ao} \\
\midrule
\endfirsthead
\toprule
\textbf{Assinatura} & \textbf{Descri\c{c}\~ao} \\
\midrule
\endhead
\code{build\_keymap() -> [u8; 128]} &
Constr\'oi a tabela de scancodes em tempo de compila\c{c}\~ao. \\
\code{apply\_shift(byte: u8) -> u8} &
Aplica Shift a letras, d\'igitos e pontua\c{c}\~ao. \\
\code{wait\_input\_clear() -> bool} &
Espera o buffer de entrada esvaziar. \code{unsafe}. \\
\code{wait\_output\_full() -> bool} &
Espera um byte na sa\'ida. \code{unsafe}. \\
\code{flush\_output()} &
Descarta at\'e 64 bytes pendentes. \code{unsafe}. \\
\code{send\_command(command: u8)} &
Envia comando ao 8042. \code{unsafe}. \\
\code{send\_data(data: u8)} &
Envia dado ao teclado. \code{unsafe}. \\
\code{init()} &
Configura tradu\c{c}\~ao, IRQ, clock e varredura. \\
\code{read\_key() -> Option<Key>} &
L\^e e traduz um scancode; devolve um caractere ou pedido de console. \\
\code{read\_char() -> Option<u8>} &
Inv\'olucro de compatibilidade que devolve apenas \code{Key::Char}. \\
\bottomrule
\end{longtable}

\section{\code{serial.rs} (5 fun\c{c}\~oes)}

\begin{longtable}{@{}p{4.8cm}p{9.7cm}@{}}
\caption{Fun\c{c}\~oes de \code{serial.rs}.}\\
\toprule
\textbf{Assinatura} & \textbf{Descri\c{c}\~ao} \\
\midrule
\endfirsthead
\toprule
\textbf{Assinatura} & \textbf{Descri\c{c}\~ao} \\
\midrule
\endhead
\code{init()} &
Configura COM1 em 38400 8N1 e detecta presen\c{c}a. \\
\code{read\_byte() -> Option<u8>} &
Leitura n\~ao bloqueante de um byte. \\
\code{write\_byte(byte: u8)} &
Envia um byte, esperando o THR. Interna. \\
\code{write\_bytes(bytes: \&[u8])} &
Envia bytes com convers\~ao de \code{\textbackslash n} para CRLF. \\
\code{write\_str(text: \&str)} &
Inv\'olucro de \code{write\_bytes} para strings. \\
\bottomrule
\end{longtable}

\section{\code{shell.rs} (9 fun\c{c}\~oes)}

\begin{longtable}{@{}p{4.8cm}p{9.7cm}@{}}
\caption{Fun\c{c}\~oes de \code{shell.rs}.}\\
\toprule
\textbf{Assinatura} & \textbf{Descri\c{c}\~ao} \\
\midrule
\endfirsthead
\toprule
\textbf{Assinatura} & \textbf{Descri\c{c}\~ao} \\
\midrule
\endhead
\code{run() -> !} &
La\c{c}o eterno de prompt, leitura e execu\c{c}\~ao. \\
\code{trim(bytes: \&[u8]) -> \&[u8]} &
Remove espa\c{c}os ASCII das extremidades. \\
\code{as\_text(bytes: \&[u8]) -> \&str} &
Converte bytes em \code{\&str} ou texto de erro. \\
\code{execute(line: \&[u8])} &
Compara e executa o comando. \\
\code{parse\_console(arg: \&[u8]) -> Option<usize>} &
Interpreta um \'indice de console de um d\'igito para \code{tty} e \code{getty}. \\
\code{getty\_demo(console: Option<usize>)} &
Executa o login em ring 3 no console escolhido ou ativo. \\
\code{demand\_demo()} &
Escreve na regi\~ao de demand paging do kernel e mostra o valor. \\
\code{user\_demo(protection: bool)} &
Executa o demo em ring 3 e imprime o status ou a falha. \\
\code{memory\_info()} &
Mostra frames, page faults, ticks e flags U/S das p\'aginas. \\
\bottomrule
\end{longtable}

\section{\code{gdt.rs} (7 fun\c{c}\~oes)}

\begin{longtable}{@{}p{6.0cm}p{8.5cm}@{}}
\caption{Fun\c{c}\~oes de \code{gdt.rs}.}\\
\toprule
\textbf{Assinatura} & \textbf{Descri\c{c}\~ao} \\
\midrule
\endfirsthead
\toprule
\textbf{Assinatura} & \textbf{Descri\c{c}\~ao} \\
\midrule
\endhead
\code{stack\_top<const N: usize>(stack: *mut Stack<N>) -> u64} &
Calcula o topo de uma pilha est\'atica. \\
\code{gdt\_reload\_cs()} &
Recarrega \code{CS} ap\'os carregar a nova GDT. \\
\code{user\_kernel\_stack\_top() -> u64} &
Topo da pilha de 16 KiB usada em interrup\c{c}\~oes vindas de ring 3. \\
\code{double\_fault\_stack\_top() -> u64} &
Topo da pilha IST1 de 8 KiB para \code{\#DF}. \\
\code{descriptor(base, limit, access, flags) -> u64} &
Monta um descritor de segmento de 64 bits. \\
\code{tss\_descriptor(base: u64, limit: u32) -> (u64, u64)} &
Monta as duas entradas de 8 bytes do descritor de TSS. \\
\code{init()} &
Carrega a GDT definitiva, a TSS e as pilhas de privil\'egio. \\
\bottomrule
\end{longtable}

\section{\code{idt.rs} (14 fun\c{c}\~oes)}

\begin{longtable}{@{}p{6.2cm}p{8.3cm}@{}}
\caption{Fun\c{c}\~oes de \code{idt.rs}.}\\
\toprule
\textbf{Assinatura} & \textbf{Descri\c{c}\~ao} \\
\midrule
\endfirsthead
\toprule
\textbf{Assinatura} & \textbf{Descri\c{c}\~ao} \\
\midrule
\endhead
\code{init()} &
Instala as 256 portas, remapeia o PIC e liga o PIT em IRQ0. \\
\code{set\_gate(idt, vector, handler, dpl, ist)} &
Preenche uma porta com seletor de c\'odigo, DPL e \'indice IST. \\
\code{io\_wait()} &
Pequeno atraso de I/O para o PIC assimilar comandos. \\
\code{remap\_pic()} &
Move as IRQs do PIC mestre/escravo para \code{0x20..0x2F}. \\
\code{pit\_init()} &
Programa o divisor do PIT para aproximadamente 100 Hz. \\
\code{eoi(irq: u8)} &
Envia \emph{End of Interrupt} ao PIC mestre e, se preciso, ao escravo. \\
\code{ticks() -> u64} &
Devolve o contador de interrup\c{c}\~oes do timer. \\
\code{exception\_name(vector: u64) -> \&'static str} &
Mapeia o vetor para um nome leg\'ivel de exce\c{c}\~ao. \\
\code{interrupt\_dispatch(frame: *mut InterruptFrame)} &
Ponto \'unico chamado pelos \emph{stubs}; despacha exce\c{c}\~oes, IRQs e a
\emph{syscall}. \\
\code{page\_fault(frame: \&mut InterruptFrame)} &
L\^e \code{CR2} e delega ao \code{paging::handle\_page\_fault}; pode encerrar
o usu\'ario. \\
\code{general\_protection(frame: \&mut InterruptFrame)} &
Trata \code{\#GP}; encerra o usu\'ario ou para o kernel. \\
\code{user\_fault(frame: \&mut InterruptFrame) -> bool} &
Verifica se a falha veio de ring 3 com um programa ativo. \\
\code{syscall(frame: \&mut InterruptFrame)} &
Encaminha \code{int \$0x80} para \code{syscall::dispatch}. \\
\code{fatal(frame: \&InterruptFrame) -> !} &
Imprime o diagn\'ostico da falha e para a CPU. \\
\bottomrule
\end{longtable}

\section{\code{paging.rs} (23 fun\c{c}\~oes)}

\begin{longtable}{@{}p{6.2cm}p{8.3cm}@{}}
\caption{Fun\c{c}\~oes de \code{paging.rs}.}\\
\toprule
\textbf{Assinatura} & \textbf{Descri\c{c}\~ao} \\
\midrule
\endfirsthead
\toprule
\textbf{Assinatura} & \textbf{Descri\c{c}\~ao} \\
\midrule
\endhead
\code{bitmap() -> *mut u8} &
Devolve o ponteiro para o bitmap est\'atico de 32 KiB. \code{unsafe}. \\
\code{frame\_is\_used(frame: usize) -> bool} &
Consulta um bit do bitmap. \\
\code{mark\_used(frame: usize)} &
Marca um frame como em uso. \\
\code{mark\_free(frame: usize)} &
Marca um frame como livre. \\
\code{free\_range(start, end, reserve\_end)} &
Libera uma faixa de RAM dispon\'ivel, preservando o fim do kernel. \\
\code{reserve\_range(start: u64, end: u64)} &
Reserva uma faixa (primeiro MiB, kernel, estruturas Multiboot2). \\
\code{init(multiboot\_magic, multiboot\_info)} &
L\^e o mapa Multiboot2 e inicializa o alocador de frames. \\
\code{alloc\_frame() -> Option<u64>} &
Aloca e zera um frame de 4 KiB; devolve \code{None} se a RAM acabou. \\
\code{free\_frame(frame: u64)} &
Devolve um frame ao bitmap. \\
\code{stats() -> Stats} &
Devolve frames alocados, livres e total. \\
\code{demand\_faults() -> u64} &
Devolve o n\'umero de page faults atendidos por demand paging. \\
\code{read\_cr3() -> u64} &
L\^e o PML4 atual de \code{CR3}. \\
\code{invlpg(virt: u64)} &
Invalida uma entrada da TLB. \code{unsafe}. \\
\code{next\_table(entry, user) -> Result<*mut PageTable, \&'static str>} &
Obt\'em ou cria a pr\'oxima tabela, aplicando U/S quando necess\'ario. \\
\code{map\_page(virt, phys, flags, user\_tables)} &
Mapeia uma p\'agina de 4 KiB e invalida a TLB. \code{unsafe}. \\
\code{page\_flags(virt: u64) -> Option<u64>} &
Devolve a entrada final que traduz um endere\c{c}o, sem criar tabelas. \\
\code{unmap\_page(virt: u64) -> Result<(), \&'static str>} &
Remove um mapeamento e libera o frame f\'isico. \\
\code{ensure\_user\_range(addr, len) -> Result<(), \&'static str>} &
Valida um intervalo do usu\'ario e materializa p\'aginas ausentes com U/S = 1. \\
\code{reset\_user\_demand()} &
Desfaz as p\'aginas de pilha/dados do demo de usu\'ario. \\
\code{paging\_page(address: u64) -> u64} &
Alinha um endere\c{c}o para o in\'icio da p\'agina. \\
\code{handle\_page\_fault(addr, error) -> Result<(), \&'static str>} &
Decide entre materializar uma p\'agina ou recusar o acesso. \\
\code{kernel\_demand\_addr(index: usize) -> u64} &
Endere\c{c}o de uma p\'agina da regi\~ao de demand paging do kernel. \\
\code{install\_user\_code(bytes: \&[u8]) -> Result<(), \&'static str>} &
Copia o programa para \code{USER\_CODE}, criando a p\'agina residente. \\
\bottomrule
\end{longtable}

\section{\code{user.rs} (10 fun\c{c}\~oes)}

\begin{longtable}{@{}p{6.0cm}p{8.5cm}@{}}
\caption{Fun\c{c}\~oes de \code{user.rs}.}\\
\toprule
\textbf{Assinatura} & \textbf{Descri\c{c}\~ao} \\
\midrule
\endfirsthead
\toprule
\textbf{Assinatura} & \textbf{Descri\c{c}\~ao} \\
\midrule
\endhead
\code{Program::new() -> Program} &
Cria um buffer vazio de 64 bytes para o c\'odigo de usu\'ario. \\
\code{Program::push(\&mut self, bytes: \&[u8])} &
Anexa bytes ao programa. \\
\code{Program::movabs\_rax(\&mut self, value: u64)} &
Emite \code{movabs rax, imm64}. \\
\code{program(demo: Demo) -> Program} &
Monta o c\'odigo de m\'aquina dos demos \code{Demand} e \code{Protection}. \\
\code{is\_active() -> bool} &
Indica se h\'a um programa de usu\'ario em execu\c{c}\~ao. \\
\code{run(demo: Demo) -> Result<i32, \&'static str>} &
Instala o c\'odigo, entra em ring 3 e devolve o status de sa\'ida. \\
\code{run\_getty(console: Option<usize>) -> Result<i32, \&'static str>} &
Troca para o console pedido, executa o blob de login e devolve o status. \\
\code{abort(status: i32) -> !} &
Chamado por exce\c{c}\~oes/\emph{syscall}; grava o status e retoma o kernel. \\
\bottomrule
\end{longtable}

\section{\code{syscall.rs} (3 fun\c{c}\~oes)}

\begin{longtable}{@{}p{5.6cm}p{8.9cm}@{}}
\caption{Fun\c{c}\~oes de \code{syscall.rs}.}\\
\toprule
\textbf{Assinatura} & \textbf{Descri\c{c}\~ao} \\
\midrule
\endfirsthead
\toprule
\textbf{Assinatura} & \textbf{Descri\c{c}\~ao} \\
\midrule
\endhead
\code{dispatch(frame: \&mut InterruptFrame)} &
Despacha \code{exit}, \code{write}, \code{read} e \code{get\_tty}. \\
\code{sys\_write(frame: \&mut InterruptFrame)} &
Valida o buffer do usu\'ario, escreve no console ativo e devolve o tamanho. \\
\code{sys\_read(frame: \&mut InterruptFrame)} &
Valida o buffer e l\^e uma linha do console, com eco opcional. \\
\bottomrule
\end{longtable}

\chapter{Mapa de portas e mem\'oria}
\label{ap:portas}

\section{Portas de I/O}

\begin{table}[H]
\centering
\caption{Portas de I/O usadas pelo kernel.}
\label{tab:portmap}
\begin{tabular}{@{}llp{7.5cm}@{}}
\toprule
\textbf{Porta} & \textbf{Dispositivo} & \textbf{Uso} \\
\midrule
\code{0x60} & 8042 dados & Scancodes e dados ao teclado \\
\code{0x64} & 8042 status/comando & Status, comandos, reset \\
\code{0x3D4} & VGA CRTC \'indice & Sele\c{c}\~ao de registrador \\
\code{0x3D5} & VGA CRTC dados & Posi\c{c}\~ao e forma do cursor \\
\code{0x20--0x21} & PIC mestre & Comandos, m\'ascara e EOI \\
\code{0xA0--0xA1} & PIC escravo & Comandos, m\'ascara e EOI \\
\code{0x40--0x43} & PIT 8253/8254 & Canal 0, modo e divisor do timer \\
\code{0x3F8--0x3FF} & UART 16550 COM1 & Dados, LSR, divisor, FIFO \\
\bottomrule
\end{tabular}
\end{table}

\section{Mapa de mem\'oria}

\begin{table}[H]
\centering
\caption{Regi\~oes de mem\'oria relevantes.}
\label{tab:memmap}
\begin{tabularx}{\linewidth}{@{}lX@{}}
\toprule
\textbf{Endere\c{c}o} & \textbf{Uso} \\
\midrule
\code{0x00000000--0x000FFFFF} & Mem\'oria baixa, BIOS, estruturas do GRUB. \\
\code{0x00100000} & In\'icio de \code{.boot} (stub, cabe\c{c}alho, GDT,
tabelas de p\'agina, pilha de 8 KiB). \\
\code{0x00107000} & \emph{Load address} de \code{.text} (virtual alto). \\
\code{0x00109518} & \emph{Load address} de \code{.rodata}. \\
\code{0x0010A000} & \emph{Load address} de \code{.data}. \\
\code{0x0010A020} & \emph{Load address} de \code{.bss}; inclui a pilha de
32 KiB do kernel. \\
\code{0x000B8000} & Framebuffer de texto VGA (80x25, 4.000 bytes). \\
\code{0xFFFFFFFF80000000} & Base virtual do kernel (metade alta). \\
\code{0xFFFFFFFF80107000} & In\'icio virtual de \code{.text}. \\
\code{0x0000008000001000} & \code{USER\_CODE}: p\'agina residente do demo. \\
\code{0x0000008000002000} & \code{USER\_DATA}: p\'agina sob demanda. \\
\code{0x0000008000007FF8} & Topo da pilha do usu\'ario; a p\'agina abaixo
\'e sob demanda. \\
\code{0xFFFFFFFFC0000000} & \code{KDEMAND\_BASE}: 2 MiB de demand paging do
kernel. \\
\bottomrule
\end{tabularx}
\end{table}

\begin{nota}[Endere\c{c}os observados]
Os valores de \code{.text}, \code{.rodata}, \code{.data} e \code{.bss}
correspondem \`a compila\c{c}\~ao de refer\^encia. Eles podem mudar com
otimiza\c{c}\~oes, vers\~ao do compilador ou altera\c{c}\~oes no c\'odigo.
O que n\~ao muda \'e o esquema: stub em 1 MiB, kernel em
\code{0xFFFFFFFF80000000}.
\end{nota}

\chapter{Gloss\'ario}
\label{ap:glossario}

\begin{description}[style=nextline,leftmargin=2em]
  \item[BIOS] Firmware legado de inicializa\c{c}\~ao de PCs, respons\'avel
        por detectar hardware e iniciar o bootloader.
  \item[Bootloader] Programa que carrega o kernel; aqui, o GRUB.
  \item[CR0/CR3/CR4] Registradores de controle da CPU x86; controlam
        pagina\c{c}\~ao, modo protegido e features.
  \item[CRTC] Controlador de tubo de raios cat\'odicos da VGA; controla
        cursor e timing de v\'ideo.
  \item[El Torito] Padr\~ao para tornar CDs/DVDs inicializ\'aveis.
  \item[Demand paging] T\'ecnica de criar a p\'agina somente no primeiro
        acesso, quando a CPU gera \code{\#PF}.
  \item[GDT] Global Descriptor Table; define segmentos de c\'odigo e dados.
  \item[GOP] Graphics Output Protocol, interface de framebuffer em UEFI.
  \item[IDT] Interrupt Descriptor Table; define handlers de interrup\c{c}\~oes
        e exce\c{c}\~oes.
  \item[Page fault] Exce\c{c}\~ao (\code{\#PF}) gerada quando uma p\'agina
        n\~ao est\'a presente ou viola as permiss\~oes U/S.
  \item[PIC] Programmable Interrupt Controller 8259; prioriza e entrega as
        IRQs \`a CPU.
  \item[PIT] Programmable Interval Timer 8253/8254; gera a IRQ0 do timer.
  \item[Ring 3] N\'ivel de privil\'egio de usu\'ario; c\'odigo com acesso
        restrito \`a mem\'oria marcada U/S.
  \item[TSS] Task State Segment; fornece \code{rsp0} e pilhas IST ao trocar
        de ring 3 para ring 0.
  \item[LME] Long Mode Enable, bit do registrador EFER que habilita o modo
        de 64 bits.
  \item[LSR] Line Status Register da UART; indica dados dispon\'iveis e
        transmissor vazio.
  \item[Multiboot2] Especifica\c{c}\~ao de protocolo entre bootloader e
        kernel; usada pelo GRUB.
  \item[PML4/PDPT/PD] N\'iveis da tabela de p\'aginas x86\_64.
  \item[Polling] Consulta repetida de hardware, sem interrup\c{c}\~oes.
  \item[Ring 0] N\'ivel de privil\'egio mais alto da CPU; c\'odigo de kernel.
  \item[Scancode] C\'odigo enviado pelo teclado para representar uma tecla.
  \item[Spinlock] Lock que faz a CPU girar at\'e a regi\~ao cr\'itica ser
        liberada.
  \item[Stub] C\'odigo inicial em assembly que prepara a CPU antes do Rust.
  \item[Triple fault] Falha de exce\c{c}\~ao durante o tratamento de outra,
        que faz a CPU resetar.
  \item[UART] Universal Asynchronous Receiver/Transmitter; controlador
        serial.
  \item[VGA texto] Modo de v\'ideo legado com matriz de caracteres em
        \code{0xB8000}.
\end{description}

\begin{thebibliography}{9}

\bibitem{intel}
Intel Corporation.
\newblock \emph{Intel 64 and IA-32 Architectures Software Developer's
Manual}.
\newblock Volume 3A: System Programming Guide, cap\'itulos sobre modo
protegido, pagina\c{c}\~ao e long mode.

\bibitem{multiboot2}
Free Software Foundation.
\newblock \emph{Multiboot2 Specification}.
\newblock Dispon\'ivel em \url{https://www.gnu.org/software/grub/manual/multiboot2/}.

\bibitem{linux}
Linux Kernel Organization.
\newblock \emph{Linux Kernel 7.3.1 Source Tree}.
\newblock Arquivos \code{arch/x86}, \code{drivers/video/console/vgacon.c},
\code{drivers/input}, \code{drivers/tty/serial/8250} e \code{kernel/}.

\bibitem{osdev}
OSDev Wiki.
\newblock \emph{PS/2 Keyboard}, \emph{Serial Ports}, \emph{VGA Text Mode} e
\emph{Setting Up Long Mode}.
\newblock Dispon\'ivel em \url{https://wiki.osdev.org/}.

\bibitem{rust}
The Rust Project.
\newblock \emph{The Rust Programming Language} e \emph{Rust Platform Support:
x86\_64-unknown-none}.
\newblock Dispon\'ivel em \url{https://doc.rust-lang.org/}.

\bibitem{qemu}
QEMU Project.
\newblock \emph{QEMU Emulator User Documentation}.
\newblock Dispon\'ivel em \url{https://www.qemu.org/documentation/}.

\bibitem{grub}
GNU GRUB.
\newblock \emph{GNU GRUB Manual}.
\newblock Dispon\'ivel em \url{https://www.gnu.org/software/grub/manual/grub/}.

\end{thebibliography}


\end{document}
